\section{The effective integral input classification}



\begin{theorem}[the orbifold quotient]\hypertarget{tag-02NY}{}\label{maps:src:F.2.7}
\StatementRecord{maps:src:F.2.7}{proved}{02NY}{the definitions and constructions stated in the proof}


For the torus fillings of an admissible datum, with affine
multiplicities $m_1,\ldots,m_k$, the fundamental group of the
total space has a canonical surjection onto
\[
 \langle x_1,\ldots,x_k\mid
       x_i^{m_i}=1,\ x_1\cdots x_k=1\rangle.
 \tag{F.27}\hypertarget{tag-02EX}{}\label{maps:eq:F.27}
\]
Consequently a simply connected example with $k=2$ requires
$\gcd(m_1,m_2)=1$.  There is no simply connected example with
$k=3$ and $m_1,m_2,m_3\geq2$.  An affine multiplicity-one
point is an ordinary smooth point and reduces to the problem
with one fewer multiple fibre.

\end{theorem}


\begin{proof}
\begin{enumerate}[label={$\langle1\rangle$\arabic*.}]
\item Over the complement of small discs at the special values,
the torus bundle projects surjectively on fundamental groups
to the punctured base.  Quotient its fundamental group by the
normal closure of all regular-fibre loops.  At a finite
filling its root cover is a smooth torus bundle over a disc,
and its free deck group is cyclic of order $m_i$.
The covering exact sequence shows that the quotient of the
filling group by its fibre subgroup is precisely this cyclic
group.  The boundary base meridian maps to its generator.
Thus adjoining that filling imposes exactly
$x_i^{m_i}=1$ after the fibre loops have been killed.
At the index-one Kuhn cusp the fundamental group is generated
by the three deck loops of its dual torus; these are regular
fibre loops.  Its base meridian contracts.  Hence its filling
imposes the ordinary meridian relation and leaves no extra
generator in this quotient.  The base relation gives the
product relation in \eqref{maps:eq:F.27}.  Van Kampen proves that \eqref{maps:eq:F.27} is
the resulting quotient, in particular the asserted surjection.

\item For two points, substitute $x_2=x_1^{-1}$; the quotient is
$\mathbb Z/\gcd(m_1,m_2)$.
For three points of orders at least two it is a nontrivial
triangle group.  Nontriviality can be seen directly.  Take a geodesic triangle
with angles $\pi/m_1,\pi/m_2,\pi/m_3$ in the sphere, Euclidean
plane or hyperbolic plane according as their sum is greater
than, equal to or less than $\pi$.
In the spherical case the possible orders are, after sorting,
$(2,2,n),(2,3,3),(2,3,4),(2,3,5)$, and the corresponding
spherical triangles exist.  The Euclidean and hyperbolic
angle conditions give the other triangles directly.
If $r_1,r_2,r_3$ are reflections in the opposite sides,
the rotations $r_2r_3,r_3r_1,r_1r_2$ have the stated
orders and product one.  They therefore give a nontrivial
representation of \eqref{maps:eq:F.27} into the orientation-preserving
isometry group.  Its quotient cannot be trivial.
Finally a free cyclic action of order one is the identity
construction on a root disc of order one, hence an ordinary
torus bundle filling and not a multiple fibre.

\end{enumerate}
\end{proof}




\begin{proposition}[exclusion of an involutory finite part]\hypertarget{tag-02NZ}{}\label{maps:src:F.2.8}
\StatementRecord{maps:src:F.2.8}{proved}{02NZ}{\SeriesRef{maps:src:F.2.7}, \SeriesRef{maps:src:F.2.15}, \SeriesRef{maps:src:F.3.1}, \eqref{maps:eq:F.3}, \SeriesRef{maps:src:F.2.9}}


A simply connected datum with $b_2=0$, $k=2$ and the stated
finite affine fillings has neither finite linear
monodromy of order one or two.  Its unordered pair of finite
linear orders must therefore belong to
\[
 \{(3,4),(3,5),(3,8),(3,10),(4,5),
       (5,6),(5,8),(5,12)\}.
 \tag{F.28}\hypertarget{tag-02EY}{}\label{maps:eq:F.28}
\]
The affine multiplicities may still be larger than these
linear orders.

\end{proposition}


\begin{proof}
\begin{enumerate}[label={$\langle1\rangle$\arabic*.}]
\item The linear orders divide the affine multiplicities and are
therefore coprime by \ref{maps:src:F.2.7}.  The fixed-rank calculation
of \ref{maps:src:F.2.9}, which uses only \ref{maps:src:F.3.1} and the
evenness in \ref{maps:src:F.2.15}, gives fixed rank two at each finite
generator.  An identity generator would make the product finite and
unipotent, hence the identity, contrary to the rank-three cusp.

\item Suppose the cohomology generators are $A,B$ with $A^2=I$.
The order of $B$ is odd.  The $(-1)$-eigenspace of $A$ has dimension
four, whereas $\ker(AB-I)$ has dimension three.  Their intersection
in the six-dimensional space contains a nonzero vector $v$.
Then $ABv=v$ and $Av=-v$ imply
\[
 Bv=Av=-v,
\]
which contradicts the odd order of $B$.  Interchanging the two
generators proves the same assertion for either finite site.

\item Delete one and two from the complete finite-order list
\eqref{maps:eq:F.3} and enumerate its coprime unordered pairs.
They are exactly the eight pairs in \eqref{maps:eq:F.28}.
Only the divisibility of the linear orders into the affine
multiplicities has been used.
\end{enumerate}
\end{proof}




\begin{proposition}[the only first-cohomology pattern for two points]\hypertarget{tag-02O0}{}\label{maps:src:F.2.9}
\StatementRecord{maps:src:F.2.9}{proved}{02O0}{\SeriesRef{maps:src:F.3.1}, \SeriesRef{maps:src:F.2.15}}


For a simply connected two-multiple-fibre datum with $b_2=0$,
the rational
cohomology system $V$ has one invariant, no coinvariants,
and fixed rank two at each finite generator.  Its quotient
$W=V/\mathbb Qe$ by the invariant line has neither invariants
nor coinvariants.  If one finite characteristic polynomial
on $W$ is $(T-1)\Phi_d(T)$ with $\deg\Phi_d=4$, then $W$
is irreducible over $\mathbb Q$.

\end{proposition}


\begin{proof}
\begin{enumerate}[label={$\langle1\rangle$\arabic*.}]
\item Write $a=\dim V^\Gamma$, $b=\dim(V^*)^\Gamma$ and $f_i$
for the finite fixed ranks.  Theorem \ref{maps:src:F.3.1} gives $a=1$,
vanishing first parabolic cohomology, and
$f_1+f_2=4+b$.  By \ref{maps:src:F.2.15} each $f_i$ is positive and even, so $b$ is
even.  The maximal square-zero cusp has
$\ker N=\operatorname{im}N$.  Therefore every globally
fixed vector belongs to $\operatorname{im}N$, whereas every
globally fixed covector kills $\operatorname{im}N$.
The global invariant and coinvariant spaces pair to zero.

\item Restriction of a globally invariant covector to the fixed
space of a finite generator is injective: averaging that
generator is the projection onto its fixed space, and the
covector is unchanged by averaging.  Its restriction kills
the $a$ global invariants in that space.  Hence
$a+b\leq f_i$ for each $i$.
If $b\geq2$ is even, the even ranks satisfy
$f_i\geq b+2$, so $4+b=f_1+f_2\geq2b+4$, a contradiction.
Thus $b=0$ and $f_1=f_2=2$.

\item The quotient $W$ has no invariant covectors, since its dual
embeds into $V^*$.  If it had a fixed line, its preimage in
$V$ would be an extension of the trivial line by the trivial
line.  Every finite generator on this two-dimensional
space would be unipotent of finite order, hence identity.
This would give two invariants in $V$, a contradiction.
Thus $W$ also has no invariant vectors.

\item If one finite characteristic polynomial is
$(T-1)\Phi_d(T)$ with the second factor irreducible of degree
four, every proper nonzero rational submodule has dimension
one or four.  In dimension one that generator acts as the
identity.  Since the product of the two finite generators
is unipotent, their induced one-dimensional characters are
inverse; hence the other generator also acts trivially.
This would be an invariant line.  In dimension four the
one-dimensional quotient is trivial by the same argument,
giving a coinvariant.  Both have just been excluded.

\end{enumerate}
\end{proof}


\paragraph{The exterior-square representation of $\mathrm{Sp}(4)$.}
We recall the classical five-dimensional orthogonal representation and the comparison of unipotent classes used below.
Let $E=\mathbb C^4$ have a symplectic basis $e_1,e_2,f_1,f_2$.
The invariant bivector $\omega=e_1\wedge f_1+e_2\wedge f_2$ has
nonzero square. The wedge pairing on $\bigwedge^2E$ is symmetric
and nondegenerate. Its orthogonal complement
$E_0=\omega^\perp$ therefore has dimension five and is nondegenerate.
Preservation of $\omega^2$ gives determinant one on $E$;
$\det(\bigwedge^2g)=(\det g)^3$ and the fixed line
$\mathbb C\omega$ therefore give determinant one on $E_0$.
The action gives
\[
 \rho:\operatorname{Sp}(E)\longrightarrow\operatorname{SO}(E_0).
\]
Its differential is injective. Indeed, a symplectic infinitesimal
transformation acting trivially on $E_0$ also acts trivially on
$\mathbb C\omega$, hence on all of $\bigwedge^2E$. In a basis,
the equations $Xe_i\wedge e_j+e_i\wedge Xe_j=0$ for every $i<j$
first make every off-diagonal coefficient of $X$ zero, and then give
$X_{ii}+X_{jj}=0$ for every distinct pair. Since there are four
indices, these equations make every diagonal coefficient zero.
Both Lie algebras have dimension ten, so the image contains a
neighbourhood of the identity in $\operatorname{SO}(E_0)$.
That complex special orthogonal group is connected. For completeness,
orthogonal elimination expresses an orthogonal matrix as a product
of reflections in vectors of nonzero norm: to fix a nonisotropic
vector, reflection in the difference of its image and itself works
when that difference is nonisotropic; otherwise reflection in their
sum followed by reflection in the vector works. At least one of
sum and difference is nonisotropic. Continue on its nondegenerate
orthogonal complement. Determinant one makes the number of
reflections even. The nonisotropic vectors of a complex vector
space of dimension at least two form a path-connected set, since
two generic complex line segments can avoid the zero set of the
norm polynomial except at finitely many parameter values, around
which the paths can be detoured. Varying one vector in a pair of
reflections to the other connects that pair to the identity.
The one-dimensional final complement contributes only the already
accounted-for sign. This proves connectedness. Thus the open image
of $\rho$ is the whole special orthogonal group.

The kernel is $\{I,-I\}$. If an element acts trivially on
$E_0$, it also fixes $\omega$ and hence every bivector.
It preserves every two-plane, as follows by applying it to a
nonzero decomposable bivector of that plane. Intersections of
two-planes show that it preserves every line. A linear map
preserving every line is scalar: apply it to the coordinate lines
and their pairwise sums. Its scalar has square one. In particular
every finite orthogonal element of order $m$ has lifts whose
$m$th powers belong to $\{I,-I\}$; the lifts are finite and
semisimple.

The symplectic transvection fixing $e_1,e_2,f_2$ and sending
$f_1$ to $f_1+e_1$ acts on the following primitive basis:
\[
 e_1\wedge e_2,\quad e_1\wedge f_2,\quad
 e_2\wedge f_1,\quad f_1\wedge f_2,\quad
 e_1\wedge f_1-e_2\wedge f_2.
\]
Its difference from the identity sends the third vector to
the negative of the first and the fourth to the second, and
annihilates the other three. Its orthogonal Jordan type is
therefore $(2,2,1)$. Conversely, for an orthogonal unipotent
$I+N$ of that type, $N^2=0$ implies $N^*=-N$, because
$(I+N)^*=(I+N)^{-1}=I-N$. Its image is an isotropic plane.
After choosing a basis of that plane and a dual isotropic
plane, use the isometry on their residual nondegenerate line
to make the identification of two such decompositions have
determinant one. Changing its sign changes that determinant
and commutes with $N$, which kills the residual line.
In these decompositions, $N$ is determined by a nonzero alternating form on
the dual two-dimensional space. All such forms are equivalent
under its general linear group. Extending that change by the
inverse transpose on the first plane gives an orthogonal
change of basis; its determinant is one. Thus these unipotents
form one special-orthogonal conjugacy class. Surjectivity of
$\rho$ lifts the conjugacy. The two lifts of $I+N$ are a
symplectic transvection and its negative. 

Finally a symplectic diagonal element with eigenvalues
$x,x^{-1},y,y^{-1}$ acts on $E_0$ with eigenvalues
\[
 1,\quad xy,\quad x/y,\quad y/x,\quad (xy)^{-1}.
\]
This gives the stated eigenvalue comparison through the primitive exterior-square representation.

\begin{theorem}[exclusion of three further order pairs]\hypertarget{tag-02MZ}{}\label{maps:src:F.2.10}
\StatementRecord{maps:src:F.2.10}{proved}{02MZ}{\SeriesRef{maps:src:F.2.9}, \eqref{maps:eq:F.31}, \SeriesRef{maps:src:F.2.8}, \SeriesRef{maps:src:F.3.1}}


No simply connected datum with $k=2$, $b_2=0$, and finite linear orders
$(3,5)$, $(3,8)$ or $(3,10)$ has the required maximal
square-zero cusp.  Among such data, the
remaining necessary unordered linear order pairs are
\[
 \{(3,4),(4,5),(5,6),(5,8),(5,12)\}.
 \tag{F.30}\hypertarget{tag-02F0}{}\label{maps:eq:F.30}
\]
When one order is five, the other finite characteristic
polynomial cannot be $(T-1)^2\Phi_4(T)^2$ or
$(T-1)^2\Phi_6(T)^2$.

\end{theorem}


\begin{proof}
\begin{enumerate}[label={$\langle1\rangle$\arabic*.}]
\item For each proposed pair, \ref{maps:src:F.2.9} gives an
irreducible rational rank-five quotient $W$.  Its cusp has Jordan
type $(2,2,1)$.  One finite polynomial is
$(T-1)\Phi_3(T)^2$ and the other is $(T-1)\Phi_d(T)$ for
$d=5,8,10$.  For the two additional assertions, replace
$\Phi_3^2$ by $\Phi_4^2$ or $\Phi_6^2$, with $d=5$.
In all these cases a commuting endomorphism of $W$ restricts to a
rational scalar on the one-dimensional fixed space of the second
generator.  This restriction is injective: a nonzero endomorphism
of an irreducible representation is invertible and cannot vanish
on that line.  Consequently
\[
 \operatorname{End}_{\Gamma}(W)=\mathbb Q.
\]

\item Consider $\mathcal H=\operatorname{Hom}(W,W^*)$.
The local fixed dimensions are $9,5,13$.  At the repeated finite
site, pairing reciprocal eigenspaces of dimensions $1,2,2$ gives
$1+4+4=9$.  At the distinct finite site it gives five.
For the cusp, the ordered pairs of blocks of sizes $2,2,1$
contribute the sum of their minimum sizes, namely thirteen.
The invariant-sheaf Euler formula therefore gives
\[
 \chi(\mathbb P^1,j_*\mathcal H)=-25+9+5+13=2.
 \tag{F.31}\hypertarget{tag-02F1}{}\label{maps:eq:F.31}
\]
Here degree zero is $\operatorname{Hom}_{\Gamma}(W,W^*)$,
and degree two is dual to
$\operatorname{Hom}_{\Gamma}(W^*,W)$.  At least one is nonzero,
so irreducibility gives a nondegenerate invariant bilinear form
$J:W\longrightarrow W^*$, inverting the map if necessary.
Its alternating part is equivariant; if nonzero it would be
invertible, which is impossible in odd dimension.  Thus $J$ is
symmetric.

\item An invariant bivector is an alternating equivariant map
$K:W^*\longrightarrow W$.  The comparison map
\[
 K\longmapsto KJ\in\operatorname{End}_{\Gamma}(W)
\]
makes it a scalar, so $K$ is a scalar multiple of the symmetric
map $J^{-1}$.  Hence $K=0$.  Applying the same argument to an
alternating map $W\longrightarrow W^*$ gives
\[
 (\Lambda^2W)^\Gamma=(\Lambda^2W^*)^\Gamma=0.
\]
Both degree zero and degree two of $j_*\Lambda^2W$ vanish.
Its Euler characteristic must therefore be nonpositive.

\item The repeated finite spectrum has four fixed bivectors and
the distinct spectrum has two.  At the cusp, write the unipotent
module as two blocks of size two and one block of size one.
Their exterior square has two trivial blocks from the individual
wedges, blocks of sizes three and one from the product of the
two size-two blocks, and two further size-two blocks.  Its fixed
rank is six.  Consequently
\[
 \chi(\mathbb P^1,j_*\Lambda^2W)=-10+4+2+6=2,
\]
contradicting the preceding step.  This proves all five stated
exclusions.  Removing the three excluded order pairs from
\eqref{maps:eq:F.28} leaves exactly \eqref{maps:eq:F.30}.
\end{enumerate}
\end{proof}




\begin{proposition}[rational uniqueness of the extension]\hypertarget{tag-02N0}{}\label{maps:src:F.2.11}
\StatementRecord{maps:src:F.2.11}{proved}{02N0}{\SeriesRef{maps:src:F.2.9}, \eqref{maps:eq:F.32}}


Let $W$ be a rational rank-five representation of a
twice-punctured affine line, with two finite generators
having fixed rank one and product with unipotent Jordan type
$(2,2,1)$.  Suppose $W$ has neither invariants nor
coinvariants.  Up to rescaling the invariant line, there is
exactly one rational extension
$0\to\mathbb Qe\to V\to W\to0$ in which the two finite
generators retain their finite orders and the cusp has
square-zero logarithm of rank three.  Thus, for the two-point data of \ref{maps:src:F.2.9}, the remaining rational
representation problem is exactly the problem of the
rank-five quotients.  

\end{proposition}


\begin{proof}
\begin{enumerate}[label={$\langle1\rangle$\arabic*.}]
\item Write the finite generators of $W$ as $S_1,S_2$ and the
extension generators as
\[
 A_i=\begin{pmatrix}1&a_i\\0&S_i\end{pmatrix}.
\]
If $S_i$ has order $m_i$, retaining that finite order is equivalent
to $a_i(I+S_i+\cdots+S_i^{m_i-1})=0$.  Averaging identifies this
condition with $a_i\in\operatorname{row}(S_i-I)$, a space of
dimension four.  A change of splitting by $b\in W^*$ sends the
two rows to $a_i+b(S_i-I)$.  This coboundary map is injective
because $W$ has no invariant covectors.  Thus the extension
space subject to the two finite conditions has dimension $8-5=3$.

\item Put $U=S_1S_2=I+N$, so $N^2=0$ and $\operatorname{rank}N=2$.
The product $A_1A_2$ has upper row $c=a_1S_2+a_2$.
Changing the splitting sends $c$ to $c+bN$.  Hence cusp restriction
is the explicit map
\[
 [(a_1,a_2)]\longmapsto[c]\in W^*/\operatorname{row}N.
\]
It is an isomorphism.  Indeed, a class in its kernel can be
represented with $c=0$, so $a_2=-a_1S_2$.  The finite conditions
then say
\[
 a_1\in\operatorname{row}(S_1-I)
              \cap\operatorname{row}(S_2-I).
\]
This intersection has dimension three: the two four-dimensional
row spaces span $W^*$ because $W$ has no common fixed vector.
The rows $b$ satisfying $bN=0$ also form a three-dimensional
space.  Their map $b\mapsto b(S_1-I)$ lands in the displayed
intersection, since $bU=b$ gives
$b(S_1-I)S_2=b-bS_2$.  It is injective, because
$bS_1=b$ and $bU=b$ together give $bS_2=b$, forcing $b=0$.
It is therefore onto.  Every pair with $c=0$ is a coboundary,
proving injectivity of cusp restriction.  Both its source and
target have dimension three, so it is an isomorphism.

\item The square of the lifted product minus identity vanishes
exactly when $cN=0$.  Under this restriction isomorphism the
square-zero extensions form
\[
 \frac{\operatorname{ann}(\operatorname{im}N)}
      {\operatorname{row}N},
 \tag{F.32}\hypertarget{tag-02F2}{}\label{maps:eq:F.32}
\]
of dimension $3-2=1$.  The lifted nilpotent has the rows of $N$
together with the row $c$, so its rank is three exactly when
$c\notin\operatorname{row}N$.  These are precisely the nonzero
classes in \eqref{maps:eq:F.32}.  Rescaling the invariant line
rescales this one-dimensional rational space and is transitive
on its nonzero classes.  This proves the asserted existence and
uniqueness with the stated finite orders and maximal cusp.
\end{enumerate}
\end{proof}




\begin{theorem}[an integral obstruction for a rational $(5,12)$ family]\hypertarget{tag-02N1}{}\label{maps:src:F.2.12}
\StatementRecord{maps:src:F.2.12}{proved}{02N1}{\eqref{maps:eq:F.33}, \SeriesRef{maps:src:F.2.1}, \SeriesRef{maps:src:F.2.7}, \SeriesRef{night:P1.11}}


Let a rational rank-six cohomology representation have generators
\[
 A=\begin{pmatrix}1&0\\0&A_5\end{pmatrix},\qquad
 B=\begin{pmatrix}1&c\\0&B_5\end{pmatrix},\qquad
 c=(0,0,0,1,0),
\]
where
\[
 A_5=\begin{pmatrix}
 1&0&1&0&0\\1&0&0&0&0\\-1&0&0&1&1\\
 1&-1&0&-1&0\\0&1&0&0&0
 \end{pmatrix},\qquad
 B_5=\begin{pmatrix}
 0&1&0&0&0\\0&0&0&0&1\\1&-2&0&0&0\\
 0&1&0&-1&-1\\0&0&1&1&2
 \end{pmatrix}.
 \tag{F.33}\hypertarget{tag-02F3}{}\label{maps:eq:F.33}
\]
For every invariant integral lattice with an index-one maximal
square-zero cusp, and every free affine filling of multiplicities
$m_1,m_2$ divisible respectively by five and twelve, a simply
connected total space has
\[
 \operatorname{Tor}_3 H_2
       \simeq\mathbb Z/3^{v_3(m_2)}.
 \tag{F.34}\hypertarget{tag-02F4}{}\label{maps:eq:F.34}
\]
In particular none of these data is three-connected.
The standard lattice nevertheless has a primitive cusp,
free affine fillings of orders five and twelve, and a simply
connected total space with rational Betti vector
$(1,0,0,0,4,0,0,0,1)$.

\end{theorem}


\begin{proof}
\begin{enumerate}[label={$\langle1\rangle$\arabic*.}]
\item Put $R=\mathbb Z_{(3)}$ and use the rational coordinates of
\eqref{maps:eq:F.33}.  Let $\mathcal P_R$ be the algebra of blocks
$\left(\begin{smallmatrix}a&b\\0&D\end{smallmatrix}\right)$,
where $a\in R$, $b\in R^{1\times5}$ and $D\in\operatorname{Mat}_5(R)$.
We first determine the integral algebra generated by $A,B$.
The following ordered words have $31$ coordinate columns in the order
$(0,j)$, $0\leq j\leq5$, followed by $(i,j)$,
$1\leq i,j\leq5$:
\[
\begin{gathered}
I,A,B,A^2,AB,BA,B^2,A^3,A^2B,ABA,AB^2,BA^2,\\
BAB,B^2A,B^3,A^3B,A^2BA,A^2B^2,ABA^2,AB^2A,\\
AB^3,BA^3,BA^2B,BAB^2,B^2A^2,B^3A,A^3BA,\\
A^3B^2,A^2BA^2,BA^3B,B^2A^3.
\end{gathered}
\]
Their coordinate determinant is $-32$.  Its adjugate divided by
$-32$ therefore expresses every elementary matrix of
$\mathcal P_R$ as an $R$-linear combination of these words.
Since $A,B\in\mathcal P_R$, this proves
$R[A,B]=\mathcal P_R$ with the coordinate map.

For a stable lattice $L_R$, the elementary idempotents split its
six coordinate fractional ideals; the matrices $E_{ij}$ for
$i,j\geq1$ make the last five equal.  Thus
\[
 L_R=3^aRe_0\oplus3^bR^5,\qquad b\geq a,
\]
where the inequality follows from the matrices $E_{0j}$.
The matrix $N=AB-I$ has square zero and nonzero Smith factors
$1,1,1$, while $A_5B_5-I$ has factors $1,1$.
In this basis of $L_R$, its upper-right row is multiplied by
$3^{b-a}$.  For $b>a$ its reduction modulo three has rank two,
contradicting the primitive rank-three cusp.  Hence $b=a$.
The scalar map $3^{-a}I$ identifies $L_R$ with $R^6$
and intertwines every exterior power of the two generators.
The given cusp operator has the same elementary divisors:
$BA-I=A^{-1}(AB-I)A$, and
$(BA)^{-1}-I=-(BA)^{-1}(BA-I)$.
This proves the asserted local marking for every stable lattice.

All arithmetic in this item is reproduced by a separate program (\S1.5).
The argument from that determinant treats all lattices, without a
bound on their indices.

\item Write $v_i=g_i^{m_i}$ for the finite power vectors.
They are fixed by the homology matrices $M_i$.
The base-orbifold quotient in \cref{maps:src:F.2.7} shows that
simple connectivity forces $\gcd(m_1,m_2)=1$.
After the cusp meridian relation the two base generators are
inverse up to a fibre translation.  Their actions on the abelian
fibre lattice are therefore inverse.  Their linear orders five
and twelve force both actions to be the identity in this quotient.
The exact Smith factors of $[A^t-I\ B^t-I]$ are five units,
with first row zero, so its fibre coinvariant module over $R$
is $R$ through the primitive character $\ell=x_0$.

Set $a_i=\ell(v_i)$.  An arbitrary cusp clutch contributes a scalar
$\delta$ to $g_1+g_2=\delta z$.  The remaining relation matrix is
\[
 \begin{pmatrix}m_1&0&-a_1\\0&m_2&-a_2\\1&1&-\delta\end{pmatrix}.
\]
Its determinant is a unit for a simply connected total space.
Since $3\mid m_2$ and $3\nmid m_1$, reduction modulo three makes
$a_2$ a unit in $R$.  This conclusion includes arbitrary clutch
translations and twists.

\item We determine the second finite-cap restriction over $R$.
Choose a primitive integral invariant character $\alpha$ whose
$R$-span is $\ell$.  The stable sublattice
\[
 \Lambda'=\mathbb Z v_2\oplus\ker\alpha
\]
has index $|\alpha(v_2)|$, prime to three.  It gives a finite torus cover; the affine action lifts because it preserves
$v_2$, preserves $\ker\alpha$, and translates by $v_2/m_2$.
The corresponding reduced cap $F'_2$ covers $F_2$ with the same
degree.  In its basis beginning with $v_2$, the action is
\[
 (u,y)\longmapsto(u+1/m_2,Dy).
\]
Thus $F'_2$ is the mapping torus of $D:T^5\to T^5$, and its
regular six-torus is the $m_2$-fold cover in the circle direction.
The coefficient change has determinant a unit in $R$; under the
scalar marking of item~1, the vertical cohomology action $D^t$
is the displayed $B_5$.

The matrices $D^t-I$ and $\bigwedge^2D^t-I$ have no nonunit
three-primary Smith factors.  Their Wang sequences therefore make
$H^q(F'_2;R)$ free for $q\leq3$.  Pullback to the generic torus
is injective in these degrees, because transfer kills its kernel
by $m_2$ and the source is free.  Deck transformations of
$F'_2\to F_2$ act by translations on the generic torus cover.
Their action there is trivial on cohomology, so the injectivity
just proved makes their action on $H^q(F'_2;R)$ trivial as well.
Transfer for this prime-to-three cover now identifies
$H^q(F_2;R)$ with $H^q(F'_2;R)$, compatibly with their restriction maps.  In particular the degree-two and degree-three
restrictions have zero kernel.

Put $T=D^t$.  Its polynomial is $(x-1)\Phi_{12}(x)$, and
\[
 E=\Phi_{12}(T)=T^4-T^2+I,\qquad E^2=E,\qquad
 \sum_{j=0}^{11}T^j=12E.
\]
The image of this integral rank-one idempotent is the primitive
invariant line.  Hence the norm for every stated $m_2$ is $m_2E$.
The Wang sequence in degree two has quotient
$\operatorname{Inv}(\bigwedge^2T)$ and kernel
$\operatorname{coker}(T-I)$.  Restriction to the $m_2$-fold cover
projects identically onto the former group.  On the latter it is
the horizontal wedge with the norm: a horizontal cell lifts to
$m_2$ successive transported prisms.
Consequently the restriction image $J_2$ projects onto the vertical
invariants, and its intersection with the horizontal invariant
summand is $e_0\wedge m_2\operatorname{Inv}(T)$.
The horizontal inclusion induces the isomorphism
\[
 \frac{\operatorname{Inv}(T)}{m_2\operatorname{Inv}(T)}
 \ \xrightarrow{\ \sim\ }\\
 \frac{\operatorname{Inv}(\bigwedge^2V_R)}{J_2}
 \ \cong\ R/3^rR.
\]
No choice of lifts of the vertical Wang classes changes this
quotient or its kernel.  At the first finite cap, $m_1$ is a unit
in $R$, so transfer and averaging give the full invariant image
and zero kernel.  At the cusp, \cref{night:P1.11} gives injective
restrictions onto the invariant lattices in degrees one, two
and three, by dualizing its saturated-kernel collapse.

\item The integral specialization-sheaf approach is used in
\cite[Lemma 7.25 and Propositions 7.26--7.27, pp.~61--63]{AC26}
and \cite[section ``The integral homology'']{E26}.
Here are the coefficient maps and calculations for the present lattice.  Put $L_q=\bigwedge^qV_R$, $A_q=\bigwedge^qA$,
$B_q=\bigwedge^qB$, and let $j$ include the thrice-punctured base.
For the full invariant-extension sheaf $j_*L_q$ the oriented
cellular complex is
\[
\begin{gathered}
 J_{1,q}\oplus J_{2,q}\oplus J_{0,q}
 \xrightarrow{d_0}L_q^2\xrightarrow{\Delta_q}L_q,\\
 J_{1,q}=\ker(A_q-I),\quad J_{2,q}=\ker(B_q-I),\quad
 J_{0,q}=\ker(B_qA_q-I),\\
 d_0(j_1,j_2,j_0)=(j_1-j_0,j_2-A_qj_0),\\
 \Delta_q(a,b)=(A_q-I)a+(B_q-I)b.
\end{gathered}
\]
These are the base attachment maps: the two cusp paths
carry $j_0$ and $A_qj_0$, whose combined boundary is
$(B_qA_q-I)j_0=0$.  Primitive invariant bases and a unimodular
kernel basis of $\Delta_q$ give the following certificates:
\[
\begin{array}{c|c|c|c|c}
q&\operatorname{rank}H^0&\operatorname{rank}H^1&
 \operatorname{SNF}(d_0\text{ in }\ker\Delta_q)&
 \operatorname{SNF}(\Delta_q)\\\hline
1&1&0&1^6&1^5,2\\
2&0&0&1^{15}&1^{15}\\
3&0&3&1^{16},2&1^{19},2.
\end{array}
\]
Here $1^s$ denotes $s$ entries equal to one.  The row-coordinate
matrices and their unimodular HNF transformations are specified by the preceding formulas.
In particular $H^0(j_*L_2)=H^1(j_*L_2)=0$,
$H^2(j_*L_1)=0$ over $R$, and $H^0(j_*L_3)=0$.

Write $\mathcal R_q=R^qf_*R$.  The stalk/fibre
identification is proper base change
\cite[Tag 09V5, Lemma 20.18.1]{StacksProject}.
The local maps of item~3 give
\[
 0\longrightarrow\mathcal R_2\longrightarrow j_*L_2
 \longrightarrow Q\longrightarrow0,
 \qquad Q_{p_2}=R/3^rR,
\]
with no other support.  Hence its connecting homomorphism gives
$H^1(\mathcal R_2)=R/3^rR$.
The degree-one quotient is also point supported, so
$H^2(\mathcal R_1)=H^2(j_*L_1)=0$.
The degree-three sheaf map is injective, so
$H^0(\mathcal R_3)=0$.

\item The only total-degree-three term in the Leray sequence
\cite[Tag 01F2, Lemma 20.13.4]{StacksProject} is therefore
$E_2^{1,2}=R/3^rR$.  It has no possible differential because
the base has cohomological dimension two.  Thus
$H^3(M;R)=R/3^rR$.
Universal coefficients
\cite[Theorem 3.2, p.~195]{Hatcher2002} identify this torsion
module with $\operatorname{Ext}_R(H_2(M;R),R)$: the quotient
$\operatorname{Hom}_R(H_3(M;R),R)$ is free and must vanish.
The elementary divisor decomposition over the PID $R$ proves
\eqref{maps:eq:F.34}.

\item For the standard lattice take
\[
 v_1=(3,0,0,0,0,0)^t,\qquad
 v_2=(-7,7,-7,7,0,7)^t,\qquad (m_1,m_2)=(5,12).
\]
They are fixed by $A^t,B^t$ respectively.
Their $\ell$-values $3,-7$ are units modulo the respective
multiplicities; applying $\ell$ to a fixed-point equation proves
that every nonidentity affine power is free.
The primitive cusp in item~1 has the integral marking of
\cref{maps:src:F.2.1}.  Glue the two finite quotient caps and
the index-one Kuhn cap to the marked torus bundle over a pair
of pants, with zero additional clutch translation.  The finite
boundary identification is the prescribed logarithmic fibre
translation; the cusp is attached in that integral marking.
These local models give a closed smooth oriented
manifold.  Its remaining van Kampen matrix is
\[
 \begin{pmatrix}5&0&-3\\0&12&7\\1&1&0\end{pmatrix},
\]
with determinant of absolute value one.  It is simply connected.

Over $\mathbb Q$ the finite images are full invariants.
Item~4 gives vanishing $H^1(j_*V_{\mathbb Q})$,
$H^0(j_*\bigwedge^2V_{\mathbb Q})$,
$H^1(j_*\bigwedge^2V_{\mathbb Q})$ and
$H^2(j_*V_{\mathbb Q})$, and the degree-three invariant group
is zero.  The five-term Leray sequence and $H^1(M;\mathbb Q)=0$
make the map from the one-dimensional $H^0(\mathcal R_1)$ to
$H^2(\mathbb P^1;\mathbb Q)$ an isomorphism.
Thus $b_2=b_3=0$.
Euler additivity gives $\chi(M)=6$ from the Kuhn cusp;
all other fibres have Euler number zero.
Poincar\'e duality therefore gives
$(1,0,0,0,4,0,0,0,1)$, as claimed.
\end{enumerate}
\end{proof}




\begin{theorem}[two further rational families excluded integrally]\hypertarget{tag-02N2}{}\label{maps:src:F.2.13}
\StatementRecord{maps:src:F.2.13}{proved}{02N2}{\SeriesRef{maps:src:F.2.11}, \SeriesRef{maps:src:F.2.7}, \SeriesRef{night:P1.9}, \SeriesRef{night:P1.11}; \cite[Tag~09V5, Lemma~20.18.1; Tag~01F2, Lemma~20.13.4]{StacksProject}; \cite[Theorem~3.2, p.~195]{Hatcher2002}}


Let $C_f,C_g$ denote the degree-four companion matrices.
Take $f=\Phi_5$ and either
\[
 \text{case A: }g=\Phi_{12},\qquad
 \text{case B: }g=\Phi_3\Phi_6.
 \tag{F.36}\hypertarget{tag-02F5}{}\label{maps:eq:F.36}
\]
Let $W$ be the quotient of
$\bigwedge^2(C_f^{-1},C_g)$ by its invariant line, and
let $V$ be its unique rational extension of \ref{maps:src:F.2.11}.
For every invariant integral lattice in $V$ with primitive
maximal square-zero cusp, every simply connected free affine
filling, with multiplicities $m_1,m_2$ divisible by five and
six respectively, fails three-connectivity.
More precisely, put $r=v_3(m_2)\geq1$.
In case A, $\operatorname{Tor}_3 H_3$ has order
$3^{2r-1}$.  In case B,
$\operatorname{Tor}_3 H_2\simeq\mathbb Z/3^r$.
For $m_2=6$ the respective obstructing group is $\mathbb Z/3$.

\end{theorem}


\begin{proof}
\begin{enumerate}[label={$\langle1\rangle$\arabic*.}]
\item We first give rational models, including the quotient
maps.  Use the exterior basis $(01,02,03,12,13,23)$ for the companion
pair.  Its common primitive invariant vectors are
\[
 \omega_{\mathrm A}=(1,2,1,2,2,1)^t,\qquad
 \omega_{\mathrm B}=(1,0,1,0,0,1)^t.
\]
Let $T$ have columns $\omega,e_1,\ldots,e_5$; its determinant is one.
The last five rows of $T^{-1}$ give the quotient map.  In that
quotient the two matrices $a,b$, and an extension row $c$, are
\[
\begin{array}{c|c|c}
 &a&b\\\hline
\mathrm A&
\begin{pmatrix}2&-1&-2&1&0\\1&0&-1&0&0\\3&-1&-2&0&1\\
3&0&-2&0&0\\1&1&-1&0&0\end{pmatrix}&
\begin{pmatrix}0&-2&0&1&0\\0&-1&0&0&1\\0&-1&0&0&0\\
1&-2&0&0&0\\0&-1&1&0&-1\end{pmatrix}\\[3pt]
\mathrm B&
\begin{pmatrix}0&-1&0&1&0\\1&0&-1&0&0\\1&-1&0&0&1\\
1&0&0&0&0\\1&1&-1&0&0\end{pmatrix}&
\begin{pmatrix}0&0&0&1&0\\0&-1&0&0&1\\0&-1&0&0&0\\
1&0&0&0&0\\0&-1&1&0&1\end{pmatrix}
\end{array}
\]
\[
 c_{\mathrm A}=(-2,-3,1,2,0),\qquad
 c_{\mathrm B}=(0,-1,1,0,0),\qquad
 A=\operatorname{diag}(1,a),\quad
 B=\begin{pmatrix}1&c\\0&b\end{pmatrix}.
\]
Exterior eigenvalue products give characteristic polynomial
$(x-1)\Phi _5$ for $a$, and $(x-1)(x+1)^2\Phi _3$ or
$(x-1)(x+1)^2\Phi _6$ for $b$.  Each has fixed rank one.
Direct multiplication gives $A^5=B^6=I$, $(AB-I)^2=0$, and
nonzero Smith factors $1,1,1$ for $AB-I$ and $1,1$ for $ab-I$.
The row conditions are $c\sum_{j=0}^5b^j=0$ and $c(ab-I)=0$.
The full quotient matrix algebra verified in the next item implies
that this rank-five quotient is irreducible, with no invariant
vectors or covectors.  Thus \cref{maps:src:F.2.11} identifies these
models with the rational extensions, up to rescaling
their invariant line.

\item Put $R=\mathbb Z_{(3)}$ and let $\mathcal P_R$ be the algebra
of blocks $\left(\begin{smallmatrix}u&v\\0&D\end{smallmatrix}\right)$,
with $u\in R$, $v\in R^{1\times5}$ and $D\in\operatorname{Mat}_5(R)$.
In the coordinate order $(0,j)$, $0\leq j\leq5$, followed by
$(i,j)$, $1\leq i,j\leq5$, the following $31$ words have coordinate
determinants $512$ and $-512$ in cases A and B:
\[
\begin{aligned}
&I,A,B,A^2,AB,BA,B^2,A^3,A^2B,ABA,AB^2,BA^2,BAB,B^2A,\\
&A^3B,A^2BA,A^2B^2,ABA^2,AB^2A,BA^3,BA^2B,BAB^2,B^2A^2,B^2AB,\\
&A^3BA,A^3B^2,A^2BA^2,ABA^3,ABA^2B,BA^3B,A^3BA^2.
\end{aligned}
\]
Their adjugates divided by the displayed determinants express
every elementary matrix of $\mathcal P_R$ as an $R$-linear
combination of these words.  Consequently $R[A,B]=\mathcal P_R$.
Its lower block also proves the irreducibility used in item~1.

For an invariant lattice, the elementary projectors and lower-block
matrix units give
\[
 L_R=3^\alpha Re_0\oplus3^\beta R^5,\qquad \beta\geq\alpha;
\]
the inequality is forced by $E_{0j}$.  In this basis the upper
nilpotent row is multiplied by $3^{\beta-\alpha}$.  If
$\beta>\alpha$, reduction modulo three has rank two, contradicting
the primitive rank-three cusp.  Thus $\beta=\alpha$ and the map $3^{-\alpha}I$ identifies $L_R$ with $R^6$, intertwining every
exterior power.  The other product conventions are related by
\[
 BA-I=A^{-1}(AB-I)A,\qquad
 (BA)^{-1}-I=-(BA)^{-1}(BA-I).
\]
This proves the assertion for all invariant lattices, with no
bound on their indices.

\item We record the consequence of simple connectivity for the
affine translations.  Write a finite affine generator,
after a translation of the torus, as
$x\mapsto M_ix+v_i/m_i$, where $M_iv_i=v_i$ and $v_i$ is its
integral power vector.  This change has an explicit map.  If $t$
is a lift of the original translation and $d$ is the linear order,
put $t_0=d^{-1}\sum_{j=0}^{d-1}M_i^jt$ and $b_0=t-t_0$.
Then
\[
 h=\frac1d\sum_{j=0}^{d-1}(d-1-j)M_i^jb_0,
 \qquad (I-M_i)h=b_0.
\]
For $T_h(x)=x+h$, the conjugate $T_h^{-1}gT_h$ has translation $t_0$, and
$v_i=m_it_0$ is integral because the given affine action has
period $m_i$.

The base-orbifold quotient in \cref{maps:src:F.2.7} shows that
$\gcd(m_1,m_2)=1$.  In particular $m_1$ is a unit in $R$.
The homology augmentation matrix $[A^t-I\mid B^t-I]$ has first
row zero and rank-five minors $-4$ and $-1$ in its last five
rows, respectively.  Its coinvariant quotient over $R$ is
therefore $R$ through $\ell(x)=x_0$.  With $a_i=\ell(v_i)$,
the finite power relations and cusp clutch give the quotient
relation matrix
\[
 \begin{pmatrix}m_1&0&-a_1\\0&m_2&-a_2\\1&1&-\delta\end{pmatrix}.
\]
Here $\delta$ is the character value of the cusp clutch.
The character annihilates the collapsed angular cycles, so these
maps on the finite caps and cusp agree on the boundary and define
a quotient of the abelianized fundamental group.  Simple
connectivity makes its determinant a unit.  Reduction modulo
three forces $a_2$ to be a unit.  Arbitrary clutch translations
are included in $\delta$.

\item Let $\alpha$ be the primitive integral invariant character
whose $R$-span is $\ell$, and let $K=\ker\alpha$.  The invariant
sublattice $\Lambda'=\mathbb Zv_2\oplus K$ has index
$|\alpha(v_2)|$, prime to three.  It gives a torus cover
and a cover $F'_2\to F_2$ of the second reduced finite
fibre.  In coordinates beginning with $v_2$, its affine map is
\[
 (\theta,y)\longmapsto(\theta+1/m_2,Dy).
\]
Thus $F'_2$ is a five-torus mapping torus, and its generic
six-torus is the $m_2$-fold cover in the circle direction.
Restriction to $K$ is the dual quotient
$(x_0,x')\mapsto x'$ from $V_R$ to $W_R$, intertwining $B$ with $b$.
Reversing the mapping-torus circle replaces $b$ by $b^{-1}$;
the identities $b^{-1}-I=-b^{-1}(b-I)$ and
$\sum_{j=0}^5b^{-j}=\sum_{j=0}^5b^j$ give the exact comparison.

Put $b_k=\bigwedge^kb$.  The nonzero Smith factors of $b_k-I$
for $0\leq k\leq4$ are units in $R$.  The Wang sequences make
$H^q(F'_2;R)$ free for $q\leq4$.  Pullback to its generic torus
is injective, since transfer kills its kernel by $m_2$ and the
source is free.  Deck transformations of $F'_2\to F_2$ act by
translations on that generic cover.  Their action on its
cohomology is trivial, and injectivity makes it trivial on
$H^q(F'_2;R)$ as well.  Transfer for this prime-to-three cover
therefore identifies the restriction maps with those for $F_2$.

The covering map on the horizontal term of the Wang
sequence is the norm
\[
 \mathcal N_{m_2,k}=\sum_{j=0}^{m_2-1}b_k^j:
 \operatorname{coker}(b_k-I)\longrightarrow\operatorname{Inv}(b_k).
\]
Indeed, the lift is the union of the $m_2$ successive prisms
with their transported coefficients.  On the vertical invariant
quotient pullback is the identity.  Consequently horizontal
inclusion induces the isomorphism
\[
 \frac{\operatorname{Inv}(b_{q-1})}
      {\mathcal N_{m_2,q-1}\bigwedge^{q-1}R^5}
 \ \xrightarrow{\ \sim\ }\
 \frac{\operatorname{Inv}(\bigwedge^q V_R)}{J_{2,q}}.
\]
Surjectivity follows by subtracting a restriction with the same
vertical component; its kernel is exactly the horizontal norm
image.  Thus no choice of a Wang splitting changes this quotient.

\item Write $\lambda=m_2/6$.  Since $b^6=I$,
$\mathcal N_{m_2,k}=\lambda\mathcal N_{6,k}$.
In case A a primitive basis of $\operatorname{Inv}(b_2)$,
in lexicographic exterior order, is the column matrix
\[
 Q=\begin{pmatrix}
 1&0\\1&4\\2&4\\0&2\\0&3\\1&4\\0&3\\1&0\\0&0\\0&2
 \end{pmatrix}.
\]
The norm coordinates in this basis are
\[
 \begin{pmatrix}
 0&3&0&-3&0&0&-4&3&4&3\\
 0&0&0&0&0&0&2&0&-2&0
 \end{pmatrix}.
\]
Their image has basis $\left(\begin{smallmatrix}1&0\\4&6\end{smallmatrix}\right)$,
and multiplication by
$\left(\begin{smallmatrix}1&0\\-4&1\end{smallmatrix}\right)$
gives $\operatorname{diag}(1,6)$.  This gives the explicit quotient
map $(x,y)\mapsto(x,y-4x)$ modulo $(\lambda,6\lambda)$.
In case B the primitive invariant vector of $b$ is
$(1,0,0,1,0)^t$, and the norm coordinates are $(3,0,0,3,0)$.
After multiplication by the remaining units of $R$, these give
\[
 Q_{2,3}^{\mathrm A}\cong R/3^{r-1}R\oplus R/3^rR,
 \qquad Q_{2,2}^{\mathrm B}\cong R/3^rR.
\]
\item We use the integral specialization and base-sheaf method
of \cite[Lemma~7.25 and Propositions~7.26--7.27, pp.~61--63]{AC26}
and \cite[section ``The integral homology'']{E26}, with the
following explicitly computed coefficients.  Put
$L_q=\bigwedge^qR^6$, $A_q=\bigwedge^qA$, $B_q=\bigwedge^qB$.
The based-sphere complex for $j_*L_q$ is
\[
\begin{gathered}
 J_{1,q}\oplus J_{2,q}\oplus J_{0,q}
 \xrightarrow{d_0} L_q^2\xrightarrow{\Delta_q}L_q,\\
 J_{1,q}=\ker(A_q-I),\quad J_{2,q}=\ker(B_q-I),\quad
 J_{0,q}=\ker(B_qA_q-I),\\
 d_0(j_1,j_2,j_0)=(j_1-j_0,j_2-A_qj_0),\\
 \Delta_q(x,y)=(A_q-I)x+(B_q-I)y.
\end{gathered}
\]
The cusp paths carry $j_0$ and $A_qj_0$, so their combined
boundary is zero.  Inverse generator conventions give a chain
isomorphism: use cusp basis $A_qJ_{0,q}$ for the inverse pair,
and maps $-I$, $\operatorname{diag}(-A_q,-B_q)$, $I$ in degrees
zero, one and two.  Both chain identities follow by multiplication.

Unimodular HNF transformations give primitive invariant bases
and kernel coordinates for $d_0$ in $\ker\Delta_q$.  In both cases
the exact Smith records are
\[
\begin{array}{c|ccc|c|c}
q&\operatorname{rank}H^0&\operatorname{rank}H^1&\operatorname{rank}H^2
 &\operatorname{SNF}(d_0\text{ in }\ker\Delta_q)
 &\operatorname{SNF}(\Delta_q)\\\hline
1&1&0&0&1^6&1^5,2\\
2&0&0&0&1^{15}&1^{14},2\\
3&0&3&0&1^{17}&1^{19},2\\
4&0&0&0&1^{15}&1^{14},2
\end{array}
\]
All displayed groups are therefore free over $R$.  These computations are done by three separate programs (\S1.5).  Their computations use the printed models.

\item Write $\mathcal R_q=R^qf_*R$.  Proper base change
\cite[Tag~09V5, Lemma~20.18.1]{StacksProject} identifies the stalks
with the fibres.  At the first finite cap, $m_1$ is a
unit in $R$, so transfer gives the full invariant restriction
image and zero kernel.  At the second finite cap, item~4 gives
injectivity and the norm quotient.  At the prescribed cusp,
dualizing the primitive-image collapse of
\cref{night:P1.11} gives the invariant restriction image, with
zero kernel in degrees one through three and a free kernel $R$
in degree four.  The latter is dual to the free degree-four
cokernel of the collapse and must be retained.

For $q=2,3$ there is consequently an exact sequence
\[
 0\longrightarrow\mathcal R_q\longrightarrow j_*L_q
 \longrightarrow Q_q\longrightarrow0,
\]
where $Q_q$ is supported at the second finite point.  Since
$H^0(j_*L_q)=0$, its connecting map embeds $Q_q$ in
$H^1(\mathcal R_q)$, with the free quotient $H^1(j_*L_q)$
computed above.  Its torsion is exactly $Q_q$.
Also $H^2(\mathcal R_1)=H^2(\mathcal R_2)=H^2(\mathcal R_3)=0$.
Injectivity gives $H^0(\mathcal R_3)=0$.  The degree-four
free cusp kernel, together with $H^0(j_*L_4)=0$, gives
$H^0(\mathcal R_4)=R$.

\item In the Leray sequence
\cite[Tag~01F2, Lemma~20.13.4]{StacksProject}, total degree three
has only $H^1(\mathcal R_2)$.  Its middle base column has no
incoming or outgoing differential.  Thus case B gives
$H^3(M;R)=R/3^rR$.
In total degree four, the surviving terms are
$H^1(\mathcal R_3)$ and $H^0(\mathcal R_4)=R$;
the degree-two base term is zero and the possible differential
from the latter term lands in $H^2(\mathcal R_3)=0$.
The extension by the free module $R$ splits.  Hence in case A
\[
 \operatorname{Tor}H^4(M;R)
   \cong R/3^{r-1}R\oplus R/3^rR.
\]
Universal coefficients
\cite[Theorem~3.2, p.~195]{Hatcher2002} identify this torsion
with the torsion of $H_3(M;R)$, and similarly identify the
case-B torsion with that of $H_2(M;R)$.  Localization retains
the finite three-primary subgroups by the identity on each
$\mathbb Z/3^s$ summand.  This proves the stated orders and
the case-B isomorphism.  At $m_2=6$ the first case-A summand
vanishes and the second is $\mathbb Z/3$; case B is also
$\mathbb Z/3$.  These groups obstruct three-connectivity.
\item The two displayed rational models have explicit smooth marked
realisations with the required free caps and simply connected filling.
Use their displayed standard cohomology lattice $\mathbb Z^6$ and its
dual homology basis, so $M_1=A^t$, $M_2=B^t$.
The independent exact calculation gives $A^5=B^6=I$, exact orders five
and six, and determinant one. Put $K=BA-I$. Then $K^2=0$, and the
third minor on rows $(0,3,4)$ and columns $(1,2,3)$ is one in both
models. Hence $K$ has rank three and all its nonzero Smith factors
are one. The marked cusp is
\[
 M_0-I=((M_1M_2)^{-1}-I)=-K^t.
\]
Set $b_1=e_0,b_2=e_3,b_3=e_4$, $a_j=(M_0-I)b_j$ and
$S=(a_1,a_2,a_3,b_1,b_2,b_3)$. Direct multiplication gives
$S\in\operatorname{GL}_6(\mathbb Z)$ and
\[
 (M_0-I)S=S\begin{pmatrix}0&I_3\\0&0\end{pmatrix}.
\]
This is the exact integral marking used to transport the index-one
Kuhn cap and its collapse maps.

The affine power vectors in these same homology coordinates are
\[
\begin{array}{c|cc}
 &v_1&v_2\\\hline
 \mathrm A&(1,0,0,0,0,0)&(-1,2,0,-1,0,0)\\
 \mathrm B&(1,0,0,0,0,0)&(-1,0,0,0,0,1).
\end{array}
\]
They satisfy $M_iv_i=v_i$. The first-coordinate covector
$\ell$ is fixed by both $M_i$ and has values $1,-1$ on $v_1,v_2$.
Consequently every nonidentity power of
$x\mapsto M_1x+v_1/5$ or $x\mapsto M_2x+v_2/6$ changes the
$\ell$-circle by $j/5$ or $-j/6$, respectively. These are free
cyclic actions of the exact stated orders. Their diagonal actions
on the product with the rotating disc therefore give the smooth finite quotient caps.

Take the affine torus bundle over the three-holed sphere with these
two boundary holonomies and their inverse product on the third
boundary. Glue the two finite caps and the proper index-one Kuhn
cap of \cref{night:P1.9}, \cref{night:P1.11} by the
specified marked mapping-torus collars and zero cusp clutch.
For an affine cusp monodromy $x\mapsto M_0x+c$, a path of translations
$d(s)$ with $d(0)=0$, $d(1)=c$ gives the collar map
$(x,s)\mapsto(x+d(s),s)$ from the linear mapping torus. The prescribed
zero clutch fixes its integral lift. The lattice map on this collar
is the displayed $S$.

The cap fundamental-group maps retain these choices. On the finite
cap the universal covering $\mathbb R^6\times\Delta$ has the deck
transformation
$(x,t)\mapsto(M_ix+v_i/m_i,e^{2\pi i/m_i}t)$; its $m_i$-th power
is the lattice translation $v_i$. Thus the boundary generators
$a,b$ satisfy $a^5=v_1$, $b^6=v_2$. On the Kuhn cap the polar
projection onto its base three-torus and angular-identity section
induce inverse maps on fundamental groups. In its dual-cell model
the one-skeleton consists of the base edges, and each angular loop
contracts after transport to a dual vertex where that angular
fibre is collapsed. The fibre map is therefore exactly
$\Lambda\to\Lambda/\operatorname{im}(M_0-I)$, expressed by the last
three rows of $S^{-1}$. The meridian bounds the toric coordinate
disc at a smooth point of a central component; the zero clutch
identifies it with $(ab)^{-1}$. Hence the cusp imposes $ab=1$.

For $J=[M_1-I\mid M_2-I]$, the first row is zero. Its lower five
rows have a unit minor on columns $(1,2,3,6,10)$ in case A and
$(1,2,3,4,6)$ in case B, with determinants $-1$ and $1$.
The inverse of that minor supplies integral preimages of
$e_1,\ldots,e_5$. Thus
$\operatorname{im}J=\ker\ell$ integrally. Van Kampen now gives
simple connectivity at the group level: from $b=a^{-1}$ and the
two power relations,
\[
 a^{-1}=a^5a^{-6}=v_1v_2\in\Lambda.
\]
So $a,b$ lie in the abelian image of the fibre lattice, their
conjugation actions become trivial, and $J$ is killed. The exact
remaining fibre quotient is $\mathbb Z$ via $\ell$, with generator
$h$ and images $v_1=h$, $v_2=-h$. The displayed identity forces
$a=1$, and $a^5=v_1$ then kills $h$. Every generator is trivial.
Equivalently, after this nonabelian reduction the residual relation
matrix on $(h,a,b)$ is
\[
 \begin{pmatrix}-1&5&0\\1&0&6\\0&1&1\end{pmatrix},
 \qquad\det=1.
\]

Finally use the checked rational invariant-extension table for these
exact models. The finite covering restrictions identify
rational cap cohomology with invariants. The primitive cusp collapse
of~\ref{night:P1.11} gives the same identification through degree
three and a one-dimensional kernel in degree four. Hence
$R^q p_*\mathbb Q=j_*\Lambda^qV$ for $1\le q\le3$, and
\[
 0\longrightarrow\mathbb Q_{\mathrm{cusp}}
 \longrightarrow R^4p_*\mathbb Q
 \longrightarrow j_*\Lambda^4V\longrightarrow0.
\]
The table therefore leaves, in total degree at most four, only
$E_2^{0,0}=E_2^{2,0}=E_2^{0,1}=\mathbb Q$,
$E_2^{1,3}=\mathbb Q^3$ and $E_2^{0,4}=\mathbb Q$.
Simple connectivity forces the $d_2:E_2^{0,1}\to E_2^{2,0}$ to be an isomorphism. There are no
other possible nonzero differentials on these terms. Thus
$(b_0,\ldots,b_4)=(1,0,0,0,4)$; oriented duality on the closed
smooth eight-manifold gives the full tuple
$(1,0,0,0,4,0,0,0,1)$.
\end{enumerate}
\end{proof}




\begin{theorem}[four order pairs and eight local spectral pairs]\hypertarget{tag-02N3}{}\label{maps:src:F.2.14}
\StatementRecord{maps:src:F.2.14}{proved}{02N3}{\SeriesRef{maps:src:F.2.15}, \SeriesRef{maps:src:F.2.9}, \SeriesRef{maps:src:F.2.7}, \eqref{maps:eq:F.40}, \SeriesRef{maps:src:F.2.2}, \SeriesRef{night:P1.11}}


Every simply connected two-finite-point datum with $b_2=0$ has
one finite linear order equal to five and the other equal
to four, six, eight or twelve.
After suppressing the common factor $(T-1)^2$, its unordered
pair of finite characteristic polynomials is exactly one
of the following necessary possibilities:
\[
\begin{array}{c|c|l}
 \text{orders}&\text{order-five factor}&\text{other factor}\\\hline
 (5,4)&\Phi_5&\Phi_2^2\Phi_4\\
 (5,6)&\Phi_5&\Phi_2^2\Phi_3,\ \Phi_2^2\Phi_6,\ \Phi_3\Phi_6\\
 (5,8)&\Phi_5&\Phi_8\\
 (5,12)&\Phi_5&\Phi_{12},\ \Phi_3\Phi_4,\ \Phi_4\Phi_6.
\end{array}
 \tag{F.39}\hypertarget{tag-02F6}{}\label{maps:eq:F.39}
\]
Its rank-five quotient is rationally irreducible.
In particular all order-three pairs, including $(3,4)$,
are impossible.  The table gives necessary conditions only.

\end{theorem}


\begin{proof}
\begin{enumerate}[label={$\langle1\rangle$\arabic*.}]
\item By \cref{maps:src:F.2.9}, \cref{maps:src:F.2.15}, the cohomology
system $V$ has a one-dimensional invariant space $I$, no invariant
covectors, and two-dimensional finite fixed spaces $F_1,F_2$.
The finite orders are among the coprime pairs of
\cref{maps:src:F.2.8}.  At a finite quotient, averaging identifies
the degree-two specialization with the fixed part of $\Lambda^2V$.
At the cusp, \cref{night:P1.11} gives the surjective homology
collapse with kernel the saturation of
$(\Lambda^2M_0-I)\Lambda^2\Lambda$.  Its dual has image
$\ker(\Lambda^2M_0^t-I)$ under the evaluation pairing.  Consequently
the natural specialization morphism identifies
\[
 R^2g_*\mathbb Q=j_*\Lambda^2V.
\]
The term $H^0(j_*\Lambda^2V)$ has no incoming Leray differential,
and its only possible outgoing differential targets
$H^2(j_*V)=H^0(V^*)^*=0$.  It survives to $H^2(M;\mathbb Q)=0$.
Thus $a_2':=\dim(\Lambda^2V)^\Gamma=0$.

\item Let $\beta:V\to V^*$ be an invariant alternating form.
Since $\beta(I)\subset(V^*)^\Gamma=0$, its radical contains $I$.
For the finite averaging projection $P_i:V\to F_i$, invariance gives
$\beta(u,x)=\beta(u,P_ix)$ when $u\in F_i$.
The alternating restriction to the two-dimensional space $F_i$
has $I$ in its radical and is therefore zero.  Hence
$F_1+F_2\subset\operatorname{rad}\beta$.  This sum has dimension
three because $F_1\cap F_2=I$.  The alternating rank is even,
so every such form has rank at most two.

Suppose $\beta,\gamma$ were independent invariant forms.
They and their sum have rank two, so
$\beta\wedge\gamma=0$.  Their image planes in $V^*$
are distinct invariant two-planes meeting in a rational line.
Each finite generator acts on that line by a sign.  At least one
finite order is odd, and the unipotent product acts trivially on the
line; both signs are therefore $+1$.  This contradicts
$(V^*)^\Gamma=0$.  It follows that
$b_2':=\dim(\Lambda^2V^*)^\Gamma\leq1$.

\item At the cusp, the exterior square of three size-two unipotent
blocks has fixed rank nine: three fixed vectors come from their
individual exterior squares and two from each of the three pairwise
tensor products.  If $d_i$ denotes a finite fixed rank on
$\Lambda^2V$, the invariant-sheaf Euler identity gives
\[
 -15+d_1+d_2+9
 =a_2'-h^1(j_*\Lambda^2V)+b_2'\leq1.
 \tag{F.40}\hypertarget{tag-02F8}{}\label{maps:eq:F.40}
\]
Here the degree-two group is dual to
$(\Lambda^2V^*)^\Gamma$, by compact-support duality and the
extension-by-zero sequence.

\item The two-dimensional unit eigenspace leaves four dimensions
for the nonunit cyclotomic factors.  Local complex compatibility
requires even multiplicity of $-1$.  Partitioning four into the
cyclotomic degrees gives the following complete list:
\[
\begin{array}{c|l}
3&\Phi_3^2\\
4&\Phi_4^2,\ \Phi_2^2\Phi_4\\
5&\Phi_5\\
6&\Phi_6^2,\ \Phi_2^2\Phi_3,\ \Phi_2^2\Phi_6,\ \Phi_3\Phi_6\\
8&\Phi_8\\
10&\Phi_{10}\\
12&\Phi_{12},\ \Phi_3\Phi_4,\ \Phi_4\Phi_6.
\end{array}
\]
The unit eigenspace contributes one fixed bivector.  A repeated
nonreal quadratic pair contributes four further fixed bivectors,
so $d_i=5$ for $\Phi_3^2,\Phi_4^2,\Phi_6^2$.
Every other listed factor contributes two, so $d_i=3$.
Inequality \eqref{maps:eq:F.40} excludes every occurrence of
$d_i=5$.  In particular all order-three pairs disappear.
The coprime order list of \cref{maps:src:F.2.8} therefore reduces
to $(5,4),(5,6),(5,8),(5,12)$; their remaining factors are exactly
the eight entries of \eqref{maps:eq:F.39}.
Each pair contains the order-five factor, so
\cref{maps:src:F.2.9} proves irreducibility of $V/I$.
\end{enumerate}
\end{proof}




\begin{proposition}[automatic local complex compatibility]\hypertarget{tag-02N4}{}\label{maps:src:F.2.15}
\StatementRecord{maps:src:F.2.15}{proved}{02N4}{\SeriesRef{maps:src:F.2.7}, \SeriesRef{maps:src:F.1.4}}


Every admissible datum with simply connected gluing and two finite affine fillings automatically has finite linear actions commuting
with real complex structures.  This assertion concerns the
individual finite actions and does not produce simultaneous
global periods.

\end{proposition}


\begin{proof}
\begin{enumerate}[label={$\langle1\rangle$\arabic*.}]
\item Let $m_i$ be the order of the free cyclic affine action
defining the $i$th filling.  Its linear order divides $m_i$.
By \ref{maps:src:F.2.7} the integers $m_1,m_2$ are coprime, so at least one
of them is odd.  The corresponding linear determinant is
one, because a finite real matrix of odd order cannot have
determinant minus one.  The relation
$M_1M_2M_0=I$ and the unipotence of $M_0$ imply
$\det M_1\det M_2=1$, so both determinants are one.

\item A finite-order real matrix is semisimple.  Its nonreal
eigenvalues occur in conjugate pairs, and its only possible
real eigenvalues are $1$ and $-1$.  Determinant one says
that the $-1$ eigenspace has even real dimension.  Since the
total dimension is six, the $1$ eigenspace also has even
real dimension.  Choose a complex structure on each of
these real eigenspaces by pairing a real basis.  On each
nonreal isotypical real summand, multiplication by $i$ in
its complex eigenspace supplies a commuting complex
structure.  Their direct sum commutes with $M_i$.
The fixed eigenspace is nonzero because the affine action
is free, as in \ref{maps:src:F.1.4}, so its even dimension is at least two.

\item Each chosen structure makes $\mathbb R^6/\mathbb Z^6$ a
complex torus and the given affine map holomorphic; thus it
supplies precisely the local complex-compatibility
hypothesis used in the subsequent reductions.  For three genuine
finite fillings \ref{maps:src:F.2.7} has already excluded simple
connectivity without this argument.

\end{enumerate}
\end{proof}


\begin{theorem}[an unbounded integral family]\hypertarget{tag-02N5}{}\label{maps:src:F.2.16}
\StatementRecord{maps:src:F.2.16}{proved}{02N5}{\eqref{maps:eq:F.U1}, \eqref{maps:eq:F.U3}, \eqref{maps:eq:F.U4}, \eqref{maps:eq:F.U5}, \eqref{maps:eq:F.U6}, \eqref{maps:eq:F.U7}, \SeriesRef{maps:src:F.2.5}}


There are infinitely many simultaneous rational conjugacy classes
of integral rank-six tuples of the linear type $(5,4)$ satisfying \cref{def:admissible},
with primitive maximal square-zero cusp.  An explicit infinite
subfamily has free affine lifts of orders five and four.
Apart from the parameter zero, the rank-five irreducible quotient
in the family is not self-dual.  Thus no finite bound on rational
conjugacy classes is valid for the unrestricted input tuples.


\end{theorem}


\begin{proof}
\begin{enumerate}[label={$\langle1\rangle$\arabic*.}]
\item Let $v\in\mathbb Z$, put $q=v^2-v+1$, and let
$Ae_j=e_{j+1}$ on $\mathbb Q^5$, with indices modulo five. Define
\[
 X=(e_0,\ e_2+ve_4),\qquad
 Y=\begin{pmatrix}
 0&-q&-v^2&-1&v\\
 0&1-vq&-vq&v^2-v&q
 \end{pmatrix},\qquad N=XY.
 \tag{F.U1}\hypertarget{tag-02IK}{}\label{maps:eq:F.U1}
\]
Multiplying $Y$ by the two columns of $X$ gives $YX=0$, so
$N^2=X(YX)Y=0$. Consequently
\[
 C=(I+N)A,\qquad B=A^4(I-N)=C^{-1}.
\]
The cycle gives $A^5=I$. Exact multiplication and determinant
expansion give
\[
 N^2=0,\qquad A^5=C^4=I_5,\qquad
 \chi_C(x)=(x-1)(x+1)^2(x^2+1).
 \tag{F.U2}\hypertarget{tag-02IO}{}\label{maps:eq:F.U2}
\]
For the determinant calculation one may substitute
\[
 (xI-A)^{-1}=
 \frac{x^4I+x^3A+x^2A^2+xA^3+A^4}{x^5-1}
\]
in the rank-two determinant identity
\[
 \det(xI-C)=\det(xI-A)
       \det\bigl(I_2-YA(xI-A)^{-1}X\bigr).
\]
The identity $C^4=I$ is checked separately; it is not inferred
from the repeated factor $x+1$ of the characteristic polynomial.
The factor $x^2+1$ shows that $B$ and $C$ have exact order four.

\item Let $P$ have columns $\mathbf1/q,e_1,e_2,e_3,e_4$, and put
$A'=P^{-1}AP$, $B'=P^{-1}BP$. Since $q>0$ and $\det P=1/q$,
these are rational changes of coordinates. Define
\[
 s=v^2+v-1,\qquad r=-\frac{(v+1)(v^2-2)}2,
 \qquad
 \beta=(s,\ s-vr,\ v^2s+vr,\ r,\ -vs-r)
 \tag{F.U3}\hypertarget{tag-02IP}{}\label{maps:eq:F.U3}
\]
and
\[
 \widehat A=\begin{pmatrix}1&0\\0&A'\end{pmatrix},\qquad
 \widehat B=\begin{pmatrix}1&\beta\\0&B'\end{pmatrix}.
 \tag{F.U4}\hypertarget{tag-02IQ}{}\label{maps:eq:F.U4}
\]

The lattice $P\mathbb Z^5$ has basis
$f_0=\mathbf1/q,f_1=e_1,\ldots,f_4=e_4$.
Its images under $A$ are $f_0,f_2,f_3,f_4$ and
$qf_0-f_1-f_2-f_3-f_4$, respectively; $A^{-1}=A^4$
therefore also preserves this lattice. The columns of $X$ are
$qf_0-f_1-f_2-f_3-f_4$ and $f_2+vf_4$, while
\[
 Y':=YP=\begin{pmatrix}
 -2&-q&-v^2&-1&v\\
 2(1-v)&1-vq&-vq&v^2-v&q
 \end{pmatrix}.
\]
Thus $X$ and $Y$ are integral maps between the lattices.
The already proved $YX=0$ shows that $I+N$ and its inverse $I-N$
preserve $P\mathbb Z^5$. Since $B=A^{-1}(I-N)$, both
$A',B'$ belong to $\operatorname{GL}_5(\mathbb Z)$.
For $r$, odd $v$ makes $v+1$ even, and even $v$ makes
$v^2-2$ even; hence $r$ and all entries of $\beta$ are integers.
The hatted matrices are therefore integral automorphisms.

The exact unit needed in the hatted lattice is
\[
 z=(-v^3+v-1,0,0,2,2(1-v^2))^t.
\]
Writing $a=-v^3+v-1$, the first coordinate of $Y'z$ is
$-2a-2+2v(1-v^2)=0$. Half its second coordinate is
$(1-v)a+(v^2-v)+q(1-v^2)=0$, since the first two terms
sum to $v^4-v^3+v-1$ and the last is its negative. Furthermore
\[
 \beta z=s(v^3-v-1)+2v^2r=1.
 \tag{F.U5}\hypertarget{tag-02IR}{}\label{maps:eq:F.U5}
\]
Indeed the two summands expand to
$v^5+v^4-2v^3-2v^2+1$ and
$-v^5-v^4+2v^3+2v^2$.
Since $AB=I-N$ and $NPz=X(Y'z)=0$, the six-dimensional
map satisfies
\[
 (\widehat A\widehat B-I)(0,z)^t=e_0.
\]
This establishes the unit in the displayed hatted coordinates.

The row identity
$\beta(I+B'+B'^2+B'^3)=0$ proves
$\widehat B^4=I$; also $\widehat A^5=I$. Their quotient actions
show that these orders are exact.

Use the lattice change
\[
 L=\frac14\begin{pmatrix}
 -2&3&-5&-1&5&-3\\
 0&2&-2&2&-2&2\\
 0&-4&0&0&4&-4\\
 0&0&0&-4&4&-4\\
 0&0&4&-4&4&-4\\
 0&0&0&-4&4&0
 \end{pmatrix},\qquad
 A_v=L^{-1}\widehat A L,\quad B_v=L^{-1}\widehat B L.
 \tag{F.U6}\hypertarget{tag-02IS}{}\label{maps:eq:F.U6}
\]
Here $\det L=-1/4$. Expanding these products gives entries in
$\mathbb Z[v]$ for both matrices. The complete expanded entries
and the exact identities
$LA_v=\widehat A L$, $LB_v=\widehat B L$ are given by the
polynomial certificate described in the last item. Thus the
defined matrices act on the stated $\mathbb Z^6$ for every integer
parameter. Their determinants are one and their inverses are
the integer powers $A_v^4,B_v^3$. Their characteristic polynomials are
\[
 (x-1)^2\Phi_5(x),\qquad (x-1)^2(x+1)^2(x^2+1),
\]
respectively.

\item Put $K=B_vA_v-I$. Multiplication gives $K^2=0$.
The following are third minors of $K$, in zero-based indices:
\[
 \begin{array}{c|c|c}
 \text{rows}&\text{columns}&\text{determinant}\\\hline
 (0,1,4)&(1,2,3)&v(v+1)\\
 (0,1,4)&(1,3,4)&v^2\\
 (0,3,4)&(1,2,3)&(v+1)^2.
 \end{array}
\]
Their integer combination $(v+1)^2-2v(v+1)+v^2$ equals one.
Hence $K$ has rank at least three, while $K^2=0$ in dimension six
gives rank at most three. The gcd of its third minors is one,
so its three nonzero Smith factors are all one. This proves
primitivity directly in the lattice, for every $v$.
For $M_1=A_v^t$ and $M_2=B_v^t$ the cusp map is
\[
 M_0=(M_1M_2)^{-1}=(I+K^t)^{-1}=I-K^t.
\]
Thus $M_0-I=-K^t$, which has the same integral Smith factors
and the required maximal square-zero type.

\item Both $A_v,B_v$ fix $e_0$. The quotient morphism to
the original rank-five pair is
\[
 R=P\,\operatorname{pr}_5 L:\mathbb Q^6\longrightarrow\mathbb Q^5,
 \qquad RA_v=AR,\quad RB_v=BR.
\]
It is surjective and has kernel $\mathbb Qe_0$, since
$Le_0=-e_0/2$. We prove that this quotient is irreducible over
$\mathbb Q$. The coprime factors $x-1$ and $\Phi_5(x)$ of
$\chi_A$ give the projector
$(I+A+A^2+A^3+A^4)/5$. Its image is $\mathbb Q\mathbf1$ and
its kernel is $S=\{\sum x_i=0\}$. Irreducibility of $\Phi_5$
shows that these are the only nonzero proper rational
$A$-invariant subspaces.

Neither is $C$-invariant. Indeed,
\[
 N\mathbf1=-2qX\binom{1}{v-1}
\]
has nonzero first entry and zero entries in positions one and three,
so $C\mathbf1=\mathbf1+N\mathbf1$ is not a multiple of
$\mathbf1$. Write $\epsilon=(1,1,1,1,1)$. If $S$ were
$C$-invariant, then $\epsilon N$ would be a multiple of
$\epsilon$; its first entry is zero, so that multiple would
vanish. But its entries in columns three and four are
$v^3-v-1$ and $v^3+v+1$. Their simultaneous vanishing would
give $v=-1$, at which the first is $-1$. This is impossible.
Since $B=C^{-1}$, the pair $(A,B)$ is irreducible, as claimed.

\item The rank-two update gives exact traces
\[
 \operatorname{tr}(AC)=\operatorname{tr}(YA^2X)=1-v^2q,
 \qquad
 \operatorname{tr}((AC)^{-1})
       =-\operatorname{tr}(YA^3X)=1+v^2q.
\]
Here $\operatorname{tr}(A^2)=\operatorname{tr}(A^3)=0$ and
$B=A^4(I-N)$ give the displayed equalities. An equivariant
isomorphism $F$ to the contragredient representation would
satisfy $FA=A^{-t}F$ and $FB=B^{-t}F$, and therefore identify
the character of the word $AC=AB^{-1}$ with the character
of its contragredient matrix. The two displayed traces would
be equal. Their difference $-2v^2q$ is nonzero when $v\ne0$,
proving the required nonselfduality.

For completeness, at $v=0$ the symmetric matrix
$Q_{ij}=1$ for $j-i\equiv1,4\pmod5$ and $Q_{ij}=0$ otherwise
has determinant two and satisfies $A^tQA=Q$, $B_0^tQB_0=Q$.
It gives the equivariant isomorphism in that case.

For the rank-six matrices, the same word has trace
\[
 h(v):=\operatorname{tr}(A_vB_v^{-1})=2-v^2q,
 \qquad h(v+1)-h(v)=-4v^3-3v^2-3v-1.
\]
Thus $h$ is strictly decreasing on the nonnegative integers,
in particular on the nonnegative multiples of twenty.
A simultaneous rational conjugacy of the ordered matrices
would preserve this trace. There are consequently infinitely
many such rational conjugacy classes already in that subfamily.

\item For $v\in20\mathbb Z$, put
$a=(1,3,-2,-1,2,0)^t$, $e=(1,0,0,0,0,0)^t$, and define
\[
 w_1=\frac15\sum_{j=0}^4(A_v^t)^j a,
 \qquad
 w_2=-\frac14\sum_{j=0}^3(B_v^t)^j e.
 \tag{F.U7}\hypertarget{tag-02IT}{}\label{maps:eq:F.U7}
\]
The polynomial matrices reduce to their parameter-zero
values modulo five and four, respectively. At zero,
$M_1a=a$ and $M_2e=e$. The two sums are therefore divisible
by their indicated denominators, so $w_1,w_2$ are integral.
Their defining averages give $M_iw_i=w_i$. The common
first-coordinate covector $\ell$ is fixed by both $M_i$ and
satisfies $\ell(w_1)=1$, $\ell(w_2)=-1$.

For $(m_1,m_2)=(5,4)$ let
$g_i(x)=M_ix+w_i/m_i$ on $\mathbb R^6/\mathbb Z^6$.
Then $g_i^j(x)=M_i^jx+jw_i/m_i$, so $g_i^{m_i}$ is the
identity. Every power with $0<j<m_i$ changes the $\ell$-circle coordinate by $j/5$ or $-j/4$, respectively.
It has no fixed point. These are free actions of the exact
orders five and four, with the affine lifts.

\item Substitution of $v=0$ in the displayed construction gives
exactly the marked pair in Theorem~\ref{maps:src:F.2.5}, and
the averages give $w_1(0)=a$, $w_2(0)=-e$. This identifies
both the matrices and the affine numerators without a further
change of basis.

A separate program computes the expanded polynomial matrices and checks every matrix, quotient, trace
and affine identity above over $\mathbb Q[v]$, with coefficient
integrality in $\mathbb Z[v]$. It checks all four hundred
third minors as well as the displayed integer combination.
Its separate Smith computations at ten stated integer
parameters are additional finite checks; the universal
primitivity assertion follows from the displayed identity
equal to one. The conclusions concern the unrestricted
input tuples in the statement.
\end{enumerate}
\end{proof}




\begin{proposition}[an odd torsion obstruction in this lattice]\hypertarget{tag-02N6}{}\label{maps:src:F.2.17}
\StatementRecord{maps:src:F.2.17}{proved}{02N6}{\eqref{maps:eq:F.U6}, \SeriesRef{maps:src:F.2.19}, \eqref{maps:eq:F.U5}, \SeriesRef{maps:src:F.2.12}}


For the integral tuple $(A_v,B_v)$ in \eqref{maps:eq:F.U6}, any simply
connected filling with free finite affine caps and the fixed
primitive Kuhn cusp has a subgroup
\[
 \mathbb Z/(v^2-v+1)\ \subset\operatorname{Tor}H_2(M;\mathbb Z).
 \tag{F.U9}\hypertarget{tag-02IU}{}\label{maps:eq:F.U9}
\]
In particular, in this specified family of lattices only
$v=0$ and $v=1$ can be three-connected, regardless of the
affine multiplicities and twists.  The stronger exclusion for every stable lattice in each
rational class is proved in Theorem \ref{maps:src:F.2.19} below.

\end{proposition}


\begin{proof}
\begin{enumerate}[label={$\langle1\rangle$\arabic*.}]
\item First compute the global coinvariants in the hatted lattice.
In the rank-five quotient, the $A'$ relations identify
$f_1=f_2=f_3=f_4=t$ and $qf_0=5t$.
The number $q=v^2-v+1$ is relatively prime to five, as its
five residue values show.  Thus these coinvariants are
infinite cyclic, with $f_0=5h$ and $t=qh$.
The $B'$ relations become $q\kappa_j h$, where
\[
 (\kappa_0,\ldots,\kappa_4)=
 (2v^2,\ v(v^3+v-1),\ v(v^3+v+1),
       -v^3+v+1,\ -v^3-v-1).
 \tag{F.U10}\hypertarget{tag-02IL}{}\label{maps:eq:F.U10}
\]
Their gcd is one: it divides $2v$ by subtracting the second
and third entries, and it divides $v^3-v-1$ by the fourth;
the latter is odd and is $-1$ modulo $v$.
Hence the rank-five global coinvariants are $\mathbb Z/q$.
The extension line is killed in rank-six coinvariants by
\eqref{maps:eq:F.U5}, since
$(\widehat A\widehat B-I)(0,z)^t=e$.
Consequently the hatted rank-six coinvariants are also
$\mathbb Z/q$.  The lattice change \eqref{maps:eq:F.U6} is an isomorphism
over every $\mathbb Z_{(p)}$ with $p$ odd, so the odd part of
the rank-six global coinvariants is precisely
$\mathbb Z/q$.

\item Let $V$ be the degree-one fibre cohomology local system.
The only difference between $R^1g_*\mathbb Z$ and $j_*V$
is a finite cokernel at finite caps; the cusp degree-one
specialization is onto.  Thus
\[
 H^2(\mathbb P^1,R^1g_*\mathbb Z)
 =H^2(\mathbb P^1,j_*V)=V_\Gamma,
\]
including its $q$-primary part just computed.  The space of
global degree-two invariants is zero: the irreducible odd
rank-five quotient has no invariant alternating tensor,
and the extension filtration
$0\to W\to\bigwedge^2V\to\bigwedge^2W\to0$
gives $(\bigwedge^2V)^\Gamma=0$.

\item It remains to check every possible source of the differential
$d_2:H^0(\mathbb P^1,R^2g_*\mathbb Z)
\to H^2(\mathbb P^1,R^1g_*\mathbb Z)$.
The cusp has no kernel in degree two.  At a finite cap the
kernel is torsion in $H^2$ of that cap, namely torsion in its
first homology.  Simple connectivity of the total filling
forces the distinguished affine numerator to be prime to the
affine multiplicity, by the determinant argument of
\ref{maps:src:F.2.12}.  The cap is therefore a mapping torus of a torus
automorphism whose order divides the linear order, five or
four.  Its first-homology torsion is supported only at the
primes five or two, respectively: averaging over the finite
order kills the torsion of the vertical coinvariants.
Hence the source of this $d_2$ has no torsion at any prime
dividing $q$, since $q$ is odd and prime to five.

\item The subgroup $\mathbb Z/q$ in bidegree $(2,1)$ survives.
There are no longer differentials on a two-dimensional base,
and this terminal filtration term is a subgroup of
$H^3(M;\mathbb Z)$.  Universal coefficients identify its
torsion with a subgroup of $\operatorname{Tor}H_2(M;\mathbb Z)$.
This proves \eqref{maps:eq:F.U9}.  Finally, $v^2-v+1=1$ for an integer $v$
exactly when $v=0$ or $v=1$.

\end{enumerate}
\end{proof}




\begin{theorem}[A three-connected datum with non-self-dual quotient]\hypertarget{tag-02N7}{}\label{maps:src:F.2.18}
\StatementRecord{maps:src:F.2.18}{proved}{02N7}{\eqref{maps:eq:F.U6}, \SeriesRef{maps:src:F.3.11}, \SeriesRef{maps:src:F.2.5}, \SeriesRef{maps:src:F.2.16}, \SeriesRef{maps:src:F.2.6}, \eqref{maps:eq:F.6}, \SeriesRef{maps:src:F.4.1}}


The parameter $v=1$ in \eqref{maps:eq:F.U6}, with affine orders five and four
and twists
\[
 w_1=(1,3,-2,-1,2,0)^t,\qquad
 w_2=(-1,0,0,0,0,0)^t,
\]
gives a three-connected smooth eight-manifold with
$H_4=\mathbb Z^4$ and all other reduced homology zero except
$H_8=\mathbb Z$.  It is almost complex and is oriented
diffeomorphic to $3\mathbb{HP}^2\#(-\mathbb{HP}^2)$.
Its irreducible rank-five quotient has no invariant bilinear
form, so integral three-connectivity does not force that
quotient to be orthogonal.  By \cref{maps:src:F.3.11} it has holomorphic periods with bounded cusp correction on the punctured base; a compact complex structure on it is not constructed here.

\end{theorem}


\begin{proof}
\begin{enumerate}[label={$\langle1\rangle$\arabic*.}]
\item The first cohomological matrix is the same $A$ as in
Theorem \ref{maps:src:F.2.5}.  The second is
\[
 B_1=\begin{pmatrix}
 1&0&0&0&0&0\\
 0&0&-1&0&1&-1\\
 0&0&-1&0&0&0\\
 0&0&-2&0&1&-2\\
 0&1&-2&0&1&-3\\
 0&0&-1&1&0&-1
 \end{pmatrix}.
 \tag{F.U11}\hypertarget{tag-02IM}{}\label{maps:eq:F.U11}
\]
Take $M_1=A^t$, $M_2=B_1^t$ and
$M_0=(M_1M_2)^{-1}$.  The displayed twists are fixed integral
vectors with $\ell(w_1)=1$, $\ell(w_2)=-1$.
They therefore give free affine actions.  The cusp is
primitive by \ref{maps:src:F.2.16}.  Exact Hermite reduction gives
\[
 \operatorname{col}_{\mathbb Z}
 [\,M_0-I\mid M_1(M_0-I)\mid M_2(M_0-I)\,]
       =\langle e_1,\ldots,e_5\rangle.
\]
Van Kampen leaves only $5a=h$, $4b=-h$, $a+b=0$,
whose relation matrix is unimodular.  This proves
$\pi_1=1$ for the gluing, including zero cusp
clutching.

\item We give the integral data needed to reuse the differential calculation of \ref{maps:src:F.2.6}.  The rational exterior invariant and parabolic ranks
are exactly the table there.  The finite cap first homology
groups are $\mathbb Z^2$ and $\mathbb Z^2\oplus\mathbb Z/2$.
On integral invariant two-form bases their suspension periods
divided by the covering orders are, respectively,
$(-1,0,-1)$ and $(0,0,1)$.  The determinants of the invariant
degree-three skew pairings are $5^4$ and $4^4$, while the
degree-two/degree-four pairings have absolute determinants
$5^3$ and $4^3$.  Thus all the finite covering indices used
in degrees two, three and four are one.

\item For exterior coefficient degrees two, three and four, the
global coboundary Smith form has one factor two and all other
nonzero factors one.  Its torsion class has peripheral
triviality flags $(\mathrm{true},\mathrm{false},\mathrm{true})$,
so the parabolic groups are torsion free, of ranks $0,3,0$.
The degree-one coboundary is primitive and the degree-one
finite covering cokernels are killed by the common invariant
covector and the Chinese remainder theorem, exactly as in
\ref{maps:src:F.2.6}.  The global coinvariants in degrees one through four
are each $\mathbb Z/2$.

\item Write $M_2=\operatorname{diag}(1,D)$.  This time use the
unimodular marking
\[
 P_1=\begin{pmatrix}
 -1&1&-1&0&-1\\
 -1&2&-1&0&-1\\
 -1&0&0&2&-1\\
 1&-1&0&-1&1\\
 0&0&2&-1&1
 \end{pmatrix},\qquad \det P_1=-1.
\]
It satisfies $DP_1=P_1\operatorname{diag}(\text{four-cycle},-1)$.
The cap therefore has the same cellular mapping-torus model
used to prove \eqref{maps:eq:F.6}.  With the unchanged twist $w_2=-e_0$,
the boundary coordinates remain $\theta=a+4k$, $h=k$;
the suspension representative restricts to $w$ in the base
direction and to $-\operatorname{Norm}_4(w)=0$ in the
auxiliary fibre direction.  Consequently the coefficient in
the torsion differential is again one.

\item The new vertical representatives in lexicographic exterior
bases are
\[
\begin{split}
 q=2:\ &(0,-1,-1,-1,0),\\
 q=3:\ &(0,0,1,0,0,0,0,0,0,1),\\
 q=4:\ &(0,0,0,0,0,0,0,-1,-1,0).
\end{split}
\]
Each has residue one in its global coinvariant
$\mathbb Z/2$.  Thus all three torsion differentials in \eqref{maps:eq:F.6}
are isomorphisms.  The integral Leray table and its free
degree-one differential are therefore exactly those of
\ref{maps:src:F.2.6}.  They give $H^2=H^3=0$ and $H^4=\mathbb Z^4$.
Universal coefficients, Hurewicz and Poincar\'e duality give
the stated integral homology and three-connectivity.

\item The relative almost-complex extension over the pair of pants
in \ref{maps:src:F.2.5} applies: the finite matrices have the same even
real eigenspace dimensions and the cusp is the same marked
Kuhn type.  The purely smooth characteristic-vector and
normal-invariant argument in \ref{maps:src:F.4.1} then gives the stated
oriented diffeomorphism type.  No integrability at the finite
punctures is used for this smooth classification.

\item Finally, Theorem~\ref{maps:src:F.2.16} gives nonselfduality of the irreducible rank-five quotient at $v=1$.
The exact program
\texttt{fourfolds\_v1\_homology.py} runs every matrix,
Smith, Hermite, suspension-period, pairing and differential
certificate just displayed.  The proof uses the explicit monomial marking displayed above.

\end{enumerate}
\end{proof}




\begin{theorem}[a residual obstruction for every lattice]\hypertarget{tag-02N8}{}\label{maps:src:F.2.19}
\StatementRecord{maps:src:F.2.19}{proved}{02N8}{\SeriesRef{maps:src:F.2.16}, \eqref{maps:eq:F.U6}, \SeriesRef{maps:src:F.2.6}, \SeriesRef{maps:src:F.2.18}, \eqref{maps:eq:F.U12}, \SeriesRef{maps:src:F.2.17}, \eqref{maps:eq:F.U1}}


Let a rational rank-six datum have finite linear orders five
and four, a primitive maximal cusp on a stable integral lattice,
and invariant line $L$, with irreducible rank-five quotient $W$
and no rational global coinvariants.  Suppose that at a prime
$p\ne2,5$ the semisimplified reduction of a stable lattice in
$W$ is the sum of the trivial module and an absolutely
irreducible module of dimension four.  Then no stable integral
lattice in that rational datum, with primitive maximal cusp,
and no free finite affine twists can give
$\pi_1=1$ and $H_2=0$ simultaneously.

For every $v\notin\{0,1\}$, the rational family in \ref{maps:src:F.2.16}
satisfies this hypothesis at each prime dividing
$v^2-v+1$.  Consequently the exclusion of those parameters
holds for \emph{all} stable lattices and all twists, not just
the lattice in \eqref{maps:eq:F.U6}.  The two remaining parameters give the three-connected manifolds of \cref{maps:src:F.2.6,maps:src:F.2.18}.

\end{theorem}


\begin{proof}
\begin{enumerate}[label={$\langle1\rangle$\arabic*.}]
\item We first give the lattice-independent mechanism.  For any
stable lattice $V_\mathbb Z$, use its primitive invariant
line $L_\mathbb Z$ and quotient lattice $W_\mathbb Z$.
The semisimplified reduction of a rational representation is
independent of its stable lattice.  One elementary way to see
this is to compare two commensurable lattices by their finite
quotient $T$.  The exact sequence obtained by reduction modulo
$p$ has end terms $T[p]$ and $T/pT$, which have the same
composition factors by the filtration of $T$ by powers of $p$.
Hence the two lattice reductions have equal classes in the
Grothendieck group of finite-length modules.  Chaining such
comparisons proves the assertion for any two lattices.

\item Thus $W_\mathbb Z/pW_\mathbb Z$ has length two, with one
trivial constituent and one irreducible four-dimensional
constituent.  It has either a trivial submodule or a trivial
quotient (and possibly both).

\item Suppose first that it has a trivial submodule.  Its inverse
image in $V_\mathbb Z/pV_\mathbb Z$ is an extension of a
trivial module by a trivial module.  Each finite generator
acts on this two-dimensional space by a unipotent matrix.
Since its order is five or four, prime to $p$, it acts as
the identity.  There are therefore at least two global
invariant vectors in degree-one fibre cohomology modulo $p$.
Write their dimension as $a_p\ge2$; this is also the
dimension of the global coinvariants of fibre homology over
$\mathbb F_p$, by the perfect pairing between fibre homology
and cohomology.

\item The abelianized van Kampen presentation now shows
\[
 b_1(M;\mathbb F_p)\ge a_p-1.
 \tag{F.U12}\hypertarget{tag-02IN}{}\label{maps:eq:F.U12}
\]
Indeed, after the monodromy relations the fibre contributes
$a_p$ dimensions, and the two finite meridians contribute
two more.  The two finite power relations and the single
base relation impose at most three linear relations.
The cusp introduces no further fibre relation in this
abelianization.  On the dual homology lattice its nilpotent
operator is still primitive and maximal, so its angular lattice is
$\ker(M_0-I)=\operatorname{im}(M_0-I)$ by primitivity, and
the latter image already vanishes among the global
monodromy relations.  Its base meridian is the base relation
just counted.  Changing affine twists changes coefficients
of these three relations, not their number.  Thus \eqref{maps:eq:F.U12}
contradicts simple connectivity.

\item Suppose instead that the reduction of $W$ has a trivial
quotient.  Its integral global coinvariants have nonzero
$p$-primary torsion, since they have rational rank zero.
The natural map from the global coinvariants of $V$ to
those of $W$ is an isomorphism.  To check its kernel, a
primitive generator $e$ of $L_\mathbb Z$ lies in
$\ker(M_0-I)=\operatorname{im}(M_0-I)$, and hence is already
zero among global coinvariants.  Therefore $H^2(R^1)$ has
nonzero $p$-primary torsion.  If the total space were simply
connected, the finite-cap argument of \ref{maps:src:F.2.17} would make
every possible torsion source in $H^0(R^2)$ supported only
at two and five.  The invariant part of $\bigwedge^2V$ is
zero by irreducibility of the odd-dimensional quotient.
Consequently the nonzero $p$-primary subgroup survives to
$H^3(M;\mathbb Z)$ and then to
$\operatorname{Tor}H_2(M;\mathbb Z)$, again a contradiction.
This proves the general criterion.

\item For its application to \ref{maps:src:F.2.16}, let $p\mid q=v^2-v+1$.
We have $p\ne2,5$.  In the original cycle lattice of \eqref{maps:eq:F.U1},
$Y\mathbf1=2q(-1,1-v)^t$, so $\mathbf1$ gives a trivial
submodule after reduction.  On the four-dimensional quotient
the first finite generator has characteristic polynomial
$\Phi_5$, and $C=B^{-1}$ has polynomial
$(x+1)^2(x^2+1)$.  This quotient is absolutely irreducible.
Here is a characteristic-independent proof over
$\overline{\mathbb F}_p$.

\item A one-dimensional composition factor would make the
eigenvalues of $A$ and $C$ equal, because $CA^{-1}$ is
unipotent.  An eigenvalue of order dividing both five and
four must be one, which is absent from $\Phi_5$.
If the quotient were reducible, its only remaining possible
composition-factor dimensions would be two and two.
On each factor, the determinants of $A$ and $C$ must agree
and have orders dividing five and four, so both determinants
are one.  The two $C$ eigenvalue pairs must therefore be
$(-1,-1)$ and $(i,-i)$ in some order.  On the former
factor $C=-I$, since its order is prime to $p$.
Then the unipotence of $CA^{-1}=-A^{-1}$ would force the
eigenvalues of $A$ to be $-1$, impossible for $\Phi_5$.
This proves absolute irreducibility and establishes the
required residual composition factors at every $p\mid q$.
When $v\notin\{0,1\}$, the integer $q$ is greater than one
and has such a prime divisor.  The criterion applies.

\end{enumerate}
\end{proof}




\begin{theorem}[the complete rational $(5,4)$ normal form]\hypertarget{tag-02NA}{}\label{maps:src:F.2.20}
\StatementRecord{maps:src:F.2.20}{proved}{02NA}{\SeriesRef{maps:src:F.2.11}, \eqref{maps:eq:F.R4}, \eqref{maps:eq:F.R6}, \eqref{maps:eq:F.R7}, \eqref{maps:eq:F.R8}}


Consider ordered rational rank-five pairs $(A,C)$ with
\[
 \chi_A=(x-1)\Phi_5,
 \qquad \chi_C=(x-1)(x+1)^2(x^2+1).
\]
Put
\[
 N=CA^{-1}-I,
 \qquad N^2=0,\quad\operatorname{rank}N=2,
\]
where $A^5=C^4=I$ and the pair is irreducible over $\mathbb Q$.
Their simultaneous rational conjugacy classes are parametrized
bijectively by
\[
 \bigl\{[r:s:t]\in\mathbb P^2(\mathbb Q):D(r,s,t)\ne0\bigr\}
 \ \sqcup\ \mathbb A^1(\mathbb Q),
 \tag{F.R1}\hypertarget{tag-02HJ}{}\label{maps:eq:F.R1}
\]
where
\[
\begin{split}
 D={}&r^3-4r^2s-r^2t+7rs^2+rst-3rt^2\\
     &{}-7s^3+11s^2t-11st^2+7t^3.
\end{split}
 \tag{F.R2}\hypertarget{tag-02HP}{}\label{maps:eq:F.R2}
\]
The second summand lies over the single exceptional point
$[4:1:3]$ of $D=0$.  The construction below is explicit and
uses only a ten-variable linear system.  By \ref{maps:src:F.2.11} it also
parametrizes the rational rank-six maximal-cusp extensions
of all these pairs, with $B=C^{-1}$.
This is a rational classification; the stable integral lattices, the affine twists and the homology of the gluings are treated below.

\end{theorem}


\begin{proof}
\begin{enumerate}[label={$\langle1\rangle$\arabic*.}]
\item Write $K=\mathbb Q(\zeta)$, $\zeta^4+\zeta^3+\zeta^2+
\zeta+1=0$, and identify the rational space with
$\mathbb Q\oplus K$.  The matrix of $A$ is multiplication by
$(1,\zeta)$:
\[
 A=\operatorname{diag}(1,J),\qquad
 J=\begin{pmatrix}
 0&0&0&-1\\1&0&0&-1\\0&1&0&-1\\0&0&1&-1
 \end{pmatrix},\qquad R=(A+I)^{-1}.
 \tag{F.R3}\hypertarget{tag-02HQ}{}\label{maps:eq:F.R3}
\]
Its rational centralizer is $\mathbb Q^*\times K^*$.
Let $I_N=\operatorname{im}N$.  If $X$ is any basis matrix
for this plane and $N=XY$, then $YX=0$.
The identity
\[
 C+I=(A+I)+XYA
\]
shows that its two-dimensional kernel has the same dimension
as the kernel of
$I_2+YARX=I_2-YRX$.  Hence
\[
 YRX=I_2,
 \qquad I_N\cap RI_N=0.
 \tag{F.R4}\hypertarget{tag-02HR}{}\label{maps:eq:F.R4}
\]
Conversely the two matrix equations in \eqref{maps:eq:F.R4}, together with
$YX=0$, give precisely that two-dimensional $-1$ eigenspace.

\item The plane $I_N$ projects nontrivially to $\mathbb Q$:
otherwise the first-coordinate covector is invariant under
both $A$ and $C$, contrary to irreducibility.
It contains no nonzero vector of $\mathbb Q\oplus0$,
because $R$ acts as $1/2$ on that line and \eqref{maps:eq:F.R4} would fail.
Write a basis of the plane as $(0,u),(1,v)$, where
$u\in K^*$.  Multiplication on $K$ by $u^{-1}$ puts it into
\[
 I_N=\langle(0,1),(1,\beta)\rangle_{\mathbb Q}.
\]
Changing the second basis vector by a rational multiple of
the first changes $\beta$ by a rational constant.
Two such planes are in the same centralizer orbit exactly
when $\beta'=a\beta+b$, with $a\in\mathbb Q^*$ and
$b\in\mathbb Q$.  Indeed a centralizer element carrying
one to the other must preserve the kernel of the projection
to $\mathbb Q$, so its $K^*$ component is itself rational;
the displayed affine change follows, and every such change
is realized.  We may therefore take
\[
 \beta=r\zeta+s\zeta^2+t\zeta^3,
 \qquad [r:s:t]\in\mathbb P^2(\mathbb Q),
 \qquad
 X=\begin{pmatrix}0&1\\1&0\\0&r\\0&s\\0&t\end{pmatrix}.
 \tag{F.R5}\hypertarget{tag-02HS}{}\label{maps:eq:F.R5}
\]
The triple is nonzero; a rational $\beta$ would make the
plane contain the forbidden first-coordinate line.
The stabilizer of such a plane in the centralizer consists
only of common scalar matrices.  Thus it introduces no
further identifications of the resulting nilpotent operators.

\item Write $Y$ as an unknown two-by-five matrix.  It is determined
by the following ten linear equations:
\[
 YX=0,\qquad YRX=I_2,\qquad
 \operatorname{tr}(YAX)=-1,\qquad
 \operatorname{tr}(YA^{-1}X)=1.
 \tag{F.R6}\hypertarget{tag-02HT}{}\label{maps:eq:F.R6}
\]
The first eight are \eqref{maps:eq:F.R4}.  The last two are exactly
$\operatorname{tr}C=\operatorname{tr}C^{-1}=-1$, since
$\operatorname{tr}A=\operatorname{tr}A^{-1}=0$ and
$C^{-1}=A^{-1}(I-XY)$.
They suffice to obtain the required characteristic
polynomial: the two $-1$ roots and determinant one leave a
monic cubic of constant term $-1$, and the two traces fix
that cubic to $(x-1)(x^2+1)$.
These remaining three roots are distinct, and the $-1$
eigenspace already has dimension two, so $C$ is semisimple
and $C^4=I$.  Also $YRX=I_2$ ensures that $N$ has rank two.
Thus every solution of \eqref{maps:eq:F.R6} gives all the asserted local
linear conditions.

\item In row-major order on the ten entries of $Y$, the determinant
of the coefficient matrix in \eqref{maps:eq:F.R6} is $D/4$.
For $D\ne0$ this proves uniqueness and gives the explicit
formula $Y=\operatorname{reshape}(E^{-1}b)$, where $E$ is
the coefficient matrix of the displayed linear equations
and $b$ their displayed right side.
The exact expanded matrix and all ten rational functions
are printed by \texttt{fourfolds\_relative\_normalform.py}.
There is no search or degree cutoff in this calculation.

\item We now treat the exceptional locus.
Set
\[
 Q_1=r^2-rs-rt+s^2+t^2,
 \qquad Q_2=r^2-2rs-3rt+4s^2-4st+4t^2.
\]
The first column of the unique generic solution is
\[
 y_{10}=\frac{2Q_1Q_2}{D},\qquad
 y_{20}=\frac{2(r-s-t)Q_1}{D}.
 \tag{F.R7}\hypertarget{tag-02HU}{}\label{maps:eq:F.R7}
\]
Here $y_{10},y_{20}$ are the first column of $Y$.
The two equalities after multiplication by $D$ are polynomial
linear combinations of \eqref{maps:eq:F.R6}, by its adjugate matrix.
The program prints the corresponding two polynomial rows
and verifies those identities, so they remain necessary
also when $D=0$.
Moreover
\[
 2Q_1=(r-s)^2+(r-t)^2+s^2+t^2>0
\]
for every nonzero rational triple.  Consequently consistency
on $D=0$ forces $r=s+t$ and $Q_2=0$.
Substitution gives
\[
 D(s+t,s,t)=(t-3s)(2t-s)^2,
 \qquad Q_2(s+t,s,t)=(3s-t)(s-2t).
 \tag{F.R8}\hypertarget{tag-02HV}{}\label{maps:eq:F.R8}
\]
There are precisely two resulting projective points,
$[3:2:1]$ and $[4:1:3]$.
At the former, $[X\mid RX]$ has rank less than four, so
\eqref{maps:eq:F.R4} is impossible.  At the latter the coefficient and
augmented matrices of \eqref{maps:eq:F.R6} both have rank nine, and every
solution is
\[
 Y_z=\begin{pmatrix}
 4z+16/5&0&-2z-1&4/5-2z&2z\\
 2z-2/5&0&-z&2/5-z&z
 \end{pmatrix},\qquad z\in\mathbb Q.
 \tag{F.R9}\hypertarget{tag-02HW}{}\label{maps:eq:F.R9}
\]
Direct substitution proves both necessity and sufficiency.
This establishes the exceptional affine line in \eqref{maps:eq:F.R1}.

\item Every pair just constructed is irreducible over $\mathbb Q$.
An $A$-stable proper rational subspace has dimension one or
four.  A one-dimensional subrepresentation or quotient must
have trivial character, since its finite eigenvalues of
orders dividing five and four must agree under the
unipotent product.  The only $A$-invariant covector is the
first-coordinate covector.  Its product with $N$ is the
second row of $Y$, which is nonzero by $YRX=I_2$.
The only $A$-invariant vector is the first-coordinate vector;
its image under $N$ vanishes exactly if both entries in
\eqref{maps:eq:F.R7} vanish.  For $D\ne0$, positivity of $Q_1$ and
\eqref{maps:eq:F.R8} rule this out.  On the exceptional line the two
entries $4z+16/5$ and $2z-2/5$ never vanish simultaneously.
There are therefore neither invariant vectors nor covectors,
which excludes both possible dimensions.

\item Finally, rescaling $(r,s,t)$ is exactly the centralizer
change already described; it gives a conjugate pair.
Conversely any conjugacy of two normalized pairs centralizes
$A$ and hence identifies their projective parameters as proved
above.  At the exceptional point its stabilizer is scalar,
so distinct $z$ give distinct pairs.  This proves the stated
bijection, including completeness and uniqueness.
The saved program checks all determinant, adjugate,
characteristic-polynomial and finite-order identities over
the rational function field, and checks the exceptional
line symbolically for every $z$.

\end{enumerate}
\end{proof}




\begin{corollary}[all good-prime lattices are forced]\hypertarget{tag-02NB}{}\label{maps:src:F.2.21}
\StatementRecord{maps:src:F.2.21}{proved}{02NB}{\SeriesRef{maps:src:F.2.20}, \SeriesRef{maps:src:F.2.19}}


For a $(5,4)$ datum with $\pi_1=1$ and $H_2=0$, the
rank-five quotient has absolutely irreducible reduction
at every prime $p\ne2,5$.  At each such prime its stable
$\mathbb Z_p$-lattice is unique up to homothety.
The same is true of the rank-six extension among lattices
with primitive maximal cusp.  Thus, once a rational pair
in \ref{maps:src:F.2.20} admits one three-connected datum, all remaining integral-lattice
ambiguity is supported at the primes two and five.


\end{corollary}


\begin{proof}
\begin{enumerate}[label={$\langle1\rangle$\arabic*.}]
\item Work first over an algebraic closure of $\mathbb F_p$.
Both finite generators are semisimple because $p\ne2,5$.
A one-dimensional composition factor must be trivial:
its $A$ and $C$ eigenvalues agree by the unipotent product,
and their orders divide five and four.
There can be at most one such factor, since the $1$
eigenvalue of $A$ occurs once.
If it occurs, the remaining four-dimensional module is
absolutely irreducible by the spectral argument in \ref{maps:src:F.2.19}.
The only other possible reducible composition pattern is
$2+3$.
On each factor, determinants of $A$ and $C$ agree and
have orders dividing five and four, hence are one.
The two-dimensional $C$ eigenvalue pair is therefore
$(-1,-1)$ or $(i,-i)$.
The first would give $C=-I$ on that factor, contradicting
the unipotence of $CA^{-1}$ and the absence of $-1$
among the eigenvalues of $A$.
In the second case the three-dimensional factor has
$C$ eigenvalues $1,-1,-1$, so $C^2=I$ there; it has
an $A$-fixed vector because the other factor has determinant
one and cannot contain the simple $1$ eigenvalue of $A$.

\item Here is the required elementary involutory identity in the
present characteristic.  Put $B=A^{-1}$ on this factor.
From $(CB-I)^2=0$ and $C^2=I$ one obtains
\[
 BC+CB^{-1}=2I.
\]
For $Bv=v$ this says $(B+I)Cv=2v$.
The operator $B+I$ is invertible, since $B^5=I$ and
$p\ne2$, so $Cv=v$.  The resulting common fixed line
contradicts simplicity of the three-dimensional factor.
We have proved that any reducible residual representation
has exactly the constituents $1$ and an absolutely
irreducible four-dimensional module.  Theorem \ref{maps:src:F.2.19}
excludes that possibility for a three-connected datum, proving absolute
irreducibility at every indicated prime.

\item Let $R=\mathbb Z_p$ and $K_p=\mathbb Q_p$.
For two stable quotient lattices, multiply one by a power
of $p$ until it is contained in the other but not in its
$p$-multiple.  Its nonzero residual image must then be the
whole quotient by absolute irreducibility.  Nakayama's
lemma, or the Smith form of the inclusion, makes the
lattices equal.  This proves quotient-lattice uniqueness.

\item For completeness the extension-lattice comparison is also
integral.  Fix a primitive-cusp lattice $V_R$, with invariant
line $Re$ and quotient $W_R$.
The coboundary map
\[
 W_R^*\longrightarrow W_R^*\oplus W_R^*
\]
is primitive, because its reduction has trivial kernel.
Thus its cokernel, the degree-one cohomology of the free
group on the two generators, is torsion free, and an
integral coboundary over $K_p$ is already integral over $R$.
The extension class $\alpha$ is primitive in this cokernel.
Indeed, divisibility by $p$ would make the extra cusp map
\[
 \ker N_W/\operatorname{im}N_W\longrightarrow Re
\]
divisible by $p$, contrary to the three unit Smith factors
of the full cusp.  Here $N_W$ has primitive image:
it is the quotient of the primitive image of $N_V$ by
$Re\subset\operatorname{im}N_V$.

\item Normalize any other stable lattice by the already proved
quotient homothety so that its quotient is $W_R$.
Its invariant line is $p^rRe$ for some integer $r$.
The extension class in that basis is $p^{-r}\alpha$;
primitivity makes it integral only when $r\leq0$.
Primitivity of the coboundary map eliminates any additional
rational change of splitting.  The extra cusp map is then
multiplied by $p^{-r}$, so its Smith factor is a unit exactly
when $r=0$.  The lattices coincide after the common homothety.
This proves the assertion.

\end{enumerate}
\end{proof}




\begin{theorem}[a terminating lattice algorithm for each rational pair]\hypertarget{tag-02NC}{}\label{maps:src:F.2.22}
\StatementRecord{maps:src:F.2.22}{proved}{02NC}{\SeriesRef{maps:src:F.2.14}, \eqref{maps:eq:F.L3}, \eqref{maps:eq:F.L5}, \SeriesRef{maps:src:F.2.11}, \eqref{maps:eq:F.L6}, \eqref{maps:eq:F.L7}}


Let $(A,C)$ be any irreducible rational rank-five pair in
one of the eight spectral cases of \ref{maps:src:F.2.14}, with
$A^5=I$, $B=C^{-1}$, and $(AB-I)^2=0$ of rank two.
There is an explicit terminating algorithm which decides
whether the pair preserves an integral lattice, and, if so,
lists all such lattices up to simultaneous integral conjugacy.
It then lists all rank-six extensions with the same finite
orders and primitive maximal square-zero cusp, up to
simultaneous integral conjugacy.  The bounds below depend on the input rational pair and lose no lattice.

\end{theorem}


\begin{proof}
\begin{enumerate}[label={$\langle1\rangle$\arabic*.}]
\item We give the algorithm and its termination proof.
The fixed line of $A$ has dimension one.  Choose a nonzero
vector $e$ on it and a nonzero fixed covector $\varphi$.
Irreducibility implies that the vectors $ge$, with $g$ a
word in $A,C$, span the whole rational space; the analogous
statement holds for the covectors $\varphi h$.
In fact words of length at most four suffice.  The spans of
words of lengths at most $j$ increase from dimension one;
if one step does not increase the span, it is invariant
under both generators and hence already the whole space.
There are only four possible increases.

\item Choose words $g_1,\ldots,g_5$ and $h_1,\ldots,h_5$ giving
bases in this way, and put
\[
 P=I+A+\cdots+A^4,\qquad T_{ij}=g_iPh_j.
 \tag{F.L1}\hypertarget{tag-02GT}{}\label{maps:eq:F.L1}
\]
The rank-one operator $P$ is a nonzero scalar multiple of
$e\varphi$.  Hence the 25 matrices $T_{ij}$ are a rational
basis of $\operatorname{Mat}_5(\mathbb Q)$.
Each belongs to $\mathcal R=\mathbb Z[A,C]$.
Let $\mathcal R_0$ be their integral span and set
\[
 \mathcal R_0^\#=
 \{X:\operatorname{tr}(XT)\in\mathbb Z
                     \text{ for all }T\in\mathcal R_0\}.
 \tag{F.L2}\hypertarget{tag-02GV}{}\label{maps:eq:F.L2}
\]
The trace pairing on the full matrix algebra is
nondegenerate, as is seen from the matrix units, so this
is another full rational lattice, computed by inverting
the 25-by-25 Gram matrix.

\item If any stable lattice exists, every element of $\mathcal R$
acts integrally on it.  Consequently
$\operatorname{tr}(XY)\in\mathbb Z$ for all
$X,Y\in\mathcal R$, and necessarily
\[
 \mathcal R_0\subset\mathcal R\subset\mathcal R_0^\#.
 \tag{F.L3}\hypertarget{tag-02GW}{}\label{maps:eq:F.L3}
\]
A nonintegral Gram entry therefore proves nonexistence.
Otherwise let
$d=[\mathcal R_0^\#:\mathcal R_0]$, the absolute Gram
determinant.  Start with the lattice generated by
$\mathcal R_0$ and $I$, and repeatedly replace a current
lattice $L$ by $L+AL+CL$, using exact Hermite reduction.
If it leaves $\mathcal R_0^\#$, \eqref{maps:eq:F.L3} proves that there
is no stable lattice.
While it remains inside, every strict enlargement reduces
its index by a factor at least two.
There can be at most $\lfloor\log_2d\rfloor$ such
enlargements starting from $\mathcal R_0$, so a bound of
$1+\lfloor\log_2d\rfloor$ tests after adjoining $I$ is
sufficient.  Stabilization gives exactly $\mathcal R$:
the stable lattice contains $I$ and all words in $A,C$,
and every element introduced belongs to their ring.
It is then an order, and
$\Lambda_0=\mathcal R\mathbb Z^5$ is an explicitly
computed stable lattice.  This proves the decision step,
including termination in the nonexistence case.

\item Use a basis of $\Lambda_0$ from now on, so
$\mathcal R\subset\operatorname{Mat}_5(\mathbb Z)$.
Smith reduction gives an integer $c\ge1$ with
\[
 c\operatorname{Mat}_5(\mathbb Z)\subset\mathcal R;
 \tag{F.L4}\hypertarget{tag-02GX}{}\label{maps:eq:F.L4}
\]
one may take the largest Smith factor of the inclusion
of the order in the full matrix lattice.
Normalize any stable lattice by a rational homothety so
that it is contained in $\mathbb Z^5$ and the gcd of all
coordinates of a lattice basis is one.
The operators $cE_{ij}$ preserve it.  The coordinate values
of its vectors generate $\mathbb Z$, so Bezout's identity,
applied to $cE_{ij}$, shows that it contains every $ce_i$.
Therefore every normalized lattice lies in the finite interval
\[
 c\mathbb Z^5\subset\Lambda\subset\mathbb Z^5.
 \tag{F.L5}\hypertarget{tag-02GY}{}\label{maps:eq:F.L5}
\]
Conversely every lattice in this interval stable under
$A,C$ is retained by the algorithm.
The only rational operators commuting with the pair are
scalars: they act by a scalar on the one-dimensional fixed
line of $A$, and then by the same scalar on the spanning
vectors $g_ie$.  Thus rational equivalence of two stable
lattices is precisely homothety.  The coordinate-gcd
normalization is unique, so \eqref{maps:eq:F.L5} gives the integral
conjugacy classification without further identifications.

\item Here is a finite enumeration bound which is also implemented.
For each $p^a\Vert c$, enumerate all subspaces of
$\mathbb F_p^5$ in reduced echelon form.
Beginning with $\mathbb Z^5$, pass through $a$ levels of
stable lattices between $pH$ and $H$, retaining only those
containing $p^a\mathbb Z^5$ and identifying equal Hermite
forms.  Every desired lattice occurs: for any target
$\Lambda$, the chain
$\Lambda+p^j\mathbb Z^5$, $0\le j\le a$, is such a
chain and ends at $\Lambda$.
If $s_p$ denotes the number of all subspaces of
$\mathbb F_p^5$, there are at most $s_p^a$ branches before
deduplication.  Intersect the resulting lattices for the
different primes and discard those of coordinate gcd
greater than one.  The Chinese remainder theorem proves
that this enumerates \eqref{maps:eq:F.L5} completely.

\item For each quotient lattice $W_{\mathbb Z}$ so obtained,
put $N_W=AB-I$.
Its image must be primitive: for a primitive full cusp,
$\operatorname{im}N_V=\ker N_V$ contains the primitive
invariant line, and quotienting that image by the line
gives a primitive image in $W_{\mathbb Z}$.
Reject a quotient whose two nonzero Smith factors are
not one.  For the remaining quotients use extension rows
$c_1,c_2\in\mathbb Z^{1\times5}$ and impose
\[
\begin{gathered}
 c_1(I+A+\cdots+A^4)=0,
 \qquad c_2(I+B+\cdots+B^{m-1})=0,\\
 (c_2+c_1B)N_W=0,
 \qquad m=\operatorname{ord}(B).
\end{gathered}
 \tag{F.L6}\hypertarget{tag-02GZ}{}\label{maps:eq:F.L6}
\]
These are integral linear equations.  Their integral kernel
modulo the coboundaries
\[
 z\longmapsto\bigl(z(A-I),z(B-I)\bigr),
 \qquad z\in\mathbb Z^{1\times5},
 \tag{F.L7}\hypertarget{tag-02H0}{}\label{maps:eq:F.L7}
\]
is a finitely generated group $G$ of rank one by \ref{maps:src:F.2.11}.
Compute its free and torsion factors by Smith reduction.
Choose an integral generator $w$ of the free rank-one
quotient $\ker N_W/\operatorname{im}N_W$.
The homomorphism
\[
 \kappa:G\longrightarrow\mathbb Z,
 \qquad [(c_1,c_2)]\longmapsto(c_2+c_1B)w
 \tag{F.L8}\hypertarget{tag-02H1}{}\label{maps:eq:F.L8}
\]
is well defined and nonzero: \eqref{maps:eq:F.L6} kills the image, and
\eqref{maps:eq:F.L7} kills $w$; rational nonvanishing is \ref{maps:src:F.2.11}.
The full cusp has unit Smith factors exactly when
$\kappa=\pm1$.
Indeed the quotient rank-two image is already primitive,
and the remaining invariant-line image is generated by
that one integer.
If $\operatorname{im}\kappa\ne\mathbb Z$, reject the
lattice.  Otherwise the inverse image of $1$ is a finite
coset of $\operatorname{Tor}G$, and lists all extensions.
Changing the sign of the invariant-line generator identifies
$\kappa=-1$ with $\kappa=1$.
There are no further identifications: an integral isomorphism
on the quotient is $\pm I$, on the invariant line is
$\pm1$, and, after $\kappa=1$ is fixed, their ratio is one.
The remaining row is exactly an integral change of splitting
\eqref{maps:eq:F.L7}.
This proves completeness and termination of the extension
step as well.

\item The order, trace-dual bound, conductor and finite-interval
steps are implemented in
\texttt{fourfolds\_lattice\_algorithm.py};
the extension-kernel and Smith steps are implemented in
\texttt{fourfolds\_two\_classes\_lattices.py} with the
finite order as an argument.  Their applications and exact
exhaustive bounds are recorded next.

\end{enumerate}
\end{proof}




\begin{theorem}[Lattice classes of the two three-connected members]\hypertarget{tag-02ND}{}\label{maps:src:F.2.23}
\StatementRecord{maps:src:F.2.23}{proved}{02ND}{\SeriesRef{maps:src:F.2.16}, \SeriesRef{maps:src:F.2.6}, \SeriesRef{maps:src:F.2.18}, \eqref{maps:eq:F.U1}, \eqref{maps:eq:F.L4}, \eqref{maps:eq:F.L5}, \eqref{maps:eq:F.L6}, \eqref{maps:eq:F.L7}, \SeriesRef{maps:src:F.2.12}, \eqref{maps:eq:F.L10}, \eqref{maps:eq:F.U6}, \SeriesRef{maps:src:F.2.19}}


For each of the two rational representations $v=0$ and $v=1$
of \ref{maps:src:F.2.16}, there are exactly two homothety classes of stable
rank-five quotient lattices, and exactly three integral
conjugacy classes of rank-six extensions with primitive
maximal cusp.  Exactly one of those three can have simply
connected finite affine fillings.  It is the integral class
already realized in \ref{maps:src:F.2.6}, respectively \ref{maps:src:F.2.18}.
Thus arbitrary stable-lattice changes produce no further
linear classes with three-connected gluing anywhere in the family \ref{maps:src:F.2.16}.
The affine twists and multiplicities are treated in \cref{maps:src:F.2.24}.

\end{theorem}


\begin{proof}
\begin{enumerate}[label={$\langle1\rangle$\arabic*.}]
\item Apply the preceding algorithm in the original five-cycle
lattice \eqref{maps:eq:F.U1}.  For both parameters the order is obtained
after four strict enlargements.  Its Smith factors in the
full matrix lattice are 21 factors one and four factors two;
thus the conductor $c=2$ in \eqref{maps:eq:F.L4} is exact.
All 374 subspaces of $\mathbb F_2^5$, including zero and
the whole space, are enumerated, with no height restriction
or neighbor-depth assumption.
After imposing stability and coordinate gcd one, the two
lattice representatives are
\[
 \mathbb Z^5,\qquad
 \Lambda_* =2\mathbb Z^5+\mathbb Z(1,1,1,1,1)^t.
 \tag{F.L9}\hypertarget{tag-02H2}{}\label{maps:eq:F.L9}
\]
The latter has index $16$.
The generic prime-by-prime implementation of \eqref{maps:eq:F.L5} and
an independent bit-vector enumeration give exactly the
same two Hermite forms.

\item For $\mathbb Z^5$, the extension group in \eqref{maps:eq:F.L6}--\eqref{maps:eq:F.L7}
is $\mathbb Z$ and $\kappa$ is a unit on its generator.
For $\Lambda_*$, that group is
$\mathbb Z\oplus\mathbb Z/2$, again with primitive
$\kappa$.  Hence there are exactly one and two primitive
extensions, respectively.
For each extension the program computes the image of the
primitive global invariant homology covector $\ell$ on
the two finite fixed lattices.  The resulting positive
divisibilities are
\[
\begin{array}{c|c|c}
 \text{quotient lattice}&\text{extension class}&
       (\operatorname{div}_1,\operatorname{div}_2)\\\hline
 \mathbb Z^5&\text{unique}&(1,2)\\
 \Lambda_*&\text{first torsion coset}&(1,1)\\
 \Lambda_*&\text{second torsion coset}&(1,2).
\end{array}
 \tag{F.L10}\hypertarget{tag-02GU}{}\label{maps:eq:F.L10}
\]
The table is identical for $v=0$ and $v=1$.
All six full matrices, the integral kernel bases, the
Smith factors and the finite divisibilities are computed
exactly by \texttt{fourfolds\_two\_classes\_lattices.py}
.

\item For a simply connected filling the determinant quotient
of \ref{maps:src:F.2.12} requires $\gcd(\ell(v_i),m_i)=1$ at each
finite cap.  The second affine multiplicity is divisible
by four.  A divisibility of two in \eqref{maps:eq:F.L10} therefore
excludes that class for every choice of twists and larger
multiplicities.
The sole remaining class has an explicit integral
intertwiner of determinant $\pm1$ with the corresponding
pair in \eqref{maps:eq:F.U6}, as printed and checked by the program.
Those pairs already have the verified simply connected,
three-connected fillings.  \ref{maps:src:F.2.19} excludes every other
rational parameter for all lattices.
Together these facts prove all the claims, including the
absence of any additional linear class with three-connected gluing in the
unbounded family.

\end{enumerate}
\end{proof}




\begin{theorem}[All three-connected twists in the unbounded family]\hypertarget{tag-02NE}{}\label{maps:src:F.2.24}
\StatementRecord{maps:src:F.2.24}{proved}{02NE}{\SeriesRef{maps:src:F.2.12}, \SeriesRef{maps:src:F.2.16}, \SeriesRef{maps:src:F.2.23}, \SeriesRef{maps:src:F.2.20}, \SeriesRef{maps:src:F.2.19}, \eqref{maps:eq:F.U6}, \SeriesRef{maps:src:F.2.5}, \SeriesRef{maps:src:F.2.18}, \eqref{maps:eq:F.T4}, \SeriesRef{maps:src:F.2.6}, \eqref{maps:eq:F.T6}, \eqref{maps:eq:F.T7}, \eqref{maps:eq:F.T5}, \eqref{maps:eq:F.T8}, \eqref{maps:eq:F.T3}, \eqref{maps:eq:F.T2}, \eqref{maps:eq:F.T9}}


Use the fixed Kuhn cusp and the affine translation clutching
convention of \ref{maps:src:F.2.12}.  In the entire rational family \ref{maps:src:F.2.16},
allow every stable lattice, every free affine finite lift,
and every pair of finite multiplicities.
The data with $\pi_1=1$ and $H_2=H_3=0$ are precisely the
following, up to the integral linear classification already
proved and the translation identifications described below:
\[
\begin{gathered}
 v\in\{0,1\},\qquad
 \text{the unique eligible integral class of \ref{maps:src:F.2.23}},\\
 (m_1,m_2)=(5,4),\qquad
 20\kappa+4r_1+5r_2=\pm1.
\end{gathered}
 \tag{F.T1}\hypertarget{tag-02I7}{}\label{maps:eq:F.T1}
\]
Here $r_i=\ell(v_i)$ is the invariant affine numerator, and
$\kappa\in\mathbb Z$ is the invariant coordinate of the
cusp translation clutching, with signs fixed by the displayed
van Kampen relation below.  The remaining fixed-direction
components of the twists impose no condition.
Every manifold in \eqref{maps:eq:F.T1} is oriented diffeomorphic to its corresponding known example and hence to
$3\mathbb{HP}^2\#(-\mathbb{HP}^2)$.
This completes the classification for this rational family.

\end{theorem}


\begin{proof}
\begin{enumerate}[label={$\langle1\rangle$\arabic*.}]
\item \ref{maps:src:F.2.19} excludes all other rational parameters, and \ref{maps:src:F.2.23}
excludes all other integral linear classes.
We therefore work with either known integral pair, in the
marking of \eqref{maps:eq:F.U6}.  Its primitive invariant homology
covector is $\ell(x)=x_0$.
Each finite fixed lattice has an integral basis
$u_i,\lambda_i$, where $\ell(u_i)=1$,
$\ell(\lambda_i)=0$, and $\lambda_i$ is primitive in
$\ker\ell$.
We may write the affine numerator as
$v_i=r_i u_i+s_i\lambda_i$.
This covers arbitrary finite affine lifts: translating the
origin makes their translation vector invariant by averaging,
and its $m_i$th iterate is then the integral vector $v_i$.
The resulting boundary translation change is included in the
specified affine clutching.

\item For both linear pairs, the normal closure of the cusp angular
lattice in the fibre lattice is exactly $\ker\ell$, by
\ref{maps:src:F.2.5} and \ref{maps:src:F.2.18}.  Van Kampen therefore leaves the abelian
presentation
\[
 m_1a=r_1h,\qquad m_2b=r_2h,\qquad
 a+b+\kappa h=0.
 \tag{F.T2}\hypertarget{tag-02IC}{}\label{maps:eq:F.T2}
\]
This also defines the sign of $\kappa$.
Its determinant is, up to sign,
\[
 p=m_1m_2\kappa+m_2r_1+m_1r_2.
 \tag{F.T3}\hypertarget{tag-02ID}{}\label{maps:eq:F.T3}
\]
Thus simple connectivity is exactly $p=\pm1$.
It forces $\gcd(m_1,m_2)=1$ and
$\gcd(r_i,m_i)=1$ by reduction modulo each multiplicity.

\item We first prove that a three-connected datum has the minimal
multiplicities.  Set $n_1=5$, $n_2=4$ and $m_i=n_i k_i$.
In the splitting $\Lambda=\mathbb Zu_i\oplus\ker\ell$,
write the finite linear map as $1\oplus D_i$.
The reduced finite fibre is a mapping torus of $T^5$ over
the circle $m_i\ell$; its fibre is connected because
$r_i$ is prime to $m_i$.
Its vertical monodromy is a power of $D_i$ prime to $n_i$,
possibly followed by a translation, which changes no
integral Wang group.

\item The primitive fixed vertical loop $\lambda_i$ gives the
only additional free two-cycle beyond vertical cycles.
Consequently an invariant integral two-form $\omega$ on
the covering torus descends integrally exactly when
\[
 \frac{\omega(v_i,\lambda_i)}{m_i}
       =\frac{r_i\omega(u_i,\lambda_i)}{m_i}
       \quad\text{is an integer}.
 \tag{F.T4}\hypertarget{tag-02IE}{}\label{maps:eq:F.T4}
\]
For precision, the torus made from the affine generator
and $\lambda_i$ projects with degree $r_i$ to the mapping-
torus base, whereas the ordinary $u_i$ suspension projects
with degree $m_i$.  A Bezout combination is a primitive
suspension.  Its integrality condition is equivalent to
\eqref{maps:eq:F.T4}, because the ordinary covering periods are already
integral and $\gcd(r_i,m_i)=1$.
The Wang sequence shows that these vertical cycles and one
primitive suspension generate the free second homology;
there is no omitted period condition.

\item The integral invariant two-form computations in \ref{maps:src:F.2.6} and
\ref{maps:src:F.2.18} give
\[
 \gcd_{\omega}\omega(u_i,\lambda_i)=n_i.
 \tag{F.T5}\hypertarget{tag-02IF}{}\label{maps:eq:F.T5}
\]
They are also independently printed for both pairs by
\texttt{fourfolds\_family\_twists.py}.
Hence the covering-image cokernel in degree two at the
$i$th finite point is exactly $\mathbb Z/k_i$.
The global invariant group of $\bigwedge^2\mathcal V$
and its integral first parabolic group both vanish.
The skyscraper cokernel sequence therefore gives
\[
 H^1(B,R^2g_*\mathbb Z)
       =\mathbb Z/k_1\oplus\mathbb Z/k_2.
 \tag{F.T6}\hypertarget{tag-02IG}{}\label{maps:eq:F.T6}
\]
Point-supported torsion kernels do not change this degree-one
base cohomology, and the cusp specialization is onto.
The term \eqref{maps:eq:F.T6}, in bidegree $(1,2)$, survives every Leray
differential on the two-dimensional base.
All rational terms of total degree three vanish, independently
of the multiplicities and twists, by the same exterior
local-system table as \ref{maps:src:F.2.6}.  Thus $H^3$ is finite.
If either $k_i>1$, its surviving nonzero filtration quotient
makes $H^3$ nonzero; universal coefficients then give
nonzero torsion in $H_2$.
This proves $m_1=5$, $m_2=4$ for every three-connected datum.

\item We now identify all twists with these multiplicities.
The following integral fact permits removal of each $s_i$:
\[
 \operatorname{Norm}_{n_i}:
 \ker\ell\longrightarrow\mathbb Z\lambda_i
 \quad\text{is onto}.
 \tag{F.T7}\hypertarget{tag-02IH}{}\label{maps:eq:F.T7}
\]
Here the norm is $I+D_i+\cdots+D_i^{n_i-1}$.
One can read \eqref{maps:eq:F.T7} directly from \eqref{maps:eq:F.T5}.
Let $\varphi_i$ generate the primitive invariant covector
lattice on $\ker\ell$.
The splitting by $u_i$ shows that the gcd in \eqref{maps:eq:F.T5} is
$|\varphi_i(\lambda_i)|$: its forms are
$\ell\wedge\varphi_i$ plus invariant vertical forms.
Thus that value is $n_i$.
The norm has the form
$x\mapsto d_i\varphi_i(x)\lambda_i$ for an integer
$d_i$; evaluating $\varphi_i$ gives
$|d_i|n_i=n_i$, hence $|d_i|=1$.
This proves surjectivity.  The saved program additionally
gives integral preimages $q_i\in\ker\ell$ with
$\operatorname{Norm}_{n_i}q_i=\lambda_i$.

\item Put
\[
 w_i=-\frac{s_i}{n_i}\sum_{j=0}^{n_i-1}jM_i^j q_i.
\]
Then
\[
 (I-M_i)w_i=-s_i\lambda_i/n_i+s_iq_i.
 \tag{F.T8}\hypertarget{tag-02II}{}\label{maps:eq:F.T8}
\]
Translation by $w_i$ conjugates the finite affine action
with numerator $r_i u_i+s_i\lambda_i$ to that with
numerator $r_i u_i$: the remaining $s_iq_i$ is integral.
It gives a diffeomorphism of the entire finite cap.
In the standard linear boundary marking, its translation
section is
$w_i+\theta s_i\lambda_i/n_i$.
Its endpoint mismatch with the linear deck map is exactly
$s_iq_i$ by \eqref{maps:eq:F.T8}, so its integral boundary residue has
$\ell$-coordinate zero.

\item We explain why these local changes extend simultaneously,
including the prescribed cusp marking.
Let $P$ be the complementary pair of pants and
$\mathcal T\to P$ its flat torus group bundle.
A section on the boundary extends over the relative
one-skeleton because the torus is connected.
The only remaining obstruction is the two-cell obstruction
with values in the fibre fundamental-group local system:
\[
 H^2(P,\partial P;\Lambda)
        \simeq H_0(P;\Lambda)=\Lambda_\Gamma.
 \tag{F.T9}\hypertarget{tag-02IJ}{}\label{maps:eq:F.T9}
\]
The equality follows from the oriented relative cellular
complex, equivalently integral Poincar\'e--Lefschetz duality
with local coefficients.  Its boundary residues map to
the indicated global coinvariants.
For both linear pairs the horizontal monodromy relation
matrix has five unit Smith factors, so
$\Lambda_\Gamma=\mathbb Z$, detected exactly by $\ell$.
This is independently checked in the saved program.
Thus the residues $s_iq_i$, with the identity section at the
cusp, have zero obstruction and extend over $P$.
A section may be smoothed relative to its boundary collars;
translation by it is a fibre-preserving diffeomorphism.
Together with the finite-cap maps it removes all $s_i$
without changing the cusp.
The same argument removes all components of translation
clutching whose $\ell$-coordinate is zero.

\item It remains to handle the integer numerators and the remaining
cusp coordinate.  If $p=-1$, the congruences from \eqref{maps:eq:F.T3}
give
\[
 r_1=1+5a,\qquad r_2=-1+4b,
 \qquad a+b+\kappa=0.
 \tag{F.T10}\hypertarget{tag-02I8}{}\label{maps:eq:F.T10}
\]
Adding $5a u_1$ or $4b u_2$ changes no finite affine action;
in its linear boundary marking it changes the translation
residue by $a u_1$ or $b u_2$.
The difference of the cusp marking has residue of
$\ell$-coordinate $\kappa$, by \eqref{maps:eq:F.T2}.
Their total obstruction in \eqref{maps:eq:F.T9} is therefore
$a+b+\kappa=0$.
The same extension argument identifies the glued manifold
with the original twists $(u_1,-u_2)$ and zero cusp
clutching.
If $p=1$, the identical calculation starts from
$(-u_1,u_2)$.

\item Finally simultaneous fibre inversion identifies these two
sign choices by an orientation-preserving diffeomorphism.
It conjugates every affine translation to its negative.
It extends over the Kuhn cusp: the lattice involution
$(x,h)\mapsto(-x,h)$ preserves its toric fan, since a
simplex given by a chain in a unit cube is taken to the
chain with reversed coordinate order in the opposite cube.
It normalizes the cusp lattice action and therefore descends
to its quotient, while fixing the base parameter.  The resulting toric map is holomorphic,
and in the regular real six-torus inversion has determinant
$(-1)^6=1$.  All translation maps also preserve orientation.
Thus every datum in \eqref{maps:eq:F.T1} gives the stated oriented smooth
type and has $H_2=H_3=0$, proving sufficiency as well as
necessity.

\end{enumerate}
\end{proof}




\begin{theorem}[the universal degree-two suspension obstruction]\hypertarget{tag-02NF}{}\label{maps:src:F.2.25}
\StatementRecord{maps:src:F.2.25}{proved}{02NF}{\SeriesRef{maps:src:F.2.14}, \SeriesRef{maps:src:F.2.7}, \SeriesRef{maps:src:F.2.9}, \SeriesRef{maps:src:F.2.17}, \SeriesRef{maps:src:F.2.12}, \eqref{maps:eq:F.T4}, \eqref{maps:eq:F.S3}, \eqref{maps:eq:F.S6}, \SeriesRef{maps:src:E.1.3}, \eqref{maps:eq:F.S2}, \eqref{maps:eq:F.S7}, \eqref{maps:eq:F.S5}, \eqref{maps:eq:F.S8}}


Every admissible datum whose gluing has $\pi_1=1$ and $H_2=H_3=0$ has
exactly two genuine finite fillings and has minimal affine
multiplicities.  After ordering the order-five point first,
its finite characteristic polynomials are one of precisely
the following two remaining necessary possibilities:
\[
\begin{array}{c|c|c}
 (m_1,m_2)&\chi_{M_1}&\chi_{M_2}\\\hline
 (5,4)&(x-1)^2\Phi_5&(x-1)^2\Phi_2^2\Phi_4\\
 (5,6)&(x-1)^2\Phi_5&(x-1)^2\Phi_2^2\Phi_3.
\end{array}
 \tag{F.S1}\hypertarget{tag-02HX}{}\label{maps:eq:F.S1}
\]
In particular all order-eight and order-twelve spectral
loci, and the other two order-six spectral loci of \ref{maps:src:F.2.14},
are excluded for every rational relative position, every
stable lattice and every free affine twist.
There is no multiplicity search bound in this statement:
the only possible affine multiplicities are the
two displayed pairs.

\end{theorem}


\begin{proof}
\begin{enumerate}[label={$\langle1\rangle$\arabic*.}]
\item Theorem \ref{maps:src:F.2.7} excludes three genuine finite fillings.
For two fillings, \ref{maps:src:F.2.9} and \ref{maps:src:F.2.14} give a primitive global
invariant homology covector $\ell$, finite fixed rank two,
and an irreducible rank-five cohomology quotient.
In particular
\[
 H^0(B,j_*\bigwedge^2\mathcal V)=0.
 \tag{F.S2}\hypertarget{tag-02HY}{}\label{maps:eq:F.S2}
\]
This is also immediate from the exterior filtration and
irreducibility in odd dimension, as in \ref{maps:src:F.2.17}.
Simple connectivity gives the determinant quotient of
\ref{maps:src:F.2.12}, so each invariant affine numerator $r_i=\ell(v_i)$
is prime to the multiplicity $m_i$, and the two
multiplicities are coprime.
As before one may conjugate any finite affine lift by a
translation so that its numerator $v_i$ is integral and
fixed by its linear part.

\item Consider one finite point.  Let $M$ have linear order $n$,
so $n\mid m$, and let its complementary degree-four factor
be $f(x)$, with
$\chi_M=(x-1)^2f(x)$.
The image of $\ell$ on the finite fixed lattice is $g\mathbb Z$
for some $g\mid n$: averaging $I+M+\cdots+M^{n-1}$ on
any vector of $\ell$-value one proves this divisibility.
Since $g\mid r$ and $\gcd(r,m)=1$, we have $g=1$.
Choose a fixed integral vector $u$ with $\ell(u)=1$.
It gives an integral invariant splitting
\[
 \Lambda=\mathbb Zu\oplus H,
 \qquad H=\ker\ell,\qquad M=1\oplus D,
 \qquad\chi_D=(x-1)f(x).
 \tag{F.S3}\hypertarget{tag-02HZ}{}\label{maps:eq:F.S3}
\]
The order of $D$ is $n$.
Let $\lambda$ generate its primitive fixed lattice and
let $\varphi$ generate its primitive invariant covector
lattice.  Choose their signs so that
\[
 h=\varphi(\lambda)>0.
 \tag{F.S4}\hypertarget{tag-02I0}{}\label{maps:eq:F.S4}
\]
Their rational pairing is nonzero, since the finite map
is semisimple and its averaging projection has rank one.

\item Both $\operatorname{Norm}_n(D)$ and $f(D)$ are integral
operators, have image on the line $\mathbb Q\lambda$,
and have invariant defining covector proportional to
$\varphi$.  Primitivity of $\lambda$ and $\varphi$
therefore implies
\[
 \operatorname{Norm}_n(D)=\frac n h\lambda\varphi,
 \qquad
 f(D)=\frac{f(1)}h\lambda\varphi,
 \qquad h\mid n,\quad h\mid f(1).
 \tag{F.S5}\hypertarget{tag-02I1}{}\label{maps:eq:F.S5}
\]
For example, the integer coefficient in the second formula
is seen by applying the integral operator to a vector on
which $\varphi$ is one; evaluating on $\lambda$ determines
that coefficient.  No lattice splitting over a localization
or semisimplification modulo a prime is used here.

\item Write $v=ru+s\lambda$.  The quotient finite fibre is a
connected $T^5$ mapping torus over the circle $m\ell$,
because $r$ is prime to $m$.
Its vertical cycles already have integral periods for an
invariant integral two-form on the covering torus.
The Wang sequence leaves one additional primitive free
suspension cycle, from $\lambda$.
Exactly as in \eqref{maps:eq:F.T4}, descent of an invariant form $\omega$
is equivalent to
$m\mid r\omega(u,\lambda)$, and hence to
$m\mid\omega(u,\lambda)$.
In the integral splitting \eqref{maps:eq:F.S3}, every invariant two-form
is an integral multiple of $\ell\wedge\varphi$ plus an
invariant vertical two-form.  Its possible values on
$(u,\lambda)$ are therefore exactly $h\mathbb Z$.
It follows that the degree-two specialization image at
this point has cokernel
\[
 C^2_p\simeq\mathbb Z/(m/h).
 \tag{F.S6}\hypertarget{tag-02I2}{}\label{maps:eq:F.S6}
\]
The integer $m/h$ is defined because $h\mid n\mid m$.
Torsion in the finite fibre cohomology is a specialization
kernel, not an additional restriction on these free periods.
The primitive suspension and vertical cycles generate the
free second homology, so \eqref{maps:eq:F.S6} is the full cokernel.

\item Let $I^2$ be the image sheaf of
$R^2g_*\mathbb Z\to j_*\bigwedge^2\mathcal V$.
The cusp specialization is onto by \SeriesRef{maps:src:E.1.3}, and all kernels
and cokernels at finite points are skyscraper sheaves.
Thus $H^1(R^2g_*\mathbb Z)=H^1(I^2)$.
Equation \eqref{maps:eq:F.S2} and the cokernel exact sequence give an
injection
\[
 \bigoplus_{p\text{ finite}}C^2_p
           \lhook\joinrel\longrightarrow H^1(B,R^2g_*\mathbb Z).
 \tag{F.S7}\hypertarget{tag-02I3}{}\label{maps:eq:F.S7}
\]
The Leray term of bidegree $(1,2)$ has no incoming or
outgoing differential on a real two-dimensional base.
If $H_2=H_3=0$, the universal coefficient theorem gives
$H^3=0$, so that term must vanish.
Every cyclic cokernel in \eqref{maps:eq:F.S7} is therefore zero.
Combining \eqref{maps:eq:F.S5}--\eqref{maps:eq:F.S6} yields the universal necessary
condition
\[
 m=h=n,\qquad n\mid f(1).
 \tag{F.S8}\hypertarget{tag-02I4}{}\label{maps:eq:F.S8}
\]
This both fixes the multiplicity and tests the spectrum.

\item For completeness, apply \eqref{maps:eq:F.S8} to the entire eight-case
list of \ref{maps:src:F.2.14}.  The order-five factor gives $\Phi_5(1)=5$.
At the other point the exact values are
\[
\begin{array}{c|c|r|r}
 n&f&f(1)&n/\gcd(n,f(1))\\\hline
 4&\Phi_2^2\Phi_4&8&1\\
 6&\Phi_2^2\Phi_3&12&1\\
 6&\Phi_2^2\Phi_6&4&3\\
 6&\Phi_3\Phi_6&3&2\\
 8&\Phi_8&2&4\\
 12&\Phi_{12}&1&12\\
 12&\Phi_3\Phi_4&6&2\\
 12&\Phi_4\Phi_6&2&6.
\end{array}
 \tag{F.S9}\hypertarget{tag-02I5}{}\label{maps:eq:F.S9}
\]
The last column is a lower bound for the nonzero local
covering-cokernel order already at the smallest allowed
multiplicity.  Only the first two rows pass.
The complete table is checked by
\texttt{fourfolds\_suspension\_filter.py}; its input is
exactly these eight spectra, with no matrix-entry or
relative-position cutoff.  This proves \eqref{maps:eq:F.S1} and all
claimed universal exclusions.

\end{enumerate}
\end{proof}




\begin{theorem}[the complete rational remaining $(5,6)$ locus]\hypertarget{tag-02NG}{}\label{maps:src:F.2.26}
\StatementRecord{maps:src:F.2.26}{proved}{02NG}{\eqref{maps:eq:F.S1}, \eqref{maps:eq:F.R2}, \SeriesRef{maps:src:F.2.25}, \eqref{maps:eq:F.R3}, \eqref{maps:eq:F.R5}, \SeriesRef{maps:src:F.2.20}, \eqref{maps:eq:F.R6}, \eqref{maps:eq:F.R7}, \eqref{maps:eq:F.R8}, \eqref{maps:eq:F.R12}, \SeriesRef{maps:src:F.2.22}}


The irreducible rational rank-five pairs in the second row
of \eqref{maps:eq:F.S1}, with square-zero rank-two cusp, are classified
up to simultaneous rational conjugacy by
\[
 \{[r:s:t]\in\mathbb P^2(\mathbb Q):
             D(r,s,t)\ne0,\ Q_-(r,s,t)\ne0\}
       \ \sqcup\ \mathbb A^1(\mathbb Q),
 \qquad Q_-=r^2-rs+st,
 \tag{F.R10}\hypertarget{tag-02HK}{}\label{maps:eq:F.R10}
\]
where $D$ is the same cubic as in \eqref{maps:eq:F.R2}.
The exceptional line again lies over $[4:1:3]$.
Thus the two rational spectral loci left by \ref{maps:src:F.2.25} both
have complete explicit normal forms, including their
entire rational exceptional loci.

\end{theorem}


\begin{proof}
\begin{enumerate}[label={$\langle1\rangle$\arabic*.}]
\item Use precisely the space $\mathbb Q\oplus\mathbb Q(\zeta_5)$,
matrices $A,R,X$ and projective plane normalization in
\eqref{maps:eq:F.R3}--\eqref{maps:eq:F.R5}.  The centralizer and image-plane arguments
in \ref{maps:src:F.2.20} require only irreducibility and the two-dimensional
$-1$ eigenspace of $C$; both hold here.
Consequently every pair has a unique projective image-plane
parameter, with only scalar stabilizer.
The unknown $Y$ now satisfies
\[
 YX=0,\qquad YRX=I_2,\qquad
 \operatorname{tr}(YAX)=-2,\qquad
 \operatorname{tr}(YA^{-1}X)=2.
 \tag{F.R11}\hypertarget{tag-02HL}{}\label{maps:eq:F.R11}
\]
Its coefficient matrix is unchanged from \eqref{maps:eq:F.R6}, so has
determinant $D/4$.
The two traces and the $-1$ eigenspace imply
\[
 \chi_C=(x-1)(x+1)^2(x^2+x+1).
\]
Its other three roots are distinct, and the $-1$ eigenspace
has its full multiplicity, so $C^6=I$.
The cusp $XY$ still has square zero and rank two.
For $D\ne0$, the ten linear equations give the unique
explicit solution.

\item The two first-column adjugate identities become
\[
 D y_{10}=2Q_-Q_2,
 \qquad D y_{20}=2(r-s-t)Q_-,
 \tag{F.R12}\hypertarget{tag-02HM}{}\label{maps:eq:F.R12}
\]
with the same $Q_2$ as in \eqref{maps:eq:F.R7}.
For $D\ne0$, these entries vanish simultaneously exactly
when $Q_-=0$.
Indeed, if $Q_-\ne0$ they would instead force
$r=s+t$ and $Q_2=0$, and \eqref{maps:eq:F.R8} would give $D=0$.
The invariant-covector argument of \ref{maps:src:F.2.20} is unchanged.
Thus the pair is rationally irreducible exactly off the
conic $Q_-=0$ in the generic locus.

\item No nonzero rational triple satisfies $D=Q_-=0$.
If $s=0$, the conic forces $r=0$, and then $D=7t^3$.
If $s\ne0$, rescale to $s=1$ and put
$r=u$, $t=u-u^2$.  Direct substitution gives
\[
 D(u,1,u-u^2)=
 -(7u^2-11u+7)(u^4-u^3+u^2-u+1).
 \tag{F.R13}\hypertarget{tag-02HN}{}\label{maps:eq:F.R13}
\]
The first factor has discriminant $-75$ and positive
leading coefficient.  The second is strictly positive,
as the identity
\[
 u^4-u^3+u^2-u+1
 =\left(u^2-\frac u2\right)^2
  +\frac34\left(u-\frac23\right)^2+\frac23
\]
shows.  In particular neither has a rational zero.
On the exceptional cubic the necessary identities \eqref{maps:eq:F.R12}
therefore force the same two points $[3:2:1]$ and $[4:1:3]$
as before.  The first is again impossible because
$[X\mid RX]$ has rank less than four.
At the second, the complete solution is
\[
 Y_z=\begin{pmatrix}
 4z+26/5&0&-2z-2&-2z-1/5&2z+1\\
 2z-2/5&0&-z&2/5-z&z
 \end{pmatrix},\qquad z\in\mathbb Q.
 \tag{F.R14}\hypertarget{tag-02HO}{}\label{maps:eq:F.R14}
\]
Its first-column entries cannot vanish simultaneously,
so each pair on this line is irreducible.
There are no further identifications, by the scalar
stabilizer argument already proved in \ref{maps:src:F.2.20}.
This establishes the claimed bijection.

\item The exact program is
\texttt{fourfolds\_relative\_six\_probe.py}.
It checks the modified characteristic polynomials, finite-order
identities, first-column adjugate expressions, conic
substitution and entire exceptional line symbolically.
Its auxiliary trace-parameter computations for the other
repeated-negative spectrum are recorded as well, although
that spectrum is now excluded integrally by \ref{maps:src:F.2.25}.
No finite sample is used to infer a rational classification.
The finite lattice and extension algorithm of \ref{maps:src:F.2.22} applies
to every point of \eqref{maps:eq:F.R10}; it does not by itself decide
which rational parameters survive the remaining integral
homology tests.

\end{enumerate}
\end{proof}




\begin{theorem}[the entire remaining order-six locus is excluded]\hypertarget{tag-02NH}{}\label{maps:src:F.2.27}
\StatementRecord{maps:src:F.2.27}{proved}{02NH}{\eqref{maps:eq:F.S1}, \SeriesRef{maps:src:F.2.25}, \eqref{maps:eq:F.S3}, \eqref{maps:eq:F.S5}, \eqref{maps:eq:F.X2}, \eqref{maps:eq:F.X3}, \eqref{maps:eq:F.S7}, \eqref{maps:eq:F.X5}, \eqref{maps:eq:E.9}}


Every simply connected datum in the second row of \eqref{maps:eq:F.S1},
at every rational relative position and every stable integral
lattice with primitive maximal cusp, has
\[
                 \operatorname{Tor}_3 H_3(M;\mathbb Z)\ne0
\]
if its degree-two suspension cokernels vanish.
Consequently no such datum is three-connected.  Together with
\ref{maps:src:F.2.25} this leaves only the minimal order pair $(5,4)$ as a possible three-connected admissible datum.
\end{theorem}


\begin{proof}
\begin{enumerate}[label={$\langle1\rangle$\arabic*.}]
\item We use cohomology with coefficients in the discrete valuation
ring $R=\mathbb Z_3$.  Suppose the degree-two suspension
cokernels vanish; otherwise \ref{maps:src:F.2.25} already gives nonzero
$H_2$ or $H_3$.  That theorem forces $m=n=h=6$ at the
order-six point, in the notation \eqref{maps:eq:F.S3}--\eqref{maps:eq:F.S5}.
We first determine its local integral representation over $R$.

\item The splitting by the fixed vector $u$ is integral:
$\Lambda_R=Ru\oplus H_R$.  On $H_R$, the finite map $D$
has characteristic polynomial
$(x-1)(x+1)^2\Phi_3(x)$.  Since two is a unit in $R$,
the idempotents $(I\pm D^3)/2$ give an integral direct sum
\[
 H_R=L_+\oplus L_-,\qquad
 \operatorname{rank}L_+=3,\quad
 \operatorname{rank}L_-=2.
\]
On $L_-$ the map is exactly $-I$, by its rational spectrum.
On $L_+$ it has order three.  The primitive fixed vector
$\lambda$ lies in $L_+$, and the primitive invariant
covector $\varphi$ vanishes on $L_-$, so their restrictions
remain primitive.  Equation \eqref{maps:eq:F.S5} with $h=6$ gives
\[
 2(I+D+D^2)|_{L_+}=\lambda\varphi|_{L_+}.
 \tag{F.X1}\hypertarget{tag-02IV}{}\label{maps:eq:F.X1}
\]
This operator is nonzero modulo three.  But modulo three
$I+D+D^2=(D-I)^2$ and $(D-I)^3=0$.  The reduction of $D$
on the rank-three space therefore has one Jordan block
of length three.  Choose $q\in L_+$ whose image has
$(D-I)^2q\ne0$ modulo three.  Then $q,Dq,D^2q$ are a
basis modulo three and hence an $R$-basis.  In this basis
$D$ is the three-cycle permutation matrix, since $D^3=I$.
We have proved, without a global lattice classification,
that the full finite representation over $R$ is
\[
                   1\oplus\operatorname{Perm}_3\oplus(-I_2).
 \tag{F.X2}\hypertarget{tag-02IW}{}\label{maps:eq:F.X2}
\]
Its dual has the same form, by reversing the three-cycle.

\item Let $u,a,b,c,d,e$ be a cohomology basis for \eqref{maps:eq:F.X2}, with
$a,b,c$ cycled and $d,e$ negated.  A full integral basis
of the invariant degree-three lattice is
\[
 u(ab+bc+ca),\qquad abc,\qquad ude,\qquad(a+b+c)de.
 \tag{F.X3}\hypertarget{tag-02IX}{}\label{maps:eq:F.X3}
\]
Here juxtaposition denotes exterior multiplication.
To check completeness, terms with an odd number of
$d,e$ are negated by the cube of the finite map and
have no invariants.  The terms with zero or two such
factors give exactly the four orbit sums in \eqref{maps:eq:F.X3},
each primitive.  With volume $uabcde$, their skew
pairing has one block with coefficient three and one
block with coefficient a unit.  Thus its determinant
has three-adic valuation two.  The program \texttt{fourfolds\_universal\_56\_obstruction.py} checks
the full invariant kernel, its primitivity, and the
pairing independently by exterior multiplication and
six-column determinants; its Smith factors are
$(1,1,3,3)$.

\item Let $B_6$ be the reduced finite fibre.  Its torus
cover has degree six.  Its free middle cohomology has
rank four and unimodular skew pairing, by integral
Poincar\'e duality.  If $I$ denotes the index of its
pullback lattice in the invariant lattice, multiplication
of pairings by the covering degree gives
\[
 2v_3(I)+2=4v_3(6)=4,
 \qquad v_3(I)=1.
 \tag{F.X4}\hypertarget{tag-02IY}{}\label{maps:eq:F.X4}
\]
Consequently the three-primary degree-three covering
cokernel is $R/3R$.  Translations of the affine action
do not enter this calculation: the cover degree,
linear invariant lattice and oriented pairings are
those of the free quotient for every allowed
affine numerator.

\item We explain why this local group must appear in total
homology even when the global integral lattice varies.
Let $\mathcal V$ be the rational rank-six cohomology
local system.  It is an extension of the irreducible
rank-five system $W$ by a trivial line.  The determinant
of $W$ is trivial.  Irreducibility in odd dimension
implies that neither $\bigwedge^2 W$ nor its dual
has a nonzero invariant: an invariant alternating
form would have an invariant radical and would have
to be nondegenerate, impossible in odd dimension.
Also $\bigwedge^3W\simeq\bigwedge^2W^*$ by its
trivial determinant.  The exterior filtration therefore
gives
\[
 (\bigwedge^3\mathcal V)^\Gamma=0,\qquad
 ((\bigwedge^2\mathcal V)^*)^\Gamma=0.
 \tag{F.X5}\hypertarget{tag-02IZ}{}\label{maps:eq:F.X5}
\]
These rational vanishing statements remain true over
$\mathbb Q_3$, since they are kernels of rational
matrices.

\item As in \eqref{maps:eq:F.S7}, take the image of
$R^3g_*R\to j_*\bigwedge^3\mathcal V_R$.
Its kernel and cokernel are point-supported sheaves.
The first vanishing in \eqref{maps:eq:F.X5} makes the local group
$R/3R$ inject into $H^1(B,R^3g_*R)$.  Bidegree $(1,3)$
has no incoming or outgoing Leray differential on
the two-dimensional base, so this injection survives
to $E_\infty^{1,3}$.

\item There is one extension issue, and it cannot remove
this torsion.  The second vanishing in \eqref{maps:eq:F.X5} says
that $H^2(B,R^2g_*R)$ has rational rank zero.
Point-supported changes do not alter this degree-two
base cohomology; it is the coinvariant module, as
in \eqref{maps:eq:E.9}.  It is a finitely generated torsion $R$-module
and hence a finite group.  Therefore
$E_\infty^{2,2}$ is finite as well.  In the filtration
of $H^4(M;R)$, take the inverse image of the injected
$R/3R\subset E_\infty^{1,3}$.  It is an extension
of $R/3R$ by the finite group $E_\infty^{2,2}$.
It is thus a nonzero finite three-primary subgroup
of $H^4(M;R)$; in particular it cannot be absorbed
as a finite quotient of a free filtration kernel.
Universal coefficients identify this torsion with
the three-primary torsion of $H_3(M;\mathbb Z)$.
This proves the claim at every rational position and
every integral lattice, and completes the exclusion.

\end{enumerate}
\end{proof}




\begin{theorem}[universal fundamental group and twist normalization]\hypertarget{tag-02NI}{}\label{maps:src:F.2.28}
\StatementRecord{maps:src:F.2.28}{proved}{02NI}{\SeriesRef{maps:src:F.2.25}, \eqref{maps:eq:F.S5}, \eqref{maps:eq:F.T8}, \SeriesRef{maps:src:F.2.24}}


In the remaining $(5,4)$ locus, fix the marked standard Kuhn
cap and allow affine translation clutching.
Let $\Lambda=\mathbb Z^6$ be fibre homology and let
$\Lambda_\Gamma$ be its global monodromy coinvariants.
Write $\Lambda_\Gamma=\mathbb Zh\oplus T$, with $T$ finite,
and let $\ell$ be the primitive invariant covector detecting
the free factor.  The following are complete reductions for
any datum that could satisfy $\pi_1=1$ and $H_2=H_3=0$.

The multiplicities must be $(5,4)$, and the local
vertical fixed pairings must be $h_1=5,h_2=4$, as in \ref{maps:src:F.2.25}.
For any affine twists and clutching the fundamental group
is abelian.  If its free presentation determinant is nonzero,
its order is
\[
 |\pi_1|=|20\kappa+4r_1+5r_2|\,|T|,
 \qquad r_i=\ell(v_i).
 \tag{F.T11}\hypertarget{tag-02I9}{}\label{maps:eq:F.T11}
\]
If that determinant is zero the group is infinite.
In particular simple connectivity is equivalent to
\[
 T=0,\qquad 20\kappa+4r_1+5r_2=\pm1.
 \tag{F.T12}\hypertarget{tag-02IA}{}\label{maps:eq:F.T12}
\]
Whenever these conditions and $h_1=5,h_2=4$ hold, every
such twist choice is oriented diffeomorphic to the normalized
filling with twists $(u_1,-u_2)$ and zero cusp clutching,
where $u_i$ is any integral fixed vector with $\ell(u_i)=1$.
Thus each eligible integral linear datum has exactly one
normalized smooth filling to test for the remaining
homology conditions; no affine-twist enumeration remains.

\end{theorem}


\begin{proof}
\begin{enumerate}[label={$\langle1\rangle$\arabic*.}]
\item The image of the fibre fundamental group is a normal
abelian subgroup in the van Kampen presentation.
After quotienting by it, the two finite meridians satisfy
$a^5=b^4=ab=1$.
That quotient is trivial, because five and four are coprime.
The entire fundamental group is therefore equal to the
image of the fibre group and is abelian.
This argument does not assume that the normal closure of
cusp vanishing cycles already exhausts $\ker\ell$.

\item In abelianization the monodromy relations give
$\Lambda_\Gamma$.
The primitive cusp angular lattice is the image of its
nilpotent operator, and this image is already zero among
those monodromy relations.
The remaining presentation is
\[
 5a=v_1,\qquad4b=v_2,\qquad a+b+v_0=0
 \quad\text{in }\Lambda_\Gamma\oplus\mathbb Za\oplus\mathbb Zb,
 \tag{F.T13}\hypertarget{tag-02IB}{}\label{maps:eq:F.T13}
\]
where $\ell(v_0)=\kappa$ fixes the sign convention.
Use any Smith presentation of $T$ and split the free part
$\mathbb Zh$.
The full relation matrix is block triangular, with the
presentation of $T$ as one diagonal block and the
three-by-three free matrix of determinant
$\pm(20\kappa+4r_1+5r_2)$ as the other.
Its determinant gives \eqref{maps:eq:F.T11}; when the latter free
matrix has rank less than three the group has positive
rank and is infinite.
For determinant $\pm1$, elimination of the three free
variables leaves precisely $T$, even if the twists have
nonzero torsion components.
This proves the exact simple-connectivity criterion.

\item We give the global twist comparison to make its scope
precise.  Under \eqref{maps:eq:F.T12}, the finite fixed covector image
is one at each point, so choose integral fixed $u_i$ with
$\ell(u_i)=1$.  The primitive vertical fixed vector
$\lambda_i$ and invariant covector $\varphi_i$ satisfy
$\varphi_i(\lambda_i)=m_i$, by the assumed maximal local
pairings.  Formula \eqref{maps:eq:F.S5} makes the finite norm onto
$\mathbb Z\lambda_i$.
Consequently every vertical numerator component
$s_i\lambda_i$ is removed by the explicit cap translation
\eqref{maps:eq:F.T8}, whose boundary residue has $\ell$-coordinate zero.

\item For the complementary pair of pants the full relative
obstruction group for extending boundary translation
sections is
$H^2(P,\partial P;\Lambda)=\Lambda_\Gamma$.
Now $T=0$, so it is exactly $\mathbb Z$ detected by
$\ell$, not just a rational or free quotient.
Prescribing the identity translation at the cusp therefore
allows all the vertical finite changes to extend over $P$.
The same argument removes the zero-$\ell$ component of
any cusp translation clutching.
For determinant $-1$, write
$r_1=1+5a$, $r_2=-1+4b$; then \eqref{maps:eq:F.T12} says
$a+b+\kappa=0$.
The two integer changes of finite boundary marking and
the remaining cusp change thus also have zero total
relative obstruction and extend.
This gives the normalized twists $(u_1,-u_2)$.
For determinant $1$ the same argument gives their negatives.
The orientation-preserving fibre inversion extending over
the Kuhn fan, proved in \ref{maps:src:F.2.24}, identifies the two choices.
All maps used in the comparison are cap translations or
translations by smooth torus-bundle sections over $P$;
the prescribed cusp is retained throughout.
This proves the claimed complete twist normalization.

\end{enumerate}
\end{proof}




\begin{theorem}[A finite integral test for three-connectivity]\hypertarget{tag-02NJ}{}\label{maps:src:F.2.29}
\StatementRecord{maps:src:F.2.29}{proved}{02NJ}{\SeriesRef{maps:src:F.2.20}, \SeriesRef{maps:src:F.2.22}, \SeriesRef{maps:src:F.2.25}, \SeriesRef{maps:src:F.2.28}, \SeriesRef{maps:src:F.2.6}, \eqref{maps:eq:F.HC7}, \SeriesRef{maps:src:E.1.3}, \SeriesRef{maps:src:F.2.16}}


Take any rational pair in \ref{maps:src:F.2.20} and any of its integral,
primitive-cusp rank-six extensions given by \ref{maps:src:F.2.22}.
Write $M_1,M_2$ for geometric transport on fibre homology
$\Lambda=\mathbb Z^6$, in the peripheral convention
$M_0M_1M_2=I$.
First require the preliminary conditions
\[
 \Lambda_\Gamma=\mathbb Z,\qquad
 \ell(\Lambda^{M_i})=\mathbb Z,\qquad h_1=5,\quad h_2=4 .
 \tag{F.HC1}\hypertarget{tag-02G9}{}\label{maps:eq:F.HC1}
\]
Here $\ell$ is the primitive global invariant covector, and
$h_i$ is the integral fixed-vector/fixed-covector pairing
on $\ker\ell$ from \ref{maps:src:F.2.25}.
Failure of any condition excludes the datum.
If they hold, use the unique normalized affine filling of
\ref{maps:src:F.2.28}, with twists $(u_1,-u_2)$ and zero cusp clutching.
The following six finite integer-matrix tests are necessary
and sufficient for this filling to have $H_2=H_3=0$.

Let
\[
 T_i=M_i^{-t},\quad T_0=(T_1T_2)^{-1},\qquad
 G_q=\operatorname{coker}
 [\,\bigwedge^qT_1-I\mid\bigwedge^qT_2-I\,]
 \quad(1\le q\le4).
 \tag{F.HC2}\hypertarget{tag-02GE}{}\label{maps:eq:F.HC2}
\]
All $G_q$ are finite.  Let $P_q$ be the integral parabolic
cohomology group of the exterior local system
$\bigwedge^q\mathcal V$, for $q=2,4$.
Let $t_q\in G_{q-1}$, $2\le q\le5$, be the natural global
coinvariant image of the unique nonzero torsion suspension
class at the order-four cap, as specified below.
Let $c_4\in G_3$ be the natural image of the unique nonzero
torsion class in
$\operatorname{coker}(\bigwedge^3T_0-I)$.
Then the six tests are
\[
\begin{array}{c|ccc}
 &\text{first test}&\text{second test}&\text{third test}\\ \hline
 H^3=0&G_1=\langle t_2\rangle&P_2=0&t_3\ne0\\
 H^5=0&G_3=\langle t_4,c_4\rangle&P_4=0&t_5\ne0 .
\end{array}
\tag{F.HC3}\hypertarget{tag-02GF}{}\label{maps:eq:F.HC3}
\]
Every group, class, and equality here is computed by
finitely many integer kernels, Smith forms, and exterior
minors.  This is a decision criterion for every rational parameter and every integral lattice in the normal form.

\end{theorem}


\begin{proof}
\begin{enumerate}[label={$\langle1\rangle$\arabic*.}]
\item We first establish the uniform local calculations needed
for the criterion, so the test does not assume that a
general finite matrix is globally conjugate to the example
in \ref{maps:src:F.2.6}.
The fixed vectors $u_i$ split the homology lattice as
$\mathbb Zu_i\oplus H_i$, where $H_i=\ker\ell$.
Write $D_i=M_i|_{H_i}$.
The maximal pairing $h_i=m_i$ gives
\[
 \operatorname{Norm}_{m_i}(D_i)=\lambda_i\varphi_i ,
 \tag{F.HC4}\hypertarget{tag-02GG}{}\label{maps:eq:F.HC4}
\]
with both factors primitive.

\item At five, work over $\mathbb Z_5$.
Modulo five the norm is $(D_1-I)^4$ and is nonzero.
Choose $q$ with $\operatorname{Norm}_5q$ primitive.
The five vectors $q,D_1q,\ldots,D_1^4q$ are independent
modulo five: otherwise the nilpotent Jordan block generated
by $q$ would have length less than five.
They are therefore a $\mathbb Z_5$ basis.
Thus $D_1$ is the regular five-cycle permutation locally.

\item At four, work over $\mathbb Z_2$.
The norm is $(D_2-I)^3$ modulo two and is nonzero.
The vectors $q,D_2q,D_2^2q,D_2^3q$ similarly generate
a primitive regular four-cycle submodule $Q$ of rank four.
The quotient is rank one with $D_2=-1$, by the
characteristic polynomial.
For a lift $z$ of its generator, put $a=(D_2+I)z\in Q$.
Since $D_2^4=I$,
\[
 (D_2^3-D_2^2+D_2-I)a=0.
\]
On the regular four-cycle module the kernel of the
displayed alternating norm is exactly $(D_2+I)Q$,
integrally.  One can see this by writing its four
coefficients; equivalently $D_2+I$ has three nonzero
Smith factors, all one, and the indicated rational image.
Subtracting a suitable element of $Q$ from $z$ makes
$D_2z=-z$.
Consequently
\[
 M_1\sim 1\oplus\operatorname{perm}_5\quad\text{over }\mathbb Z_5,
 \qquad
 M_2\sim1\oplus\operatorname{perm}_4\oplus(-1)
 \quad\text{over }\mathbb Z_2 .
 \tag{F.HC5}\hypertarget{tag-02GH}{}\label{maps:eq:F.HC5}
\]
These local comparisons apply to cap cohomology, not just
to abstract representation lattices.
Indeed the integral lattice of intertwiners is the kernel
of an integer matrix and its tensor with $\mathbb Z_p$
is the lattice of local intertwiners.
Approximate a local invertible intertwiner by an integral
one whose determinant is a $p$-adic unit.
Keeping the $u_i$ summand fixed produces an equivariant
torus isogeny of degree prime to $p$, preserving the
normalized translation $\pm u_i/m_i$.
The resulting cap and boundary coverings induce
isomorphisms on $\mathbb Z_p$ cohomology: the vertical
torus maps are cohomology isomorphisms over $\mathbb Z_p$,
so the mapping-torus Wang sequences give the assertion.
At primes not dividing $m_i$, transfer already makes
finite-cover pullback an isomorphism onto invariants.

\item Here are the full required integral results for these
two monomial models.
The ranks of their invariant exterior lattices in
degrees zero through six are
$(1,2,3,4,3,2,1)$.
The order-five cap has no cohomology torsion in degrees
two through five.
The order-four cap has exactly $\mathbb Z/2$ in each
of those four degrees.
This follows from the Wang sequence for the mapping
torus of the vertical $D_i$:
its torsion in degree $q$ is
\[
 \operatorname{Tor}\operatorname{coker}
       (\bigwedge^{q-1}D_i^t-I).
 \tag{F.HC6}\hypertarget{tag-02GI}{}\label{maps:eq:F.HC6}
\]
In particular it introduces no additional torsion primes.

\item Integral pullback from either finite cap is onto the
invariant lattice in degrees two through five.
For degree two this is precisely the $m_i/h_i=1$
suspension calculation of \ref{maps:src:F.2.25}.
For completeness, the determinants of the pairings between
the invariant lattices in complementary degrees
$(1,5),(2,4),(3,3)$ are respectively
$m_i,m_i^3,m_i^4$ in absolute value.
The degree-one cover image has index $m_i$.
The free downstairs complementary pairings are unimodular,
and pullback multiplies each pairing by the covering
degree $m_i$.
The middle pairing first gives degree-three index one.
The degree-two index one then gives degree-four index one;
the degree-one index $m_i$ gives degree-five index one.
These determinant identities and all the Wang Smith
forms are verified by the fixed, exhaustive computation
in \texttt{fourfolds\_local\_caps.py}, for both orders and
all exterior degrees zero through six.
Thus no unexamined finite matrix is being inferred from
an example.

\item We next identify the boundary maps.
Choose any unimodular matrix $P_2=[u_2,H_2]$ realizing
the splitting at four, and put
$J=P_2^{-t}|_{\mathbb Z^5}$ for the inclusion of vertical
cohomology.
Choose the unique nonzero torsion representative
\[
 w_q\in\operatorname{coker}
            (\bigwedge^{q-1}D_2^t-I),\qquad
 t_q=[(\bigwedge^{q-1}J)w_q]\in G_{q-1}.
 \tag{F.HC7}\hypertarget{tag-02GJ}{}\label{maps:eq:F.HC7}
\]
A Smith transformation supplies an integral $w_q$.
Its norm is exactly zero, since it lies rationally in
the image of $\bigwedge^{q-1}D_2^t-I$.
The finite boundary coordinates of \ref{maps:src:F.2.6} satisfy
$\theta=a+4k$, $h=k$, hence $a=\theta-4h$.
The pullback of the suspension cocycle has $\theta$
component $w_q$ and $h$ component
$-\operatorname{Norm}_4(w_q)=0$.
This holds in every exterior degree, including $q=5$.
It is an integral cellular assertion for the monomial
model, where the torus cellular differentials vanish;
the prime-to-two isogeny comparison transports it to
every eligible cap.
The base connecting map therefore gives exactly the
class $t_q$ in \eqref{maps:eq:F.HC7}.
Different representative choices change it by a global
monodromy relation.  Any sign change is immaterial for
an order-two class.

\item There is one further boundary map, from the free
degree-four cusp kernel.
The primitive cusp is integrally conjugate to
$\left(\begin{smallmatrix}I&I\\0&I\end{smallmatrix}\right)$.
Its degree-three exterior difference has ten Smith
factors one and one Smith factor two; its cokernel is
$\mathbb Z^9\oplus\mathbb Z/2$.
The other exterior differences in degrees one, two,
four and five have torsion-free cokernels.
The exact matrices and a torsion representative are
saved by \texttt{fourfolds\_cusp\_boundary.py}.
Let $C$ be the smooth cusp cap and $Y=\partial C$.
The boundary Wang sequence gives
$\operatorname{Tor}H^4(Y)=\mathbb Z/2$.
In the Poincar\'e--Lefschetz exact sequence its image in
\[
 H^5(C,Y)\cong H_3(C)=\mathbb Z^9
\]
is zero, so this torsion class lifts to $H^4(C)$.
Its restriction to an ordinary fibre is zero.
By \SeriesRef{maps:src:E.1.3} the kernel of integral specialization is
a primitive copy $K=\mathbb Z$ in degree four.
The lift is in $K$, and the image of this cyclic group
in $H^4(Y)$ contains the nonzero torsion class.
A cyclic image containing nonzero torsion cannot
also have a nonzero free part.
Its primitive generator therefore restricts to the
unique nonzero torsion class of $H^4(Y)$.
Consequently the cusp contribution to $d_2$ is
exactly $c_4$ as defined above.  In particular this
column must not be discarded merely because it
vanishes in the first examples.

\item It remains to prove that these are all the tests.
For the irreducible odd-dimensional quotient $W$
of $\mathcal V$, invariant alternating forms and
bivectors vanish: a nonzero one would give an
isomorphism $W\cong W^*$ with an odd-dimensional
nondegenerate alternating pairing.
The exact exterior sequences for
$0\to\mathbb Q\to\mathcal V_{\mathbb Q}\to W\to0$,
together with $\det W=1$, give
\[
 H^0(B,j_*\bigwedge^q\mathcal V)=0
       \quad(2\le q\le5),\qquad
 G_q\otimes\mathbb Q=0\quad(1\le q\le4).
 \tag{F.HC8}\hypertarget{tag-02GK}{}\label{maps:eq:F.HC8}
\]
For the endpoint assertions, the extension cannot split:
on $\mathbb Q\oplus W$ a square-zero nilpotent has rank at
most two, whereas the stipulated rank is three.
An invariant covector of $\mathcal V_{\mathbb Q}$ would
split this extension, since $W$ has no trivial quotient.
Thus $(\mathcal V_{\mathbb Q}^*)^\Gamma=0$, which gives
both $H^0(\bigwedge^5\mathcal V_{\mathbb Q})=0$ and
$G_1\otimes\mathbb Q=0$.
Here and below the first assertion is integral as well,
since the coefficient lattices are torsion free.
The local invariant ranks in exterior degrees two
and four are $3,3,9$ at the three punctures.
The Euler formula for the rank-fifteen local system
therefore gives $P_2\otimes\mathbb Q=P_4\otimes\mathbb Q=0$.
In degree one it similarly gives rational parabolic
cohomology zero.
The only global degree-one section is the invariant
line.  Simple connectivity and the universal
coefficient theorem make its free transgression
onto $H^2(B,\mathbb Z)$; equivalently the rational
second Betti number is zero, as in \ref{maps:src:F.2.6}.
Thus the normalized filling has rational $b_2=b_3=0$.

\item All specializations in exterior degrees two through
five are onto, by the local calculations and \SeriesRef{maps:src:E.1.3}.
Their only point-supported kernels are
\[
\begin{array}{c|cccc}
q&2&3&4&5\\ \hline
\text{kernel}&\mathbb Z/2&\mathbb Z/2&
                 \mathbb Z\oplus\mathbb Z/2&\mathbb Z/2 .
\end{array}
\tag{F.HC9}\hypertarget{tag-02GL}{}\label{maps:eq:F.HC9}
\]
The free summand in degree four is at the cusp; the
others are at four.
Consequently $H^1(B,R^qg_*\mathbb Z)=P_q$ in the
degrees used here, while
$H^2(B,R^qg_*\mathbb Z)=G_q$.
The latter equality also holds in degree one:
changing a sheaf at isolated points changes no
degree-two cohomology.  The equality with coinvariants
follows from compactly supported top cohomology on
the punctured sphere and the sequence $j_!\to j_*$.

\item The base has real dimension two.
Thus the only possible nonzero differential in the
integral Leray spectral sequence is
$d_2:E_2^{0,q}\to E_2^{2,q-1}$.
The complete associated graded groups in total
degrees three and five are therefore
\[
\begin{split}
 \operatorname{gr}H^3&:
   G_1/\langle t_2\rangle,\quad P_2,\quad
   \ker(\mathbb Z/2\xrightarrow{\,t_3\,}G_2),\\
 \operatorname{gr}H^5&:
   G_3/\langle t_4,c_4\rangle,\quad P_4,\quad
   \ker(\mathbb Z/2\xrightarrow{\,t_5\,}G_4).
\end{split}
\tag{F.HC10}\hypertarget{tag-02GA}{}\label{maps:eq:F.HC10}
\]
These are all finite.  A finite filtered group
vanishes exactly when every associated graded group
vanishes, regardless of extension classes.
Since $\pi_1=1$ and rational $b_2=b_3=0$, the universal
coefficient theorem gives
$H^3=\operatorname{Ext}(H_2,\mathbb Z)$, so
$H_2=0$ exactly when $H^3=0$.
Integral Poincar\'e duality gives $H^5\cong H_3$.
This proves \eqref{maps:eq:F.HC3}.

\item For reproducibility, no sheaf calculation is left
as a black box in the integer test.
Given $A=\bigwedge^qT_1$, $B=\bigwedge^qT_2$ on a
free lattice $L$, form the saturated integer kernel
\[
 K_q=\ker[\,A-I\mid A(B-I)\mid -(AB-I)\,]\subset L^3.
 \tag{F.HC11}\hypertarget{tag-02GB}{}\label{maps:eq:F.HC11}
\]
Its image under
$(x,y,z)\mapsto((A-I)x,(B-I)y)$ is precisely the
integral parabolic cocycle lattice.
Quotient by the columns of
$\left(\begin{smallmatrix}A-I\\B-I\end{smallmatrix}\right)$
to obtain $P_q$.
This uses the integral images of peripheral
differences, not their saturations.
Smith transformations also supply every $w_q$, the
cusp order-two representative, and the tests for
their images and generated subgroups in $G_q$.

\item The program \texttt{fourfolds\_survivor\_test.py}
implements exactly these operations.
Its executed test set is precisely $v=0,1,2$ from
\ref{maps:src:F.2.16}.  For $v=0,1$ all six tests pass;
$G_1,\ldots,G_4$ are each $\mathbb Z/2$,
$P_2=P_4=0$, all four finite-cap images are nonzero,
and the cusp image is zero.
For $v=2$ they are each $\mathbb Z/6$,
$P_2=P_4=\mathbb Z/3$, and the first, second,
fourth and fifth tests fail.  The computation thus
checks both accepted and rejected data.

\item We finish by making the convention dictionary explicit.
The imported family matrices called $A,B$ in that
program are geometric cohomology pullbacks, so the
geometric homology inputs are $(A^t,B^t)$.
Their genuine left cohomology transport in \eqref{maps:eq:F.HC2}
is $(A^{-1},B^{-1})$, which is what the code uses.
For comparison with computations written using
$(A,B)$, simultaneous reversal of the two peripheral
generators sends
\[
 g_1\mapsto g_1^{-1},\qquad
 g_2\mapsto g_2^{-1},\qquad
 g_0=(g_1g_2)^{-1}\mapsto
     g_2g_1=g_1^{-1}(g_1g_2)g_1 .
 \tag{F.HC12}\hypertarget{tag-02GC}{}\label{maps:eq:F.HC12}
\]
This is an orientation-reversing marking of the
pair of pants fixing its three ends: each peripheral
conjugacy class is inverted.
Thus the new cusp is conjugate to the inverse old
cusp, not an independently chosen third matrix.
Inversion preserves invariant and coinvariant
lattices, and the indicated conjugating generator
acts trivially on global coinvariants.
All boundary classes in \eqref{maps:eq:F.HC3} have order two, so
reversal signs do not change the tests.
The implementation uses \eqref{maps:eq:F.HC2} directly and hence
does not depend on this comparison for its validity.

\end{enumerate}
\end{proof}




\begin{corollary}[Classification of the three-connected data]\hypertarget{tag-02NL}{}\label{maps:src:F.2.30}
\StatementRecord{maps:src:F.2.30}{proved}{02NL}{\SeriesRef{maps:src:F.2.20}, \SeriesRef{maps:src:F.2.22}, \eqref{maps:eq:F.HC1}, \eqref{maps:eq:F.HC3}, \eqref{maps:eq:F.T12}, \SeriesRef{maps:src:F.3.11}, \SeriesRef{maps:src:F.2.7}, \SeriesRef{maps:src:F.2.14}, \SeriesRef{maps:src:F.2.25}, \SeriesRef{maps:src:F.2.27}, \SeriesRef{maps:src:F.2.28}, \SeriesRef{maps:src:F.2.29}, \SeriesRef{maps:src:E.1.3}, \SeriesRef{maps:src:F.2.15}, \SeriesRef{maps:src:F.2.5}, \SeriesRef{sec:V-known-inputs}, \SeriesRef{maps:src:F.4.1}}


For the marked Kuhn cap and affine translation clutchings
used throughout this section, the following is a complete constructive
classification of the data with
$\pi_1=1$ and $H_2=H_3=0$, up to integral conjugacy of their
linear tuples and oriented diffeomorphism of their fillings.

There are no data with three genuine finite fillings.
With two fillings their affine orders are five and
four.  Order the order-five point first.
The complete, unbounded parameter set is the rational
projective normal form of \ref{maps:src:F.2.20} together with its
exceptional rational affine line.
For each parameter, run the terminating order,
conductor, stable-lattice, and primitive-extension algorithm
of \ref{maps:src:F.2.22}.  Discard precisely the integral classes failing
\eqref{maps:eq:F.HC1} or one of \eqref{maps:eq:F.HC3}.
The classes that pass these tests are exactly all
integral linear tuples with three-connected gluing.
For each such class all simply connected affine
twists are exactly the choices satisfying \eqref{maps:eq:F.T12};
their fillings are all oriented diffeomorphic to its
normalized filling.
If ordered labels with the order-four point first are
desired, interchange the finite labels and recompute
$M_0=(M_1M_2)^{-1}$.

Every such datum has
\[
 \pi_1=1,\qquad
 H_j=\begin{cases}
  \mathbb Z,&j=0,8,\\
  \mathbb Z^4,&j=4,\\
  0,&\text{otherwise},
 \end{cases}
 \tag{F.HC13}\hypertarget{tag-02GD}{}\label{maps:eq:F.HC13}
\]
and its oriented smooth manifold is
$3\mathbb{HP}^2\#(-\mathbb{HP}^2)$.
It has bounded-cusp holomorphic periods in the sense of \ref{maps:src:F.3.11}.

The fundamental group and the integral homology of every other gluing are computed by \cref{maps:src:F.2.31,maps:src:F.2.35}.

\end{corollary}


\begin{proof}
\begin{enumerate}[label={$\langle1\rangle$\arabic*.}]
\item The exclusions \ref{maps:src:F.2.7}, \ref{maps:src:F.2.14}, \ref{maps:src:F.2.25}, and \ref{maps:src:F.2.27}
leave exactly the stated minimal orders and rational
quotient problem.
The rational normal form is injective on conjugacy
classes, including the exceptional line.
The finite interval and extension-coset enumeration
in \ref{maps:src:F.2.22} gives every integral class over each rational
point, with all bounds proved for that input.
No distinct rational points can become integrally
conjugate, since an integral conjugacy induces a
rational quotient conjugacy.
Theorem \ref{maps:src:F.2.28} and the necessary-and-sufficient
criterion \ref{maps:src:F.2.29} prove both directions of the
retention rule and account for every affine twist.
This is also an exhaustive enumeration if desired:
enumerate primitive projective integer triples with
first nonzero coordinate positive by increasing height,
and enumerate reduced rational exceptional parameters
by increasing numerator and denominator height.
Every finite stage has finitely many inputs, each
terminates by \ref{maps:src:F.2.22}, and every rational input occurs.
No finite cutoff is imposed or needed.

\item Hurewicz makes these simply connected manifolds three-connected.  The only singular contribution to
Euler number is the Kuhn cusp, whose value is six
by \SeriesRef{maps:src:E.1.3}; torus bundles and the free finite caps have
Euler number zero.
Integral Poincar\'e duality and the universal coefficient
theorem then give \eqref{maps:eq:F.HC13}, including torsion freeness.

\item For the smooth classification, \ref{maps:src:F.2.15} supplies a
compatible complex structure at each finite cap, and
the cusp already has its complex model.
Choose their common fibre orientation to be the one
induced by the cusp.
Commuting complex structures at the finite caps can
be chosen in that component: their two-dimensional
fixed real summands allow the orientation to be
reversed when necessary.
The fibre of the bundle of orientation-compatible
complex structures retracts to $\mathbb{CP}^3$.
Its connectedness and simple connectivity extend
the prescribed boundary structures across the
two-dimensional pair of pants, exactly as in \ref{maps:src:F.2.5}.
A compatible horizontal distribution gives an
almost complex structure on the total manifold.
Theorem \S\ref{sec:V-known-inputs} therefore gives signature two,
intersection lattice $I_{3,1}$, and $c_2^2=2$.
Since $H^2=0$, $p_1/2=-c_2$ is primitive,
characteristic, and has square two.
The complete integral reflection, Wall-form and Kramer--Stolz argument in \ref{maps:src:F.4.1}
depends only on these almost complex characteristic
numbers and three-connectivity.
It applies without alteration and gives the claimed
exact oriented diffeomorphism.
Finally \ref{maps:src:F.3.11} applies to the resulting determinant-one
monodromy and primitive cusp.
It gives the punctured-base periods with bounded cusp correction.

\end{enumerate}
\end{proof}




\begin{theorem}[complete fundamental-group data for every input]\hypertarget{tag-02NM}{}\label{maps:src:F.2.31}
\StatementRecord{maps:src:F.2.31}{proved}{02NM}{\SeriesRef{maps:src:E.1.3}; \cite[Theorems 2.3--2.4, p.~425; Euler-characteristic comparison, pp.~427--428]{rootScott83}}


Let $k=2$ or $3$, and let an integral rank-six tuple satisfy the
finite-order and primitive maximal-cusp conditions of \cref{def:admissible}.  No
irreducibility, determinant-one, or three-connectivity hypothesis is
imposed.  Specify each finite affine action as
\[
 \tau_i(x)=M_ix+v_i/m_i,\qquad v_i\in\Lambda=\mathbb Z^6,
 \quad \tau_i^{m_i}=1\text{ on }\Lambda_{\mathbb R}/\Lambda,
 \tag{F.G1}\hypertarget{tag-02FO}{}\label{maps:eq:F.G1}
\]
and require its cyclic action to be free.  The numerator need not be
fixed by $M_i$.  Define its integral power vector and the cusp
lattice by
\[
 w_i=\frac1{m_i}\sum_{j=0}^{m_i-1}M_i^jv_i\in\Lambda^{M_i},
 \qquad U=\operatorname{im}\bigl((M_1\cdots M_k)^{-1}-I\bigr).
 \tag{F.G2}\hypertarget{tag-02FR}{}\label{maps:eq:F.G2}
\]
Use left conjugation transport for the peripheral generators $s_i$,
and record the integral translation clutching by $\kappa\in\Lambda$
in $s_1\cdots s_kx^\kappa=1$.  Then the full fundamental group is
\[
 \left\langle x^\lambda,s_1,\ldots,s_k\ \left|\
 \begin{array}{l}
 x^\lambda x^\mu=x^{\lambda+\mu},\quad x^u=1\ (u\in U),\\
 s_ix^\lambda s_i^{-1}=x^{M_i\lambda},\quad
 s_i^{m_i}=x^{w_i},\quad s_1\cdots s_kx^\kappa=1
 \end{array}\right.\right\rangle .                         \tag{F.G3}\hypertarget{tag-02FS}{}\label{maps:eq:F.G3}
\]
Only a basis of $\Lambda$ and of $U$ is needed in this finite
presentation.  The program below also determines exactly which
quotient of the apparent fibre subgroup injects.

Put
\[
 L=\sum_{g\in\langle M_1,\ldots,M_k\rangle}gU,
 \qquad A=\Lambda/L.                                      \tag{F.G4}\hypertarget{tag-02FT}{}\label{maps:eq:F.G4}
\]
This is an integral sum, without saturation.  For two points, put
$d=\gcd(m_1,m_2)$, $M=M_1$, and
$S_j(M)=I+M+\cdots+M^{j-1}$.  Choose integers $a,b$ with
$am_1+bm_2=d$, and set, in $A$,
\[
 \begin{split}
 w'_2&=-w_2-S_{m_2}(M)\kappa,\\
 e&=(m_2/d)w_1-(m_1/d)w'_2,\qquad
 z=aw_1+bw'_2,\qquad \overline A=A/\langle e\rangle .
 \end{split}                                             \tag{F.G5}\hypertarget{tag-02FU}{}\label{maps:eq:F.G5}
\]
There is an exact sequence and a complete marked extension description
\[
 1\longrightarrow\overline A\longrightarrow\pi_1
   \longrightarrow C_d\longrightarrow1,
 \qquad txt^{-1}=M(x),\quad t^d=z.                         \tag{F.G6}\hypertarget{tag-02FV}{}\label{maps:eq:F.G6}
\]
Thus its order is $d|\overline A|$ if $\overline A$ is finite;
otherwise it is infinite.  In particular $d=1$ gives precisely
the abelian group $\overline A$.

For three genuine points, so $m_i\ge2$, let
\[
 \Gamma=\langle q_1,q_2,q_3\mid q_i^{m_i}=1,
                         \ q_1q_2q_3=1\rangle,
 \qquad \epsilon=\sum_i1/m_i-1.                           \tag{F.G7}\hypertarget{tag-02FW}{}\label{maps:eq:F.G7}
\]
The action $q_i\mapsto M_i$ is well defined on $A$.
If $\epsilon\le0$, the exact sequence is
\[
                    1\longrightarrow A\longrightarrow\pi_1
                         \longrightarrow\Gamma\longrightarrow1. 
 \tag{F.G8}\hypertarget{tag-02FX}{}\label{maps:eq:F.G8}
\]
If $\epsilon>0$, $\Gamma$ is the finite spherical triangle group.
Using its action on $A$, form the exact transfer vector
\[
 e=\sum_{g\in\Gamma}g\kappa+
       \sum_{i=1}^3\ \sum_{g\in\Gamma/\langle q_i\rangle}g w_i.
 \tag{F.G9}\hypertarget{tag-02FY}{}\label{maps:eq:F.G9}
\]
The second sums use left coset representatives and are independent
of their choice.  In this case the exact sequence has kernel
$A/\langle e\rangle$, with the same lifted relations \eqref{maps:eq:F.G3}.
Consequently \eqref{maps:eq:F.G3}, \eqref{maps:eq:F.G6}, \eqref{maps:eq:F.G8}, and \eqref{maps:eq:F.G9} give every such fundamental group by finite operations, including the precise
injected abelian kernel and all extension data.  These statements
do not claim a conjugacy classification of the original tuples.

For every input, the full first integral homology is the Smith
cokernel of the integer matrix
\[
 \begin{pmatrix}
 M_1-I&\cdots&M_k-I&-w_1&\cdots&-w_k&\kappa\\
 0&\cdots&0&&\operatorname{diag}(m_1,\ldots,m_k)&&\boldsymbol1
 \end{pmatrix}.                                           \tag{F.G10}\hypertarget{tag-02FP}{}\label{maps:eq:F.G10}
\]
Here the bottom left zero block has $k$ rows.  The code records
the full matrix, not only its invariant factors.

\end{theorem}


\begin{proof}
\begin{enumerate}[label={$\langle1\rangle$\arabic*.}]
\item \emph{Actual finite and cusp caps.}
The affine lift on $\mathbb R^6$ has $m_i$th power translation by
$w_i$, as the displayed sum shows.  It commutes with that power,
so $M_iw_i=w_i$.  The group generated by lattice translations
and this lift is the fundamental group of the reduced flat
six-manifold.  Its presentation consists of the lattice, its
$s_i$ conjugation relations, and $s_i^{m_i}=w_i$.
The finite cap deformation retracts onto that reduced manifold
by the equivariant radial contraction of its root disc.
Thus these are the cap relations, even for a numerator
which is not an invariant lattice vector.

\item For the primitive Kuhn cap, the compact toric strata in its
normalization and all their intersection components are simply
connected.  Applying van Kampen successively to their tubular
neighborhoods identifies the fundamental group of the central
fibre with that of its dual three-torus.  The nearby fibre map
kills the angular lattice $U$ and sends the long-period lattice
$\Lambda/U$ isomorphically onto this dual-torus group.  This is
also the explicit toric collapse used in \SeriesRef{maps:src:E.1.3}.  The primitive
index-one condition gives $U=\ker(M_0-I)$ and makes $\Lambda/U$
free of rank three.  A section through a smooth central stratum
bounds the base meridian by its parameter disc.  The smooth cap
retracts onto the central fibre, so it adds exactly these
relations, with no extra meridian relation left unspecified.

\item Over the punctured sphere the bundle has a section up to homotopy,
because the base retracts to a graph.  Its group is the appropriate
semidirect product.  Choose the cap sections and retain their
relative translation clutching as $\kappa$.  Van Kampen with the
boundary maps just described gives exactly \eqref{maps:eq:F.G3}.
Changing the transport convention causes no group ambiguity:
reverse the ordered list and replace each affine map by its
inverse, and replace $\kappa$ by $-\kappa$.  The new linear
matrices are $M_k^{-1},\ldots,M_1^{-1}$, the power vectors are
$-w_k,\ldots,-w_1$, and the cusp image lattice is unchanged.
The substitutions $s_i\mapsto s_{k+1-i}^{-1}$ give the explicit
isomorphism between these two presentations.  This supplies the
dictionary when geometric transport is recorded with opposite
path-composition order.

\item \emph{Integral closure and two-point elimination.}
In a normal abelian fibre subgroup, the normal closure of $U$
is precisely the integer sum \eqref{maps:eq:F.G4}.  Start with its column
Hermite basis and repeatedly append its images under every
$M_i$ and $M_i^{-1}$, then take Hermite normal form again.
If the rational rank increases, it can do so at most three
times.  At the first nonincrease the rational span is invariant
under every generator and its inverse, hence never increases
again.  If the current lattice has saturation index $D$ in
that span, every subsequent strict enlargement divides this
index by an integer at least two.  There can therefore be at
most $\lfloor\log_2D\rfloor$ such enlargements.  This proves
termination and gives the computable bound used by the code;
it never replaces a finite-index lattice by its saturation.

\item On $A$ one has $M_2=M^{-1}$ in the two-point case.  The product
relation eliminates $s_2=t^{-1}x^{-\kappa}$, where $t=s_1$.
Multiplying this word $m_2$ times gives
\[
        (t^{-1}x^{-\kappa})^{m_2}
            =t^{-m_2}x^{-S_{m_2}(M)\kappa}.
\]
Thus the two power relations become $t^{m_1}=w_1$ and
$t^{m_2}=w'_2$.  Both right sides are fixed by $M$ on $A$.
Also $M^{m_1}=M^{m_2}=1$, hence $M^d=1$ there.
Bezout elimination gives exactly $e=0$ and $t^d=z$ in
\eqref{maps:eq:F.G5}.  Replacing the Bezout pair changes $z$ by a multiple
of $e$ and therefore changes no extension datum.

\item To prove the asserted injection, construct
$\overline A\rtimes_M\mathbb Z$ and quotient by its central
element $t^dz^{-1}$.  This cyclic subgroup has nonzero
projection $d\mathbb Z$ to the infinite cyclic factor, so it
intersects $\overline A$ trivially.  The quotient has precisely
the presentation obtained by elimination.  This proves \eqref{maps:eq:F.G6},
including the absence of any further hidden fibre relation.

\item \emph{The three-point kernel and the finite transfer.}
For completeness this step computes the fibre injection; it
does not infer it from the appearance of a normal subgroup in
a presentation.  Let $\mathcal B$ be the orbifold classifying
space obtained from the three-holed sphere by attaching, at
its three boundary circles, the mapping cylinders of
$S^1\to BC_{m_i}$.  Over the punctured sphere take the bundle
with fibre $BA=K(A,1)$ and the action.  For each
boundary, the cyclic extension
\[
 E_i=\langle A,s_i\mid s_i a s_i^{-1}=M_i(a),\ s_i^{m_i}=w_i\rangle
\]
contains $A$ injectively: the same semidirect-product construction
as above applies because $M_i^{m_i}=1$ and $M_iw_i=w_i$.
The associated fibration $BA\to BE_i\to BC_{m_i}$ restricts
on $S^1$ to the required mapping-torus fibration.  Replace the
maps by their mapping-path fibrations before homotopy gluing,
so that the fibre is preserved on every attaching cylinder.
Glue these three fibrations, using the relative clutching
$\kappa$ on the punctured sphere.  Van Kampen gives \eqref{maps:eq:F.G3}
for the resulting total space.  Thus the homotopy sequence
of this fibration over $\mathcal B$ calculates precisely
the image of $A$ in the group under discussion.

\item When $\epsilon\le0$, the universal orbifold cover is the plane
or hyperbolic plane.  Its contractibility identifies
$\mathcal B$ with $B\Gamma$: indeed its Borel construction
has the just-described boundary-cylinder model and contractible
universal fibre.  The Euclidean and hyperbolic alternatives
are the oriented two-orbifold classification in
\cite[Theorems 2.3--2.4, p.~425; Euler-characteristic comparison, pp.~427--428]{rootScott83}.
Since $\pi_2(B\Gamma)=0$, the homotopy sequence proves \eqref{maps:eq:F.G8}.

\item When $\epsilon>0$, the universal orbifold cover is $S^2$.
Pulling the constructed $BA$ fibration to this finite cover
trivializes its local system.  Its transgression
$\pi_2(S^2)\to A$ is its integral clutching vector.
There is one lifted ordinary base cell for each $g\in\Gamma$,
contributing $g\kappa$, and one lifted cone cap for each left
coset of $\langle q_i\rangle$, contributing $gw_i$.
To check the sign, the base relation has fibre term $+\kappa$
and each cone relation has term $-w_i$; the oriented sphere
chain subtracts the cone cells from the base cell.  The result
is therefore the plus sign in both sums of \eqref{maps:eq:F.G9}.
Changing a coset representative by $q_i^j$ does not change
its contribution because $w_i$ is fixed.  The result is
$\Gamma$ invariant, as also checked by the code.  The long
homotopy sequence now says that the only additional kernel
is exactly $\langle e\rangle$.  This proves the spherical
assertion, with the orientation fixed by this cell calculation.

\item There is no implicit infinite search for a spherical group.
After sorting its orders the possibilities are
$(2,2,n),(2,3,3),(2,3,4),(2,3,5)$, with groups of orders
$2n,12,24,60$, respectively.  These follow directly from
the positive spherical-triangle angle condition.
The code constructs the dihedral, tetrahedral, octahedral,
or icosahedral permutation group and finds the marked
peripheral generators in that finite group.  Breadth-first
words supply their actions on $A$; every generator edge is
checked modulo the lattice $L$.  It then computes
the two finite sums in \eqref{maps:eq:F.G9} and their integer Smith form.

\item An order-one site is not a genuine cone.  Its matrix is the
identity and its relation is $s_i=x^{w_i}$.  Eliminate it
from \eqref{maps:eq:F.G3}; moving this vector to the right of the remaining
peripheral product replaces the clutching by
\[
       \kappa+(M_{i+1}\cdots M_k)^{-1}w_i.                 \tag{F.G11}\hypertarget{tag-02FQ}{}\label{maps:eq:F.G11}
\]
This reduces the three-point case with such a site to the
two-point computation.  A primitive rank-three cusp cannot
remain with only one finite linear monodromy: a finite-order
unipotent matrix is the identity.  Thus no smaller case has
been omitted within the hypotheses of this theorem.

\item \emph{Freeness and abelianization.}
Freeness of \eqref{maps:eq:F.G1} is itself decidable exactly.  A nontrivial
stabilizer in a cyclic group contains a subgroup of prime
order, so it is enough to check $\tau_i^{m_i/p}$ for every
prime $p\mid m_i$.  For an affine map $Rx+b$ of a torus,
a fixed point exists precisely when every integral covector
in $\ker(R^t-I)$ takes an integral value on $b$.
Smith form proves this criterion: the image of $R-I$ on
the torus is the connected subtorus annihilated by those
primitive covectors.  The code computes the powers in the
rational affine matrix group and performs these finite tests.

\item Abelianizing \eqref{maps:eq:F.G3} yields \eqref{maps:eq:F.G10}.  No cusp column was lost:
\[
 \operatorname{im}(M_0-I)
       \subset\sum_i\operatorname{im}(M_i-I)
\]
follows by telescoping the product on the right; inversion
preserves the image of its difference from the identity.
The Smith form therefore computes every free and torsion
summand of $H_1$ exactly.

\item The program \texttt{fourfolds\_vankampen.py} saves all input
matrices, power vectors, integral closure traces, group
extension data, and homology presentation matrices.
It checks the nine explicitly listed three-point parameters
of its original test set, four two-point order/clutching
choices, a genuinely nonfixed affine numerator, and a
finite-index closure step.  It also checks a nonabelian
rank-six two-point tuple, a spherical $(2,4,2)$ example
with injected kernel $\mathbb Z^2\oplus\mathbb Z/4$, and
the infinite perfect hyperbolic example with orders
$(5,4,19)$.  The nonabelian test deliberately belongs to
the full $\operatorname{GL}_6$ scope and has orientation-
reversing finite matrices; it is not claimed to admit a
complex compactification.  All bounds above depend on the input.

\end{enumerate}
\end{proof}




\begin{theorem}[a complete three-point logarithmic family]\hypertarget{tag-02NN}{}\label{maps:src:F.2.32}
\StatementRecord{maps:src:F.2.32}{proved}{02NN}{\SeriesRef{maps:src:F.3.7}, \SeriesRef{maps:src:F.2.28}, \eqref{maps:eq:F.LT4}, \eqref{maps:eq:F.LT5}, \SeriesRef{maps:src:F.2.5}, \eqref{maps:eq:F.HC8}, \eqref{maps:eq:F.HC9}, \eqref{maps:eq:F.LT8}, \SeriesRef{maps:src:F.2.29}, \eqref{maps:eq:F.LT6}, \eqref{maps:eq:F.LT9}}


Start with the compact complex fourfold of \ref{maps:src:F.3.7} and its
normalized order-five and order-four affine fillings.
Let $\ell$ be the primitive invariant covector on fibre
homology, with the sign chosen in \ref{maps:src:F.2.28}.
At any additional ordinary fibre choose
$v\in\mathbb Z^6$ and an integer $m\ge2$ such that
\[
 r=\ell(v),\qquad \gcd(r,m)=1.
\tag{F.LT1}\hypertarget{tag-02H3}{}\label{maps:eq:F.LT1}
\]
Perform the marked logarithmic transformation with
translation $v/m$ described below.
This gives a smooth compact complex fourfold $X_{m,v}$
with Euler number six and three genuine finite affine
fillings, whose linear orders are five, four and one
and whose affine orders are five, four and $m$.
Its fundamental group is
\[
 \pi_1(X_{m,v})=
 \langle a,b,c,h\mid h\text{ central},\
   a^5=h,\ b^4=h^{-1},\ c^m=h^r,\ abc=1\rangle .
\tag{F.LT2}\hypertarget{tag-02H4}{}\label{maps:eq:F.LT2}
\]
This is the group of the oriented Seifert manifold
with invariants
$(0;(5,-1),(4,1),(m,-r))$, in the convention
$x_i^{\alpha_i}h^{\beta_i}=1$ and product $x_i=1$.
Its base orbifold is hyperbolic, and its central fibre
$h$ has infinite order.

Put $C_{m,r}=\mathbb Z/(20r-m)\mathbb Z$, with
$\mathbb Z/0\mathbb Z=\mathbb Z$, and let
$\varepsilon=1$ when $20r=m$ and zero otherwise.
All the integral homology groups are
\[
\begin{array}{c|c}
j&H_j(X_{m,v};\mathbb Z)\\ \hline
0,8&\mathbb Z\\
1,6&C_{m,r}\\
2&\mathbb Z^\varepsilon\oplus(\mathbb Z/m)^5\\
3&(\mathbb Z/m)^{10}\\
4&\mathbb Z^4\oplus(\mathbb Z/m)^{10}\\
5&(\mathbb Z/m)^5\\
7&\mathbb Z^\varepsilon .
\end{array}
\tag{F.LT3}\hypertarget{tag-02H5}{}\label{maps:eq:F.LT3}
\]
Under \eqref{maps:eq:F.LT1} the exceptional free-rank case is exactly
$(m,r)=(20,1)$.
In particular $(m,r)=(19,1)$ and $(21,1)$ have
$H_1=0$ but nontrivial, infinite fundamental groups.
More generally this occurs whenever $|20r-m|=1$.
The result holds for every vertical component of $v$ in $\ker\ell$.

\end{theorem}


\begin{proof}
\begin{enumerate}[label={$\langle1\rangle$\arabic*.}]
\item We first specify the complex construction and its marking.
Over a small disc at the ordinary value use a holomorphic
period marking
$T_t=\mathbb C^3/\Lambda_t$.
Let $\omega_v(t)$ be the holomorphic period vector
corresponding to $v$.
On the root disc $t=s^m$ quotient the pulled-back family by
\[
 (s,Z)\longmapsto
   (\zeta_m s,\ Z+\omega_v(s^m)/m).
\tag{F.LT4}\hypertarget{tag-02H6}{}\label{maps:eq:F.LT4}
\]
Condition \eqref{maps:eq:F.LT1} makes $v$ have exact order $m$
in $\mathbb Z^6/m\mathbb Z^6$.
Every nontrivial power therefore acts freely, including
over $s=0$.
The quotient is smooth, has a proper map to the disc,
and has central fibre of multiplicity $m$.
On the punctured disc its holomorphic identification with
the original family is
\[
 (s,Z)\longmapsto
 \left(s^m,\ Z-\frac{\log s}{2\pi i}\omega_v(s^m)\right).
\tag{F.LT5}\hypertarget{tag-02H7}{}\label{maps:eq:F.LT5}
\]
Changing the branch of $\log s$ changes the second
coordinate by a lattice period.
Changing $s$ by $\zeta_m s$ changes its logarithm by
$2\pi i/m$ and cancels the translation in \eqref{maps:eq:F.LT4}.
Thus \eqref{maps:eq:F.LT5} is a well-defined holomorphic isomorphism
of the quotient with the original torus family off zero.
Gluing on a fixed annulus supplies the asserted compact
complex manifold: its charts are the original charts
and the free finite quotient charts \eqref{maps:eq:F.LT4}, and it is
the union of compact pieces over closed base regions.
The gluing is an isomorphism on collars, so the total
space is Hausdorff and smooth.
The Euler number is unchanged, since the removed
ordinary torus cap, the new cap, and their boundary
all have Euler number zero.
This marking fixes the new relation $c^m=v$ and
zero extra cusp clutching.

\item In the original datum the invariant normal closure
of the cusp angular lattice is exactly $\ker\ell$,
as proved in \ref{maps:src:F.2.5}.
Adding the identity linear generator does not change
that lattice.
The fibre image in van Kampen therefore reduces to
the central generator $h$, while the original two
finite relations remain $a^5=h$, $b^4=h^{-1}$.
The new relation is $c^m=h^{\ell(v)}$ and the base
relation is $abc=1$.
This proves \eqref{maps:eq:F.LT2}, including the absence of any
dependence on the vertical part of $v$.

\item One may verify the Seifert identification directly:
take a product circle bundle on a three-holed sphere
and attach its three solid tori with primitive meridian
slopes $(5,-1),(4,1),(m,-r)$.
They are primitive by \eqref{maps:eq:F.LT1}, and van Kampen gives
exactly \eqref{maps:eq:F.LT2}.
Its orbifold base has
$1/5+1/4+1/m<1$ for every $m\ge2$.
Pulling back the circle fibration to its contractible
hyperbolic orbifold cover gives a circle bundle on
the plane, hence the trivial circle bundle.
The subsequent universal cover is the plane times
$\mathbb R$, and the fibre subgroup is infinite cyclic.
This also proves directly that none of these groups
is trivial.
The rational Euler number in these conventions is
$r/m-1/20$.

\item In abelianization the relations $5a=h$, $4b=-h$
give $a=4x$, $b=-5x$, $h=20x$ because five and four
are coprime.
The product relation gives $c=x$.
The remaining relation is $(m-20r)x=0$.
Thus $H_1=C_{m,r}$ exactly, including its integral
generator and the exceptional infinite-cyclic case.

\item We now compute the higher integral groups, including
all extensions.
Write $\mathcal R^q$ for the higher direct-image
sheaf of the original fourfold and
$\mathcal R_m^q$ for that of $X_{m,v}$, after identifying
the bases and their ordinary local systems.
At the new point the reduced finite quotient is a
torus with lattice
$\mathbb Z^6+\mathbb Z(v/m)$.
Its covering by the original torus is cyclic of
degree $m$.
The degree-one pullback has Smith factors
$1,1,1,1,1,m$, so its degree-$q$ cokernel is
\[
 C^q=(\mathbb Z/m)^{\binom5{q-1}},
 \qquad 1\le q\le6 .
\tag{F.LT6}\hypertarget{tag-02H8}{}\label{maps:eq:F.LT6}
\]
There is no cohomology torsion in this torus cap.
The sheaf comparison, which does not assert a map
between the total manifolds, is therefore
\[
 0\longrightarrow\mathcal R_m^q
  \longrightarrow\mathcal R^q
  \longrightarrow C^q_{\{p\}}\longrightarrow0 .
\tag{F.LT7}\hypertarget{tag-02H9}{}\label{maps:eq:F.LT7}
\]
It is the identity away from the new point and the
covering image at that point.

\item For $2\le q\le5$, all original global sections of
$\mathcal R^q$ are supported at the original finite
point or the cusp, by \eqref{maps:eq:F.HC8}--\eqref{maps:eq:F.HC9}.
Their map to $C^q_{\{p\}}$ is zero.
Consequently their $H^0$ groups are unchanged,
their $H^2$ groups are unchanged, and
\[
 0\longrightarrow C^q
  \longrightarrow H^1(\mathcal R_m^q)
  \longrightarrow H^1(\mathcal R^q)
  \longrightarrow0 .
\tag{F.LT8}\hypertarget{tag-02HA}{}\label{maps:eq:F.LT8}
\]
The original degree-one base cohomology of these
sheaves is zero for $q=2,4,5$ and is $\mathbb Z^3$
for $q=3$.
The assertion for $q=5$ also follows immediately
from the original $H^6=0$: its bidegree $(1,5)$
has no incoming or outgoing differential.
Thus \eqref{maps:eq:F.LT8} gives $C^q$ in degrees two, four and
five, and $\mathbb Z^3\oplus C^3$ in degree three.

\item The original four order-two torsion transgressions
remain isomorphisms.
Indeed their boundary suspension representatives and
their images in global coinvariants are the ones in
\ref{maps:src:F.2.29}; adding the identity monodromy changes none
of those coinvariant quotients.
These point-supported boundary maps do not depend
on the free degree-one transgression, which changes
under the new logarithmic gluing.
The free degree-four cusp column likewise remains
zero, since its boundary torsion representative has
zero image in that same unchanged global quotient.
It follows that the full filtrations in degrees
three and five have just one nonzero term:
\[
 H^3=C^2=(\mathbb Z/m)^5,\qquad
 H^5=C^4=(\mathbb Z/m)^{10}.
\tag{F.LT9}\hypertarget{tag-02HB}{}\label{maps:eq:F.LT9}
\]
In degree four there is an exact sequence
\[
 0\longrightarrow
       \mathbb Z^3\oplus(\mathbb Z/m)^{10}
  \longrightarrow H^4
  \longrightarrow\mathbb Z\longrightarrow0.
\]
It splits as an abelian-group extension because
its quotient is free, giving the stated middle group.
Thus no torsion extension has been inferred from
rational ranks.

\item For the remaining free ranks,
$H^1(X_{m,v})=\operatorname{Hom}(C_{m,r},\mathbb Z)$
has rank $\varepsilon$.
The degree-one global section group has rank one,
and the only rational differential involving it is
$d_2:H^0(\mathcal R_m^1)\to H^2(\mathcal R_m^0)$.
Its kernel has rank $\varepsilon$, hence its
cokernel also has rank $\varepsilon$.
The degree-one parabolic group is rationally zero,
and there are no rational degree-two global sections.
Therefore $b_2=\varepsilon$.
For an integral check of this degree-one row, its
finite sheaf cokernel is
$\mathbb Z/\gcd(20,m)$: the original invariant
section restricts as $20\ell$, and its value at
the new covering cokernel is $20r$ modulo $m$.
This finite group changes no rational rank.

\item Poincar\'e duality gives
$H_3=H^5$, $H_4=H^4$, and $H_5=H^3$.
The universal coefficient theorem now identifies
the torsion of $H_2$ with $H^3$ and its free rank
with $\varepsilon$.
It also gives $H^2=C_{m,r}$: when $20r\ne m$ this
is $\operatorname{Ext}(H_1,\mathbb Z)$; when
$20r=m$ it is the free rank-one group.
Duality gives $H_6=C_{m,r}$ and
$H_7=\mathbb Z^\varepsilon$.
This proves the complete table \eqref{maps:eq:F.LT3}.

\item The program \texttt{fourfolds\_three\_point\_family.py}
checks the original four torsion images and zero
cusp image once, then computes the cyclic
covering Smith forms in all exterior degrees one
through six and the exact van Kampen abelianization
for the thirty specified pairs
$m\in\{2,3,4,5,6,8,10,19,20,21\}$,
$r\in\{-1,1,m+1\}$.
It includes a nonzero vertical component in every
test vector.
The unbounded assertion follows from \eqref{maps:eq:F.LT6}--\eqref{maps:eq:F.LT9},
not from extrapolation of those finite tests.

\end{enumerate}
\end{proof}




\begin{theorem}[arbitrary finite affine caps]\hypertarget{tag-02NO}{}\label{maps:src:F.2.33}
\StatementRecord{maps:src:F.2.33}{proved}{02NO}{\SeriesRef{maps:src:F.2.31}}


Let $g(x)=Mx+v/m$ be a free cyclic affine action of affine order
$m\geq2$ on $T^r$, with $M\in\mathrm{GL}(r,\mathbb Z)$ and
$v\in\mathbb Z^r$.  For every specified degree bound $d$, there is
an explicit terminating integer algorithm producing chains and
boundary maps through degree $d+1$ for the pair
\[
 B=(D^2\times T^r)/\langle(s,x)\mapsto
 (e^{2\pi i/m}s,g(x))\rangle,
 \qquad Y=\partial B.                                      \tag{F.CA1}\hypertarget{tag-02FE}{}\label{maps:eq:F.CA1}
\]
The cap complex has at most $2^r$ generators in each degree and
the boundary complex at most $2^{r+1}$.  These are complexes of
finite rank in each degree; only their indicated finite prefixes
are required to compute $H_j(B)$, $H_j(Y)$ and $H_j(B,Y)$ for
$j\leq d$.  Their maps include the affine translations and all integral higher coherence terms.

\end{theorem}
\begin{proof}
\begin{enumerate}[label={$\langle1\rangle$\arabic*.}]
\item We give the complete recursive algorithm.  All chains below are
finite integer dictionaries, so the displayed operations are
effective over $\mathbb Z$.

\item First consider the normalized bar resolution of $\mathbb Z^r$
over $R=\mathbb Z[t_1^{\pm1},\ldots,t_r^{\pm1}]$ and its Koszul
resolution $K=R\otimes\bigwedge^*\mathbb Z^r$.
The latter differential is contraction with $(t_i-1)_i$.
Write $t^a=t_1^{a_1}\cdots t_r^{a_r}$ and put
$S_n(t)=(t^n-1)/(t-1)$, interpreted as its finite integral
Laurent sum for positive, zero or negative $n$.
The following $\mathbb Z$-linear contraction of the augmented
Koszul complex fixes a completely explicit choice:
\[
 s_K(t^a e_I)=
 \sum_{j<\min I}
 t_{j+1}^{a_{j+1}}\cdots t_r^{a_r}S_{a_j}(t_j)
 e_j\wedge e_I,                                           \tag{F.CA2}\hypertarget{tag-02FF}{}\label{maps:eq:F.CA2}
\]
where $\min\varnothing=r+1$ and indices in $I$ are increasing.
The bar contraction is
$s_B(t^a[g_1|\cdots|g_q])=[a|g_1|\cdots|g_q]$; terms with a
zero bar entry are zero.  Set $p_0=1$ and define the
$R$-linear projection and homotopy on coefficient-one bar
generators recursively by
\[
 p_q=s_Kp_{q-1}d_B,
 \qquad
 h_q=s_B(1-ip_q-h_{q-1}d_B).                               \tag{F.CA3}\hypertarget{tag-02FG}{}\label{maps:eq:F.CA3}
\]
Here $i(e_I)$ is the alternating sum of the bars of the ordered
unit vectors in $I$.  The recursion decreases degree.  The
identities for the two contractions give, by induction,
$pi=1$, $dh+hd=1-ip$, $ph=hi=h^2=0$.
For example $s_K^2=s_B^2=0$, $ps_B=s_Kp$, and the alternating
unit-vector formula gives $pi=1$; substituting these into
\eqref{maps:eq:F.CA3} proves the three side identities in the same induction.
Augmenting $R\to\mathbb Z$ gives a strong chain contraction
of the ordinary torus bar complex to $\bigwedge^*\mathbb Z^r$.

\item Translations must enter before this contraction.  Let
$b_j=\sum_{a=0}^{j-1}M^av/m$ and let $p_j$ be its coordinatewise
fractional part.  The $m$ points $p_j$ are distinct because the
action is free.  Use the full fundamental groupoid of $T^r$
on these points: an arrow from object $a$ to object $b$ is an
integer vector $\lambda$, represented by the affine path from
$p_a$ to $p_b+\lambda$.  Composition adds the integer vectors.
The generator acts strictly by shifting objects and, for its
$s$th power, by
\[
 \lambda\longmapsto M^s\lambda+z_s(b)-z_s(a),
 \qquad z_s(a)=M^sp_a+b_s-p_{a+s}\in\mathbb Z^r.             \tag{F.CA4}\hypertarget{tag-02FH}{}\label{maps:eq:F.CA4}
\]
The subscripts on objects are taken modulo $m$.
The nerve of this groupoid is equivariantly homotopy equivalent
to the torus: the simplex with successive arrow labels
$\lambda_i$ maps affinely to the torus with lifted vertices
$p_{a_j}+\sum_{i\leq j}\lambda_i$.
Forgetting the objects contracts its chains to the ordinary
torus bar complex.  An explicit homotopy is the negative
alternating prism which, at each possible position, inserts
the zero-labelled arrow to object zero and replaces all
subsequent objects by zero.  Combining this with \eqref{maps:eq:F.CA3}
gives a contraction $(i,p,h)$ of the groupoid chains to
$\bigwedge^*\mathbb Z^r$ with the same side identities.

\item Use the normalized cyclic bar construction on this strict
groupoid action.  Its differential is $d_v+\delta$, where
$\delta$ is the ordinary group-bar boundary, multiplied by
$(-1)^q$ on vertical degree $q$; its first face applies \eqref{maps:eq:F.CA4}.
The vertical contraction transfers this to cyclic bar chains
with exterior coefficients by the finite formulas
\[
 I=(1+h\delta)^{-1}i,\qquad
 P=p(1+\delta h)^{-1},\qquad D=p\delta I.                  \tag{F.CA5}\hypertarget{tag-02FI}{}\label{maps:eq:F.CA5}
\]
Inverses mean their alternating geometric series.  On a chain
of cyclic bar degree $p$ each series has at most $p+1$ terms,
since $\delta$ lowers that degree.  Multiplication and the
contraction identities directly give $(d_v+\delta)I=ID$ and
$P(d_v+\delta)=DP$,
and a chain homotopy inverse.  Equivalently these are the
finite homological perturbation identities, here derived by
expanding the finite series.  The first term of $D$ is the
cyclic bar differential for the exterior action of $M$.
Every remaining term lowers cyclic bar degree by at least two.

\item For completeness, the second contraction used in the algorithm
is also explicit.  The cyclic periodic resolution has one
$\mathbb Z[C_m]$ generator in each degree, with differential
$g-1$ in odd degree and $1+g+\cdots+g^{m-1}$ in even degree.
On $g^a$, $0\leq a<m$, its augmented contraction is
$\sum_{j=0}^{a-1}g^j$ in even degree and is $1$ precisely when
$a=m-1$ in odd degree, zero otherwise.
Apply \eqref{maps:eq:F.CA3} with this contraction to compare it with the
cyclic bar resolution.  Tensor its coefficients with the
exterior action of $M$.  Transferring the higher terms of $D$
once more by \eqref{maps:eq:F.CA5} terminates: each higher term followed by
the horizontal homotopy decreases cyclic degree by at least
one.  The resulting basis in total degree $n$ is
$\{(p,I):p+|I|=n,\ I\subseteq\{1,\ldots,r\}\}$.
This proves the claimed $2^r$ bound and termination in every
specified degree.

\item \par\medskip\noindent\textbf{Explicit comparison with the periodic cyclic resolution.}\quad 
\hypertarget{tag-02J1}{}\label{maps:mapsF:repair:cyclic-inclusion}
Write $C_m=\langle g\rangle$, $R_m=\mathbb Z[C_m]$, and let
$e_n$ generate the periodic resolution in degree $n$.
Let $B_*$ be the normalized $R_m$-bar resolution, with the
insertion contraction $s_B$.  Besides the periodic contraction
$s_P$ just specified, fix the inclusion recursively by
\[
 i_0(e_0)=[],\qquad
 i_n(e_n)=s_B\bigl(i_{n-1}(d_Pe_n)\bigr),
 \quad n>0,
\]
and extend it $R_m$-linearly.  This is a finite formula.  More
explicitly, with every summation index between $1$ and $m-1$,
\[
\begin{split}
 i_{2j}(e_{2j})
   &=\sum_{a_1,\ldots,a_j}
       [g^{a_1}|g|\cdots|g^{a_j}|g],\\
 i_{2j+1}(e_{2j+1})
   &=\sum_{a_1,\ldots,a_j}
       [g|g^{a_1}|g|\cdots|g^{a_j}|g].
\end{split}
\]
The empty sum of indices gives $i_0(e_0)=[]$ and
$i_1(e_1)=[g]$.  The contraction identity for $s_B$ proves
$d_Bi=id_P$ by induction: the chain to which $s_B$ is applied
is a cycle, since the preceding comparison commutes with
differentials and $d_P^2=0$.

\item Define $p$ and $h$ on coefficient-one bar generators by the
recursion already used for the torus comparison, substituting
$s_P$ for $s_K$, and extend them $R_m$-linearly.  Their values
can be checked without an unspecified comparison theorem.
For $1\leq a_i<m$, put
$\epsilon(a,b)=\lfloor(a+b)/m\rfloor$.  The projection is
\[
\begin{split}
 p_{2j}[g^{a_1}|\cdots|g^{a_{2j}}]
   &=\prod_{\nu=1}^{j}\epsilon(a_{2\nu-1},a_{2\nu})\,e_{2j},\\
 p_{2j+1}[g^{a_1}|\cdots|g^{a_{2j+1}}]
   &=\left(\sum_{b=0}^{a_1-1}g^b\right)
       \prod_{\nu=1}^{j}\epsilon(a_{2\nu},a_{2\nu+1})\,e_{2j+1}.
\end{split}
\]
These expressions follow by induction from
$p_n=s_Pp_{n-1}d_B$: the terms other than the first bar face
are annihilated after the preceding contraction identity is
used.  In $pi(e_{2j})$ only the summand with every
$a_\nu=m-1$ survives.  The same observation, with initial
factor $1$, proves $pi(e_{2j+1})=e_{2j+1}$.
Thus $pi=1$ in every degree.

\item The recursion
$h=s_B(1-ip-hd_B)$ gives $d_Bh+hd_B=1-ip$ by induction.
Moreover $p s_B=s_Pp$ follows from the displayed projection
formulas, while $s_P^2=s_B^2=0$ and $s_Bi=0$ follow directly
from their definitions.  Substitution in the same recursion
gives $ph=hi=h^2=0$.  We therefore have the strong comparison
required for the second finite perturbation series, including
its inclusion, projection, homotopy and side identities.
After tensoring with the exterior coefficient module, the
same identities hold.  Every higher perturbation followed by
the horizontal homotopy lowers cyclic degree, so this added
description changes neither the termination proof nor the
$2^r$ generator bound.

\item The free action makes the Borel construction equivalent to
$T^r/C_m$, and radial contraction identifies this with $B$.
For $Y$ use exactly the same algorithm on $T^{r+1}$ with
matrix $\operatorname{diag}(1,M)$ and numerator $(1,v)$.
The projection forgetting its first circle is equivariant and
represents $Y\to B$ under the radial contraction.
More generally an equivariant integral matrix $L$ acts on
groupoid arrows as
$\lambda\mapsto L\lambda+k_b-k_a$, where
$k_a=Lp_a-p'_a\in\mathbb Z^{r'}$.
Transferring this strict bar map by $P'LI$ gives the required
integer chain map.  Its mapping cone computes the pair.
Every series is finite with the proved degree bound, and
integer kernels and Smith normal forms finish the calculation.
\end{enumerate}
\end{proof}



The programs \href{https://github.com/KokunoYumeto/s6-complex-structure-workbench/blob/83d76022bcd442278b760db9ff438afceed6c6ea/pro_returns/S6_10_PAPERS_checkpoint_15_20261119T1045Z/S6_10_PAPERS/code/V/replay_chains/code/fourfolds_bar_koszul.py}{\texttt{fourfolds\_bar\_koszul.py}} and
\href{https://github.com/KokunoYumeto/s6-complex-structure-workbench/blob/83d76022bcd442278b760db9ff438afceed6c6ea/pro_returns/S6_10_PAPERS_checkpoint_15_20261119T1045Z/S6_10_PAPERS/code/V/replay_chains/code/fourfolds_affine_cap_chains.py}{\texttt{fourfolds\_affine\_cap\_chains.py}} implement these
formulas, including the prime-subgroup freeness check of \ref{maps:src:F.2.31}.
The first checks all 281 abelian and 498 cyclic bar simplices
in its explicitly stated alphabets and degree ranges.
The second checks translated circles of orders two and three,
two affine presentations of the Klein bottle (one with a
nonfixed numerator), and the three-dimensional half-turn
manifold.  The separate rank-six program checks every degree
zero through seven for the order-five and order-four caps,
against the independent integral Wang sequence of their
five-torus mapping-torus descriptions.
The boundary-map program checks both translated circles and
the Klein bottle, including the full chain-map identities and
the integral Thom-shift homology of each cone.  These are
finite verification sets, not an asserted exhaustive test over
all input matrices; the terminating formulas prove generality.







\begin{theorem}[the marked Kuhn cusp pair]\hypertarget{tag-02NP}{}\label{maps:src:F.2.34}
\StatementRecord{maps:src:F.2.34}{proved}{02NP}{\SeriesRef{maps:src:F.2.33}; \cite[Theorem 4.1.1 and Lemmas 4.1.2--4.1.3, pp.~480--482]{ffAchingerOgus20}; \cite[Theorem 8.7, pp.~57--59, and Corollary 8.8, p.~59, arXiv v1]{Yang2026}}


The fixed primitive rank-three Kuhn cusp admits an explicit
finite integer chain complex for its cap $W$, its boundary
$Y$, and the inclusion $Y\to W$.  It is obtained from a
category with 26 objects, 170 nonidentity morphisms, and
respectively $26,170,288,144$ strict flags of lengths
$0,1,2,3$.  The construction records the full integral
collapse map.  In particular it gives
\[
\begin{array}{c|rrrrrrrrr}
j&0&1&2&3&4&5&6&7&8\\\hline
\operatorname{rank}H_j(W)&1&3&9&9&10&3&1&0&0\\
\operatorname{rank}H_j(Y)&1&4&12&18&18&12&4&1&0\\
\operatorname{rank}H_j(W,Y)&0&0&1&3&10&9&9&3&1
\end{array}                                                \tag{F.CK1}\hypertarget{tag-02FJ}{}\label{maps:eq:F.CK1}
\]
The only torsion in these three rows is
$\operatorname{Tor}H_3(Y)=\mathbb Z/2$.

\end{theorem}
\begin{proof}
\begin{enumerate}[label={$\langle1\rangle$\arabic*.}]
\item Let $\mathcal K$ be the periodic Kuhn triangulation of
$\mathbb R^3$.  Normalize every simplex by translating its
lexicographically first vertex to zero.  The objects are its
simplex orbits; a morphism $\tau\to\sigma$ is every strict
translated face inclusion $\tau+w\subset\sigma$.
Keep distinct translated inclusions as distinct morphisms.
Composition adds the translation vectors.  Enumerating all
nonempty faces of the six chains in the unit cube gives
$(1,7,12,6)$ simplex orbits in dimensions zero through three,
and the displayed morphism and flag counts.  The category
nerve is the barycentric cell model of the dual decomposition
of $\mathbb R^3/\mathbb Z^3$; the direction of the face order
is reversed when interpreted as dual cells.

\item For an object $\tau$ let $L_\tau\subset\mathbb Z^3$ be the
primitive span of its vertex differences.  Its primitivity
follows from unimodularity of every Kuhn simplex.  Choose by
Smith reduction a quotient map
$Q_\tau:\mathbb Z^3\to\mathbb Z^{3-\dim\tau}$ with kernel
$L_\tau$, and an integral section $R_\tau$.
Define two strict diagrams of tori on this category:
\[
\begin{array}{c|c|c}
 &\text{object at }\tau&\text{map for }\tau+w\subset\sigma\\\hline
\mathcal W&T^{3-\dim\tau}&Q_\sigma R_\tau\\
\mathcal Y&T^4&S_w=\begin{pmatrix}1&0\\ w&I_3\end{pmatrix}.
\end{array}                                                \tag{F.CK2}\hypertarget{tag-02FK}{}\label{maps:eq:F.CK2}
\]
The first maps compose because $Q_\sigma$ vanishes on
$L_\tau$, and the second compose because translations add.
The maps $[\,0\ Q_\tau\,]$ form a natural transformation
$\mathcal Y\to\mathcal W$: the only additional term is
$Q_\sigma w=0$, since $w$ is a vertex of $\sigma$.

\item Here these diagrams describe the local topology.
The toric moment quotient on the dual cell to $\tau$ has
angular torus $T^3/(L_\tau\otimes\mathbb R/L_\tau)$;
passing to a face collapses exactly the next angular
subtorus.  Barycentrically subdividing the dual cells gives
the homotopy colimit of $\mathcal W$ for the central fibre.
For its real oriented logarithmic boundary, retain the
argument $\theta_0$ of the base parameter and the three
angular coordinates $\theta$ of the toric chart.  Translating
the chart by $w$ changes them to
$(\theta_0,\theta+w\theta_0)$, which is precisely $S_w$.
Thus the logarithmic boundary is the homotopy colimit of
$\mathcal Y$, and forgetting the logarithmic phases on each
stratum is exactly the natural transformation just given.
This can also be read directly in the unimodular charts
$x_0\cdots x_k=t$: the moment coordinates form the dual
simplex, the four logarithmic phases remain, and the ordinary
stratum retains their quotient by the vertex differences.
The finite face identifications give \eqref{maps:eq:F.CK2}, so there are
no further gluing constants in this fixed marked cap.

\item The Kuhn map with its divisor log structures is proper,
log smooth and vertical after restricting to a small disc.
Achinger--Ogus, Theorem~4.1.1, identifies its topological germ
with the mapping cylinder of this logarithmic collapse
\cite[Theorem 4.1.1 and Lemmas 4.1.2--4.1.3, pp.~480--482]{ffAchingerOgus20}.  Consequently the preceding
homotopy-colimit diagrams and natural transformation model
the pair $(W,Y)$ itself.  Fix the marking by letting $b_i$
be the category loop with translation label $e_i$.  A
homotopy-colimit edge with map $S_{e_i}$ has the relation
$b_i^{-1}\theta_0b_i=\theta_0+a_i$; equivalently
$\theta_0b_i\theta_0^{-1}=b_i+a_i$.  This specifies the
required primitive cusp
$U=\begin{pmatrix}I_3&I_3\\0&I_3\end{pmatrix}$.

\item \par\medskip\noindent\textbf{Hypotheses and marking of the logarithmic comparison.}\quad 
\hypertarget{tag-02J2}{}\label{maps:mapsF:repair:log-cylinder-scope}
Give the small disc its standard divisorial log structure and
the Kuhn total space the divisorial structure of the reduced
central fibre.  In a unimodular chart the characteristic
monoid map is
\[
 \mathbb N\longrightarrow\mathbb N^{k+1},\qquad
 1\longmapsto(1,\ldots,1).
\]
Both monoids are fine and saturated.  The group cokernel is
free, and the underlying chart is the monomial semistable
map $x_0\cdots x_k=t$ with additional smooth coordinates;
this is log smooth.  The quotient monoid is a group, since
the negative of each coordinate generator equals the sum of
the other generators after the diagonal generator is killed.
The morphism is therefore vertical.  The periodic toric
construction is proper over the sufficiently small disc.

\item These are precisely the hypotheses of
\cite[Theorem~4.1.1, pp.~480--482]{ffAchingerOgus20}.
Its proof also notes that exactness is automatic over the
standard log disc.  It identifies the germ with the mapping
cylinder of the logarithmic collapse.  This homeomorphism
depends on a trivialization; Lemma~4.1.3 fixes that
trivialization on the logarithmic central fibre.  We keep
the resulting radial boundary identification and its stated
angular and tropical generators.  This suffices to model
the pair and its inclusion up to marked homotopy.  It does
not replace the cusp marking by another integral
basis.

\item We specify finite integer chains.  Replace every
torus by its ordinary normalized bar complex.  The affine
simplex map from the bar construction to the torus is natural
for all the integer homomorphisms in \eqref{maps:eq:F.CK2}, so this replacement
does not change either homotopy colimit or its collapse map.
Take the simplicial replacement over all strict flags of the
finite category.  Its vertical differential is the torus bar
boundary; its horizontal differential composes adjacent
arrows, with the first face applying the diagram map, and
has the sign $(-1)^q$ in vertical degree $q$.
Apply the first contraction and the finite series \eqref{maps:eq:F.CA5}
object by object.  There are no flags longer than three,
so every series has at most four terms.  The resulting
finite basis consists of a strict flag beginning at $\tau$
and an exterior basis vector of its torus.  Transfer the
natural transformation by $P_{\mathcal W}FI_{\mathcal Y}$.
This proves both the chain-map identity and completeness of
the collapse map, including all integral coherence corrections.

\item The raw cap chain ranks in degrees zero through six are
\[
 (26,199,688,1231,1226,648,144),
\]
and the raw boundary ranks in degrees zero through seven are
\[
 (26,274,1124,2420,3010,2186,864,144).
\]
The program enumerates every object, every translated
incidence and every flag, verifies composition, checks
$d^2=0$ on every generator, and checks the collapse-chain
identity on every generator.  Integral cancellation of a
unit pivot is a chain homotopy equivalence: clear its row by
integer column operations, then remove the resulting
contractible two-generator summand.  The adjacent differential
updates are forced by $d^2=0$.  After respectively 2063,
4988 and 7087 such cancellations for the cap, boundary and
collapse cone, Smith reduction gives \eqref{maps:eq:F.CK1}.  All matrices
and the remaining differentials are saved by
\texttt{fourfolds\_kuhn\_cap\_chains.py}; the reduced cap and
relative differentials vanish, and the boundary has one
nonzero Smith factor, equal to two in degree four to three.
This computes the integral groups without a filtration
extension problem.
\end{enumerate}
\end{proof}


The finite-pair and cusp-pair chain models enter the global
gluing cone of \cref{maps:src:F.2.35}. Its maps to the common
torus-bundle complement retain the peripheral word
and clutching translation.

The periodic toric quotient used for the cusp is the construction in
Yang's theorem \emph{Higher-dimensional periodic quotient}, and its
Kuhn-grid Euler number is the formula in his corollary
\emph{The standard grid family}
\cite[Theorem 8.7, pp.~57--59, and Corollary 8.8, p.~59, arXiv v1]{Yang2026}.
For the present primitive rank-three grid, this gives Euler number
$3!=6$.  The full cap and boundary complexes, their integral inclusion,
and the torsion groups are calculated in the proof above.







\begin{theorem}[complete integral homology of a marked gluing]\hypertarget{tag-02NQ}{}\label{maps:src:F.2.35}
\StatementRecord{maps:src:F.2.35}{proved}{02NQ}{\SeriesRef{maps:src:F.2.31}, \SeriesRef{maps:src:F.2.33}, \SeriesRef{maps:src:F.2.34}}


For every admissible datum there is a terminating exact
integer algorithm computing every integral homology group of
the resulting eight-manifold, including every torsion invariant
factor.  The input consists of the finite homology matrices
$M_1,\ldots,M_k$, free affine numerators $v_i/m_i$ of their
affine orders, the integral clutch vector $\kappa$, and the
chosen primitive cusp marking
\[
 M_0=(M_1\cdots M_k)^{-1},\qquad
 M_0P=P\begin{pmatrix}I_3&I_3\\0&I_3\end{pmatrix},
 \qquad P\in\mathrm{GL}(6,\mathbb Z).                       \tag{F.GC1}\hypertarget{tag-02FZ}{}\label{maps:eq:F.GC1}
\]
Here $k=2$ or $3$ and the genuine orders are at least two;
order-one sites are first eliminated as in \ref{maps:src:F.2.31}.
If no marking is given, the algorithm chooses one by Smith reduction.
The complete chain prefix through degree nine has at most
4003 generators in any degree.  

\end{theorem}
\begin{proof}
\begin{enumerate}[label={$\langle1\rangle$\arabic*.}]
\item We supply both global boundary maps which were left outstanding
after \ref{maps:src:F.2.34}.  Let $F_k$ be the free group with ordered
generators $s_1,\ldots,s_k$, and let $M(w)$ be the product
specified by a reduced word $w$.  The torus-bundle complement
over the punctured sphere retracts to the corresponding bundle
over a graph.  Its fundamental group and multiplication are
\[
 \Gamma=\mathbb Z^6\rtimes F_k,
 \qquad (x,w)(y,u)=(x+M(w)y,wu).                            \tag{F.GC2}\hypertarget{tag-02G0}{}\label{maps:eq:F.GC2}
\]
The fibre and base are aspherical and the base has no
two-cells.  Hence a chosen section identifies the complement
up to homotopy with $B\Gamma$.  These section and peripheral
choices are fixed by the input clutch convention below.

\item \emph{A finite complement complex with its full bar projection.}
Use reduced words for $F_k$.  The tree resolution of the free
group has $K_0=\mathbb Z[F_k]$ and
$K_1=\bigoplus_i\mathbb Z[F_k]e_i$, with
$de_i=s_i-1$, and no higher terms.  Its explicit augmented
contraction sends a word to its signed path in the Cayley
tree: for each positive letter $s_i$ add its preceding prefix
times $e_i$, and for a negative letter $s_i^{-1}$ subtract
the prefix including that letter times $e_i$.
This is the Fox derivative, and its boundary is $w-1$.
Together with the bar insertion contraction it gives, by
the recursion \eqref{maps:eq:F.CA3}, a strong contraction of the free-group
bar resolution to this tree resolution.  Every recursion
decreases bar degree and every word is finite.

\item For a bar simplex $[(x_1,w_1)|\cdots|(x_n,w_n)]$ of $\Gamma$,
put $y_j=M(w_1\cdots w_{j-1})x_j$.
The exact Alexander--Whitney chain map to the torus/free-group
double bar complex is
\[
 \sum_{p+q=n}(-1)^{pq}
 [w_1|\cdots|w_p]\otimes[y_{p+1}|\cdots|y_n].              \tag{F.GC3}\hypertarget{tag-02G1}{}\label{maps:eq:F.GC3}
\]
Terms with an identity bar entry are zero.  The horizontal
first face acts on torus vectors by
$\rho(w)=M(w)^{-1}$, and the horizontal boundary has sign
$(-1)^q$.  Notice that $\rho$ is a right action:
$\rho(wu)=\rho(u)\rho(w)$.  This order is essential when
the matrices do not commute.  Formula \eqref{maps:eq:F.GC3} follows from
the diagonal bisimplicial bar construction; alternatively,
its chain-map identity follows by pairing adjacent interior
faces, with the first face using
$M(w_1)^{-1}y_j=M(w_2\cdots w_{j-1})x_j$.
The usual shuffle inverse is finite in each degree, so the
map is a chain equivalence.

\item Apply the torus contraction and \eqref{maps:eq:F.CA5}, then the free-group
tree contraction and the same finite perturbation formula.
Higher terms decrease group bar degree at least two, whereas
the tree complex only has degrees zero and one.  Thus the
resulting complement complex is
\[
 C_n(E)=\bigwedge^n\mathbb Z^6\ \oplus\
 \bigoplus_{i=1}^k\bigwedge^{n-1}\mathbb Z^6,
 \qquad
 d(e_i\otimes\xi)=(-1)^{n-1}
   (\bigwedge^{n-1}M_i^{-1}-1)\xi.                         \tag{F.GC4}\hypertarget{tag-02G2}{}\label{maps:eq:F.GC4}
\]
Its differential alone has this short form; its projection
from group bars retains the finite higher perturbation terms.
Those projection terms are used for every boundary map.

\item \emph{Finite affine boundary maps.}
For the $i$th finite cap use the orbit groupoid of \ref{maps:src:F.2.33}.
Its one-edge homotopy colimit for the affine generator is
the boundary mapping torus.  Vertically contracting it gives
a bounded complex with
$C_n(Y_i)=\bigwedge^n\mathbb Z^6\oplus
\bigwedge^{n-1}\mathbb Z^6$ and differential
$(-1)^{n-1}(\bigwedge^{n-1}M_i-1)$ on the second summand.
Its inclusion into the cyclic bar construction, followed by
the cap projection in \eqref{maps:eq:F.CA5}, is the full map $Y_i\to B_i$.
In particular this one-edge model avoids an unnecessary
periodic resolution of the seven-dimensional boundary.

\item Here is an explicit group-bar map for its other inclusion.
Send a torus-groupoid arrow labelled $\lambda$ to the
translation $x^\lambda\in\Gamma$.  With the integral offsets
$z_1(a)$ in \eqref{maps:eq:F.CA4}, put
$\eta_a=s_i^{-1}x^{z_1(a)}$.
For an arrow $f:a\to b$, the exact identity
\[
 \eta_a\,\psi(gf)=\psi(f)\,\eta_b                        \tag{F.GC5}\hypertarget{tag-02G3}{}\label{maps:eq:F.GC5}
\]
is the natural-transformation identity: it is just
$s_i^{-1}x^{M_i\lambda+z_1(b)}
 =x^\lambda s_i^{-1}x^{z_1(b)}$.
For a $q$-simplex, the edge map is $(-1)^q$ times the
alternating prism which inserts $\eta$ at every position,
uses the original arrows before the insertion and the
transformed arrows afterwards.  Its boundary is precisely
the signed difference of the two end maps, by \eqref{maps:eq:F.GC5}.
Compose this prism with the finite vertical lift of the
one-edge complex, then with \eqref{maps:eq:F.GC3} and the complement
projection.  This proves the full chain-map identity.
The offsets record the integer affine lift: replacing
$v_i$ by $v_i+m_i\lambda$ changes the offsets and the map,
even if the affine map on the torus has not changed.
An optional unimodular local basis simply applies its matrix
to every translation in \eqref{maps:eq:F.GC5}.

\item We make the distinction between a bar cell and the geometric
positive root explicit.  Put
$a_i=z_1(0)=\lfloor v_i/m_i\rfloor$, coordinatewise, and
let $\tau_i$ be the positive one-edge bar cell closed using
the zero-labelled return arrow to object zero.  With the
first-face convention above its geometric loop is
\[
 \tau_i=g_i^{-1}x^{a_i},\qquad
 F_E(\tau_i)=-e_i+M_i^{-1}a_i\in C_1(E),\qquad
 g_i=x^{a_i}\tau_i^{-1}.                                  \tag{F.GC5a}\hypertarget{tag-02G4}{}\label{maps:eq:F.GC5a}
\]
Thus the bar cell has the opposite peripheral orientation,
with the displayed integer correction when the numerator
has a nonzero integer part.  Formula \eqref{maps:eq:F.GC5} sends its
group element to $s_i^{-1}x^{a_i}$; consequently the
geometric positive root in the last expression of \eqref{maps:eq:F.GC5a}
maps to $s_i$ exactly.  This is the stated conversion between
the boundary and complement chain conventions.  In a local basis
$Q$, replace $a_i$ in the group and complement formulas by
$Q\lfloor Q^{-1}v_i/m_i\rfloor$.
Writing $\mathcal N_i=\sum_{j=0}^{m_i-1}M_i^j$, the cyclic
degree-two cell has boundary
$m_i\tau_i+w_i-\mathcal N_i a_i$.
Indeed $\tau_i^{m_i}=x^{\mathcal N_i a_i-w_i}$ follows
directly from $g_i^{m_i}=x^{w_i}$.
This independently checks the affine power relation and
the degree-one prism sign, including negative numerators.

\item There is an optional exact geometric shortcut.  If a verified
local integral basis gives $M_i=\operatorname{diag}(1,D)$ and
$v_i=\epsilon e_0$, $\epsilon\in\{1,-1\}$, eliminate the
first fibre generator from the cap group using
$g^{m_i}=u^\epsilon$.  The cap is the mapping torus with
group $\mathbb Z^5\rtimes_D\mathbb Z$.
The fibre vector $(u,a)$ maps to $x^a g^{\epsilon m_i u}$,
and the root generator maps to $g$.  Substituting this
homomorphism for $\psi$ in the same prism gives the boundary
map to this bounded geometric cap model.  The shortcut is
used only after checking the displayed integral normal form;
the general periodic construction remains available for
every other input.

\item To check the peripheral convention directly, the universal
cover of $Y_i$ has the root generator $g_i$ with linear part
$M_i$ and affine displacement $v_i/m_i$.  Its integral
translations and $g_i$ map to the same fibre translations
and $s_i$ in \eqref{maps:eq:F.GC2}.  The full base-circle translation obeys
$g_i^{m_i}=x^{w_i}t$, where $w_i$ is the power vector of
\ref{maps:src:F.2.31}.  Filling the disc kills $t$, giving exactly
$s_i^{m_i}=x^{w_i}$.  The straight paths between the chosen
orbit points give \eqref{maps:eq:F.GC5}, so the prism is this marked map,
including numerators with nonzero integer parts.

\item \emph{The marked cusp boundary map.}
Write $P=[A\ B]$ in \eqref{maps:eq:F.GC1}, and set
\[
 w_0=(s_1\cdots s_k)^{-1},\qquad
 c=w_0x^{-\kappa}\in\Gamma.
                                                               \tag{F.GC6}\hypertarget{tag-02G5}{}\label{maps:eq:F.GC6}
\]
Send the three angular generators to the columns of $A$,
the three tropical generators to the columns of $B$, and
the base-phase generator to $c$.  Relation \eqref{maps:eq:F.GC1} gives
$c b_i c^{-1}=b_i+a_i$ and all the other nilpotent boundary
relations.  Filling kills the angular generators and $c$,
so \eqref{maps:eq:F.GC6} gives exactly
$s_1\cdots s_kx^\kappa=1$, in agreement with \ref{maps:src:F.2.31}.

\item This homomorphism has a direct chain implementation on the
full Kuhn diagram.  For an incidence with translation label
$w$, assign the group element $x^{Bw}$.  A vertical phase
$(n,a)\in\mathbb Z\oplus\mathbb Z^3$ maps to $c^nx^{Aa}$.
The incidence shear is $S_w(n,a)=(n,a+n w)$, and
\[
 (c^nx^{Aa})x^{Bw}=x^{Bw}(c^nx^{A(a+n w)}).                \tag{F.GC7}\hypertarget{tag-02G6}{}\label{maps:eq:F.GC7}
\]
For a flag of $p$ incidence arrows and a vertical bar
simplex of length $q$, sum over all $(p,q)$ shuffles.
Each incidence step contributes its assigned group element;
each vertical step contributes its phase transformed by
all preceding incidence shears.  Give the shuffle the sign
$(-1)^{\#\{\text{incidence steps before vertical steps}\}}$.
Interior face cancellations use exactly \eqref{maps:eq:F.GC7}, and adjacent
incidence faces use addition of their translation labels.
This is therefore a chain map from the double bar complex
to $C_*(B\Gamma)$.  Precompose with the finite Kuhn lift of
\ref{maps:src:F.2.34} and postcompose with \eqref{maps:eq:F.GC3} and the complement
projection.  The other cusp map is the full logarithmic
collapse already computed in \ref{maps:src:F.2.34}.
Both the cusp boundary and the complement are aspherical,
so the displayed marked homomorphism also verifies the
homotopy class of this geometric inclusion.

\item \emph{The final finite integer calculation.}
Take the mapping cone of the just constructed maps
\[
 \bigoplus_{i=0}^k C_*(Y_i)\longrightarrow
 C_*(E)\oplus\bigoplus_{i=0}^k C_*(B_i),
 \qquad y\longmapsto(f_E(y),-f_B(y)),                     \tag{F.GC8}\hypertarget{tag-02G7}{}\label{maps:eq:F.GC8}
\]
where $B_0=W$ is the Kuhn cap.  Gluing collars makes this
the homotopy pushout model of the closed manifold.
Consequently its cone computes integral homology directly;
there is no remaining filtration extension to choose.

\item For clarity the exact finite size bound can be read off
before running any matrix calculation.  Put
$t_n=\sum_{q=0}^{\min(n,6)}\binom6q$, with binomial
coefficients outside their range interpreted as zero.
Let $w_n$ and $y_n$ be the raw cap and boundary ranks of
\ref{maps:src:F.2.34}, zero outside their listed ranges.  Then the cone rank
in degree $n$ is
\[
 \binom6n+k\binom6{n-1}+k t_n+w_n+
 k\binom7{n-1}+y_{n-1}.                                 \tag{F.GC9}\hypertarget{tag-02G8}{}\label{maps:eq:F.GC9}
\]
For $0\leq n\leq9$ its maximum is 3890 when $k=2$ and
4003 when $k=3$.  Every bar comparison and perturbation
series used to produce these matrices terminates by its
degree filtration.  Word reductions and affine orbit
enumerations are finite, with no imposed size cutoff.
Computing integer kernels followed by Smith normal forms
now gives $H_0,\ldots,H_8$; higher groups vanish by dimension.
Unit cancellations are optional exact integral chain
equivalences which substantially reduce the matrices before
Smith reduction.  This proves the algorithm and the bound.
\end{enumerate}
\end{proof}


The programs \texttt{fourfolds\_global\_chains.py}, \texttt{fourfolds\_total\_homology.py} and \texttt{fourfolds\_marked\_gluing.py} implement this algorithm (\S1.5).

A third full integration uses
$A=I_3$, $B=\operatorname{diag}(0,1,0)$ in \eqref{maps:eq:F.IV2},
affine orders $4,4$, numerators $a_1,a_2$, and zero
clutching.  These two linear actions reverse orientation.
Its 7345 integral unit cancellations give free ranks
\[
                  (1,0,4,11,10,1,6,3,0).
\]
The nonzero torsion groups are
\[
\begin{aligned}
 \operatorname{Tor}H_1=\operatorname{Tor}H_5
       &=(\mathbb Z/2)^3\oplus\mathbb Z/4,&
 \operatorname{Tor}H_2&=(\mathbb Z/2)^2,\\
 \operatorname{Tor}H_3=\operatorname{Tor}H_7
       &=\mathbb Z/2,&
 \operatorname{Tor}H_4&=(\mathbb Z/2)^4\oplus(\mathbb Z/4)^5.
\end{aligned}
\]
The independent van Kampen $H_1$ agrees.  Its Euler number
is six, its mod-two Betti vector is
$(1,4,10,14,20,14,10,4,1)$, and its integral top groups
have the required nonorientable form.  Thus this integration
also tests the orientation-reversing part of the stated scope.


The free-bar and semidirect tests, all four periodic-boundary
comparisons, all six one-edge comparison/Thom tests, all seven
peripheral-sign checks, the complete 10048-generator cusp
comparison, and the cached generic three-connected datum computation have
also been independently rerun successfully.  The
reduced-chain, van Kampen $H_1$, Euler and mod-two duality
checks have additionally been independently replayed on
all three complete integration outputs.  The general
algorithm is proved by the explicit chain equivalences and
finite operations above; these finite execution sets are
reported with their exact scope.







\begin{theorem}[simultaneous integral conjugacy for every input]\hypertarget{tag-02NR}{}\label{maps:src:F.2.36}
\StatementRecord{maps:src:F.2.36}{proved}{02NR}{\cite[Algorithm 8.3, p.~21]{rootBHJ22}; \cite[Algorithm 8.3 and its correctness proof, p.~21]{rootBHJ22}; \cite[Theorem 3.4, p.~6; Propositions 8.1--8.2, p.~21]{rootBHJ22}}


There is an unconditional exact algorithm which decides simultaneous
$\mathrm{GL}_6(\mathbb Z)$ conjugacy of any two admissible linear tuples and returns a conjugating matrix when one exists. The algorithm
does not impose a bound on their entries, on the affine multiplicities,
or on the integral conjugator. More generally it applies to any integral
tuple containing a matrix $1+N$ with $N^2=0$ and $\operatorname{rank}N=3$;
the primitive-image assumption is unnecessary for this statement.

\end{theorem}


\begin{proof}
\begin{enumerate}[label={$\langle1\rangle$\arabic*.}]
\item We use column vectors and seek $P$ satisfying $PX_i=Y_iP$ for every
marked index, including the cusp index. Put
\[
\begin{aligned}
 H&=\{P\in\operatorname{Mat}_6(\mathbb Z):PX_i=Y_iP\ (\forall i)\},\\
 \mathcal O&=\{C\in\operatorname{Mat}_6(\mathbb Z):CX_i=X_iC\ (\forall i)\},
 \qquad E=\mathbb Q\mathcal O.
\end{aligned}                                                    \tag{F.IC1}\hypertarget{tag-02GM}{}\label{maps:eq:F.IC1}
\]
Both lattices are saturated integer kernels of explicitly displayed
linear equations in thirty-six unknowns. Thus Hermite or Smith
reduction computes full bases, rather than sublattices of finite index.

\item We first decide rational isomorphism. For a basis $h_1,\ldots,h_s$ of
$H$, form the polynomial $f(c)=\det(\sum c_jh_j)$. If it is zero,
there is no rational isomorphism. Otherwise successively substitute
one of $0,1,\ldots,6$ for each variable while keeping the remaining
polynomial nonzero. At each stage one of these seven values works,
because the degree in that variable is at most six. The final matrix
$D\in H$ is invertible over $\mathbb Q$. This is a terminating exact
test, including the negative case, and not a search for a lucky
invertible intertwiner.

\item The full lattice $L=D^{-1}H\subset E$ is a right $\mathcal O$-module.
Let $J$ be the Jacobson radical of $E$, let $B=E/J$, and use bars
for images under the quotient map. We next verify the splitting
hypothesis needed for the principal-ideal algorithm. In a rational
$N$-adapted basis,
\[
 N=\begin{pmatrix}0&I_3\\0&0\end{pmatrix},\qquad
 E\subseteq
 \left\{\begin{pmatrix}A&T\\0&A\end{pmatrix}:A,T\in
                                      \operatorname{Mat}_3(\mathbb Q)\right\}.
                                                               \tag{F.IC2}\hypertarget{tag-02GN}{}\label{maps:eq:F.IC2}
\]
Projection to $A$ has square-zero kernel and image an algebra
$S\subseteq\operatorname{Mat}_3(\mathbb Q)$. A nilpotent ideal is
contained in the radical, so $E/J\cong S/J(S)$.

\item Choose a composition series of the faithful $S$-module $\mathbb Q^3$.
The kernel of its action on the associated graded module lowers this
filtration and is therefore nilpotent. Conversely the radical kills
every simple factor. Thus this kernel is exactly $J(S)$ and every
simple factor of $S/J(S)$ acts nontrivially on a composition factor.
For a simple algebra $\operatorname{Mat}_n(D)$, where $D$ has centre
$K$ and degree $d$ over that centre, every nonzero simple module has
$\mathbb Q$-dimension $n d^2[K:\mathbb Q]$. If $D$ were
noncommutative this would be at least four, which is impossible in
$\mathbb Q^3$. Consequently
\[
 B\cong\prod_j\operatorname{Mat}_{n_j}(K_j),\qquad
 n_j[K_j:\mathbb Q]\le3.                                    \tag{F.IC3}\hypertarget{tag-02GO}{}\label{maps:eq:F.IC3}
\]
Every noncommutative factor is therefore
$\operatorname{Mat}_2(\mathbb Q)$ or
$\operatorname{Mat}_3(\mathbb Q)$; the remaining factors are number
fields of degree at most three. This proves hypothesis (H) of
\cite[Algorithm 8.3, p.~21]{rootBHJ22} for every tuple under consideration.

\item For clarity this decomposition is effective without a split-algebra
oracle. Rational linear algebra computes $J$ as the radical of the
regular trace pairing and computes the quotient multiplication table.
Its simple factors are obtained from its centre. A commutative factor
is a number field. In a noncommutative factor of dimension $n^2$, we
enumerate integral combinations $e$ of a rational basis until the
left ideal $Be$ has dimension $n$. Such an $e$ exists because the
factor is $\operatorname{Mat}_n(\mathbb Q)$, so its rank-one matrices
are rational points; clearing denominators expresses one among the
enumerated combinations. Left multiplication on $Be$ then supplies
an explicit isomorphism to $\operatorname{Mat}_n(\mathbb Q)$.
The program verifies multiplication on every pair of basis elements.

\item The images $\overline{\mathcal O}$ and $\overline L$ are a full order
and a full right lattice in $B$. If the right order of
$\overline L$ differs from $\overline{\mathcal O}$, this lattice
cannot be principal over that order. Otherwise Algorithm 8.3 of
\cite[Algorithm 8.3 and its correctness proof, p.~21]{rootBHJ22}, in its right-ideal form, decides principality and
returns a generator $\overline\alpha$ when one exists. The cited
algorithm passes through a maximal order and its central conductor,
tests local freeness and principality over the maximal order, and
solves the remaining reduced-norm congruence in a finite quotient.
Its correctness proof uses its Theorem 4.5(b) and Theorem 7.9; its
Theorem 8.4 specifies the number-field, finite-field and integer
factorization subproblems. These have unconditional exact algorithms.
In the recorded implementation the class and unit groups of every
number-field component are certified with \texttt{GRH=false} before
the principal-ideal decision is used. The negative decision thus
does not presume the generalized Riemann hypothesis.

\item If $\overline L$ is not principal, $L$ cannot be principal. If it is
principal, solve the integer linear equations expressing
$\overline\alpha$ as a combination of the images of the chosen basis
of $L$. This produces a lift $\alpha\in L$ and an integral
intertwiner $P=D\alpha\in H$. We claim that a simultaneous integral
conjugacy exists if and only if $|\det P|=1$.

\item The sufficient direction is immediate. For the necessary direction,
suppose $P_0\in H$ is unimodular. Every $h\in H$ has
$P_0^{-1}h\in\mathcal O$, so
$L=\alpha_0\mathcal O$, where $\alpha_0=D^{-1}P_0$.
Write $\alpha=\alpha_0u$ with $u\in\mathcal O$.
Because both image elements generate $\overline L$,
$\overline u$ is a unit of $\overline{\mathcal O}$.
Lift its inverse to $v\in\mathcal O$. Then $uv=1+j$ with
$j\in J\cap\mathcal O$, and the finite geometric series for
$(1+j)^{-1}$ shows that $u^{-1}=v(1+j)^{-1}\in\mathcal O$.
Thus $u$ is an integral unimodular matrix and
$P=P_0u$ is unimodular. This also proves that any chosen lift of the
image generator gives the same yes-or-no answer. It is the
right-module form of \cite[Theorem 3.4, p.~6; Propositions 8.1--8.2, p.~21]{rootBHJ22},
with the needed determinant condition written out.

\item Every branch either proves nonexistence or returns a matrix whose
integrality, determinant and all intertwining equations are verified
exactly. This completes the decision procedure.

\end{enumerate}
\end{proof}




\begin{theorem}[a complete canonical-representative algorithm]\hypertarget{tag-02NS}{}\label{maps:src:F.2.37}
\StatementRecord{maps:src:F.2.37}{proved}{02NS}{\SeriesRef{maps:src:F.2.2}, \SeriesRef{maps:src:F.2.36}, \SeriesRef{maps:src:F.2.38}, \SeriesRef{maps:src:F.2.16}}


Every admissible linear tuple, for either $k=2$ or $k=3$,
has a computable canonical integral representative and a
returned simultaneous unimodular conjugator. Two inputs have
the same representative if and only if they are simultaneously
$\mathrm{GL}_6(\mathbb Z)$ conjugate. The algorithm uses a proved
finite bound depending on the input. No fixed global entry bound
can represent every class.

This is an effective classification. The
ordered involution sector has the separate complete finite
normal form \ref{maps:src:F.2.38}, and all of its $56$ classes have been
computed. 

\end{theorem}


\begin{proof}
\begin{enumerate}[label={$\langle1\rangle$\arabic*.}]
\item We remove the existential free-lift condition by an explicit
construction. Let $d>1$ be the order of $M$ and choose a primitive
$u\in\ker(M-I)\cap\mathbb Z^6$. Such a vector is obtained from a
nonzero rational kernel vector by clearing denominators and dividing
by the greatest common divisor of its integer coordinates; both
operations retain the equation $(M-I)u=0$. Put
$g(x)=Mx+u/d^2$ and $S_d=\sum_{a=0}^{d-1}M^a$.
Then $g^j(x)=M^jx+ju/d^2$, $S_dM^j=S_d$ and $S_du=du$.
A fixed point of $g^j$ gives
$(M^j-I)x+ju/d^2=\lambda\in\mathbb Z^6$.
Applying $S_d$ gives $ju/d=S_d\lambda\in\mathbb Z^6$.
A Bezout row pairing to one with $u$ therefore gives $d\mid j$.
Writing $j=dr$, the original equation becomes $ru/d=\lambda$,
so the same row gives $d\mid r$. Thus only multiples of $d^2$
can have fixed points. Since $g^{d^2}$ is translation by the
integer vector $u$, its torus action is the identity. This proves
both its affine order $d^2$ and freeness of every nonidentity power.
For $M=I$, translation by $u/2$ for primitive $u$ has affine order
two and is free. Conversely, when $I-M$ is invertible over
$\mathbb R$, every affine map $Mx+a$ fixes the torus class of
$(I-M)^{-1}a$. The free-lift test is consequently exactly
$\ker(M-I)\ne0$.

\item The finite linear orders here are the already proved list of
\cref{maps:src:F.2.2}:
\[
 d\in\{1,2,3,4,5,6,8,10,12\}.
 \tag{F.CN1}\hypertarget{tag-02FL}{}\label{maps:eq:F.CN1}
\]
Taking these finitely many powers, checking the fixed kernel,
the product relation and the cusp Smith factors therefore gives
an exact admissibility test. The affine order is the
separate retained gluing datum; the construction in item1 gives
its exact relation to the chosen finite lift.

\item For an admissible full tuple $T=(M_1,\ldots,M_k,M_0)$ let
$H(T)$ be its maximum absolute entry. There is at most one identity
finite slot in a three-point tuple: two such identities would make
$M_0$ finite order, whereas $M_0=I+N$, $N^2=0$ and
$\operatorname{rank}N=3$ give $(I+N)^d=I+dN\ne I$ for every
positive integer $d$. At a fixed identity slot, delete that matrix
and reinsert $I$ at the same slot after classifying the remaining
ordered tuple. These two maps are inverse, and both commute with
$T\mapsto (QM_1Q^{-1},\ldots,QM_0Q^{-1})$ for the same
$Q\in\mathrm{GL}(6,\mathbb Z)$. They therefore retain the returned simultaneous conjugator. The remaining invertible integral
matrices have height at least one, so inserting $I$ leaves height
unchanged. Every entry in the inserted block is constant; the first
unequal row-major entry of two tuples is therefore the same before
and after insertion. This proves preservation of the lexicographic comparison as well. The affine order and
numerator at that slot are retained as gluing data and are transported
by this same $Q$. In the two-involution sector use the canonical
representative of \ref{maps:src:F.2.38}, with its own specified
convention. In all remaining sectors, order full integer tuples
first by $H$, then lexicographically by their row-major entries in
the displayed marked order, and take the first admissible member
of the simultaneous conjugacy class.

\item To compute the representative, enumerate heights
$1,\ldots,H(T)$ and all integer entries in the stated order.
There are at most
\[
 (2H(T)+1)^{36(k+1)}
 \tag{F.CN2}\hypertarget{tag-02FM}{}\label{maps:eq:F.CN2}
\]
integer lists in this box. Retain the invertible tuples satisfying
the ordered product relation and the stated maximal square-zero
primitive cusp condition, which are finite integer tests, and
apply the decision-and-witness algorithm of
\cref{maps:src:F.2.36} to each candidate and the tuple.
Its general rank-three square-zero input permits precisely this
call. A positive answer transfers all finite-order, orientation
and admissibility conditions through its integral
conjugator. The tuple itself occurs, with the identity
as a possible witness, so a positive answer is reached in
finitely many tests. The first one is the required representative.
Record the returned convention as
$P_T^{-1}T_iP_T=\operatorname{Can}(T)_i$.
The displayed candidate bound makes no assumption about entries
of $P_T$; the terminating witness construction is the stated
input from \cref{maps:src:F.2.36}.

\item Every simultaneous orbit has a least element in this order.
Its height is at most the height of each given member, so
\eqref{maps:eq:F.CN2} includes that global least element.
Equivalent inputs have exactly the same ordered orbit and
therefore the same representative. Conversely, if $T,T'$ have
the same representative, the matrix
$Q=P_TP_{T'}^{-1}$ satisfies $Q^{-1}T_iQ=T'_i$ for every slot.
The separate involution convention has the same property by
\cref{maps:src:F.2.38}, with the product of its returned
witnesses giving the corresponding map. Conjugation preserves
the ordered finite orders, so the involution sector is preserved
by those same maps. This proves the classification assertion
and its finite stopping rule.

\item Finally, \ref{maps:src:F.2.16} already gives an unbounded admissible integral
family. In its homology convention direct symbolic
multiplication gives
\[
             \operatorname{tr}(M_1M_2^3)
                       =-v^4+v^3-v^2+2.                 \tag{F.CN3}\hypertarget{tag-02FN}{}\label{maps:eq:F.CN3}
\]
This is unbounded even on the previously constructed free
affine subfamily $v\in20\mathbb Z$. It is invariant under
simultaneous conjugacy. A fixed finite entry box contains
only finitely many tuples and hence only finitely many such
traces; it cannot contain a representative of every class.
The input-dependent nature of \eqref{maps:eq:F.CN2} is therefore necessary
for this particular exhaustive procedure.

\end{enumerate}
\end{proof}




\begin{theorem}[the complete ordered involution stratum]\hypertarget{tag-02NT}{}\label{maps:src:F.2.38}
\StatementRecord{maps:src:F.2.38}{proved}{02NT}{\eqref{maps:eq:F.IV1}, \eqref{maps:eq:F.IV3}, \eqref{maps:eq:F.IV4}, \eqref{maps:eq:F.IV2}, \SeriesRef{maps:src:F.2.31}}


There are exactly 56 simultaneous $\mathrm{GL}_6(\mathbb Z)$
conjugacy classes of ordered tuples $(M_1,M_2,M_0)$ satisfying
the primitive maximal-cusp conditions of \cref{def:admissible} with both finite linear orders equal to two.  An explicit integral normalizer
returns a canonical representative and a checked unimodular
conjugating matrix for every such input, without any bound on
its entries.  Every representative has free affine lifts of
affine order four at both finite points.  None can give a
simply connected filling, for any choices of free affine
orders or twists.

\end{theorem}


\begin{proof}
\begin{enumerate}[label={$\langle1\rangle$\arabic*.}]
\item First choose an integral cusp basis, so that
\[
 M_0=I_6+N,\qquad N=\begin{pmatrix}0&I_3\\0&0\end{pmatrix}.
                                                               \tag{F.IV1}\hypertarget{tag-02GP}{}\label{maps:eq:F.IV1}
\]
This is effective by Smith reduction.  In fact choose vectors
$\beta_i$ whose images $\alpha_i=N\beta_i$ form a primitive
basis of $\operatorname{im}N=\ker N$.
The index-one hypothesis says that their classes form a basis
of the quotient, so $[\alpha_1,\alpha_2,\alpha_3,
\beta_1,\beta_2,\beta_3]$ is unimodular.

\item Since $M_2=M_1(I-N)$ and $M_1^2=M_2^2=I$, we have
$M_1NM_1=-N$.  Solving this block equation gives precisely
\[
 M_1=\begin{pmatrix}A&B\\0&-A\end{pmatrix},\qquad
 M_2=\begin{pmatrix}A&B-A\\0&-A\end{pmatrix},\qquad
 A^2=I_3,\quad AB=BA.                                     \tag{F.IV2}\hypertarget{tag-02GQ}{}\label{maps:eq:F.IV2}
\]
Every integral conjugator preserving \eqref{maps:eq:F.IV1} is
$\begin{pmatrix}C&D\\0&C\end{pmatrix}$, with
$C\in\mathrm{GL}_3(\mathbb Z)$ and $D$ arbitrary integral.
Thus one first classifies $A$ integrally; after fixing it,
$B$ is considered modulo $AD+DA$ and modulo conjugation by
the integral centralizer of $A$.

\item We include the integral involution classification needed here.
For the lattice $L=\mathbb Z^3$, let $L_+$ and $L_-$ be the
primitive kernels of $A-I$ and $A+I$.  They satisfy
$2L\subseteq L_+\oplus L_-\subseteq L$.
The quotient $L/(L_+\oplus L_-)$ is a vector space over
$\mathbb F_2$.  Its projections into $L_+/2L_+$ and
$L_-/2L_-$ are injective: a glue vector with only one
nonzero projection would give a new integral vector in the
corresponding eigensublattice, contradicting its definition.
It is consequently the graph of an isomorphism between two
subspaces of equal dimension $r$.
Unimodular changes of bases in $L_\pm$ put this graph in
coordinate form.  For each glued pair $e_+,e_-$ the vectors
$(e_++e_-)/2$ and $(e_+-e_-)/2$ form a swap block.
The unglued vectors give trivial and sign blocks.  Hence
there is a unique type
\[
 A\cong I_p\oplus(-I_q)\oplus
 \begin{pmatrix}0&1\\1&0\end{pmatrix}^{\oplus r},
 \qquad p+q+2r=3.                                         \tag{F.IV3}\hypertarget{tag-02GR}{}\label{maps:eq:F.IV3}
\]
The dimensions of the two rational eigenspaces and the
index $[L:L_+\oplus L_-]=2^r$ prove uniqueness.
Here $r$ is zero or one, so this normalization only requires
integer kernels, one Smith reduction and, in the glued case,
one explicit half-sum pair whose integrality is checked.

\item For fixed \eqref{maps:eq:F.IV3}, compute the quotient of the commuting
matrices by the subgroup $\{AD+DA\}$.
On the trivial and sign diagonal blocks the norm is
respectively $2D$ and $-2D$.  The commuting off-diagonal
blocks between the trivial and sign summands are zero.
Every commuting block involving a swap summand is a norm.
For example a trivial-to-swap block is a row $(a,a)$ and
is obtained as $(d,e)+(e,d)$; a sign-to-swap block has the
same calculation with a difference.  For two swap summands,
the commuting block is
$\begin{pmatrix}a&b\\b&a\end{pmatrix}$, and applying the
norm to $\begin{pmatrix}d&e\\f&g\end{pmatrix}$ gives
$\begin{pmatrix}e+f&d+g\\d+g&e+f\end{pmatrix}$.
The reverse mixed blocks have the transposed calculations.
Therefore the full integral quotient is exactly
\[
 \operatorname{End}_{\mathbb Z[A]}(L)/\{AD+DA\}
 \cong \operatorname{Mat}_p(\mathbb F_2)
       \oplus\operatorname{Mat}_q(\mathbb F_2).             \tag{F.IV4}\hypertarget{tag-02GS}{}\label{maps:eq:F.IV4}
\]
No finite-index sublattice has been substituted for a norm
image; the displayed block calculations are surjective
integer formulas.

\item The centralizer action on \eqref{maps:eq:F.IV4} is precisely independent
similarity in the two matrix factors.  Indeed, in a product
of commuting block matrices, every contribution factoring
through a swap block is a norm.  The remaining diagonal
blocks are just the reductions of the trivial and sign
blocks of the centralizer matrix.  If that matrix is an
integral unit, its inverse shows that both reductions are
invertible.  Conversely block-diagonal unimodular changes
on the trivial and sign summands induce every element of
$\mathrm{GL}_p(\mathbb F_2)\times
\mathrm{GL}_q(\mathbb F_2)$; elementary matrices give
unimodular lifts.  There is no additional identification
between the two factors, since an intertwiner of a trivial
and a sign lattice is zero over $\mathbb Z$.

\item It remains a genuinely finite classification.  Enumerate
all $2^{n^2}$ matrices of size $n=0,1,2,3$ over $\mathbb F_2$
and conjugate by every invertible matrix of that size.
Their respective numbers of similarity classes are
$1,2,6,14$.  Taking the smallest bit-coded representative
in each orbit gives a fixed canonical choice.
The six integral involution types and their counts are
\[
\begin{array}{c|rrrrrr}
 (p,q,r)&(0,3,0)&(1,2,0)&(2,1,0)&(3,0,0)&(0,1,1)&(1,0,1)\\\hline
 \text{classes}&14&12&12&14&2&2
\end{array}
\]
and their sum is 56.  A representative is \eqref{maps:eq:F.IV2} with
$B$ equal to the chosen $0$--$1$ representatives in its
trivial and sign blocks, and zero in every other block.

\item This also gives an effective normalizer for an arbitrary
input, not merely a list of representatives.  Normalize its
cusp and then its involution by the preceding Smith and
glue constructions.  Use the finite-field orbit lookup to
normalize its two reductions.  In dimensions at most three
the binary lift of an invertible binary matrix is already
unimodular: its determinant is odd and has absolute value
at most two.  The latter bound follows by adjoining a row
and column of ones and applying the $4\times4$ Hadamard
bound.  Finally solve the full integral linear equation
$AD+DA=B-B_{\rm can}$ by Smith reduction.  The preceding
norm calculation proves solvability.  Conjugating by
$\begin{pmatrix}I&D\\0&I\end{pmatrix}$ subtracts this norm
from the upper-right block.  The product of these explicitly
recorded conjugators is integral unimodular and carries all
three marked matrices to their canonical representative.

\item Both matrices in \eqref{maps:eq:F.IV2} have a three-dimensional fixed
space.  Choose a primitive integral fixed vector $v_i$.
Then $g_i(x)=M_ix+v_i/4$ has affine order four:
$g_i^2$ is the nontrivial translation $v_i/2$.
This order-two subgroup is free, and every nontrivial
subgroup of a cyclic group of order four contains it, so
the full action is free.  For any other finite affine lift
the affine order is still even, since its linear part has
order two.  The two-point orbifold quotient of \ref{maps:src:F.2.31} is
therefore the nontrivial cyclic group of order
$\gcd(m_1,m_2)$.  Simple connectivity is impossible for
every choice of twists and affine orders.

\end{enumerate}
\end{proof}




\begin{theorem}[the full effective classification and all marked gluings]\hypertarget{tag-02NU}{}\label{maps:src:F.2.39}
\StatementRecord{maps:src:F.2.39}{proved}{02NU}{\SeriesRef{maps:src:F.2.37}, \SeriesRef{maps:src:F.2.38}, \SeriesRef{maps:src:F.2.31}, \SeriesRef{maps:src:F.2.35}, \SeriesRef{maps:src:F.2.3}, \SeriesRef{maps:src:F.2.30}, \SeriesRef{maps:src:F.2.1}, \SeriesRef{maps:src:F.2.2}, \SeriesRef{maps:src:F.2.36}, \eqref{maps:eq:F.CN3}, \SeriesRef{maps:src:F.2.33}, \SeriesRef{maps:src:F.2.7}}


The admissible data and the topology of their gluings are classified as follows.
For $k=2$ and $k=3$, every admissible ordered integral tuple
has the canonical representative of \ref{maps:src:F.2.37}, with an explicit
simultaneous unimodular conjugator.  Equality of these
representatives is necessary and sufficient for simultaneous
integral conjugacy.  There is no fixed global height cutoff;
the finite stopping bound for an input of height $H$ is
$(2H+1)^{36(k+1)}$ candidates before the proved rejection
filters.  The separate involution normal form has exactly
56 classes and is computed completely in \ref{maps:src:F.2.38}.

For every choice of free affine orders and integer
numerators, integral clutch vector, and specified primitive
Kuhn cusp marking, \ref{maps:src:F.2.31} and \ref{maps:src:F.2.35} compute the full
fundamental group and every integral homology group of the
gluing.  The algorithms have no cutoff on affine orders,
twist entries or lattice indices.  The complete list of data
with $\pi_1=1$ and $H_2=H_3=0$ is the constructive
parametrization of \ref{maps:src:F.2.30}, with its exact lattice enumeration
and necessary-and-sufficient finite tests.  Each such filling has the homology and oriented smooth type stated
there.

\end{theorem}


\begin{proof}
\begin{enumerate}[label={$\langle1\rangle$\arabic*.}]
\item The input is finite, so the prime-subgroup test in \ref{maps:src:F.2.31}
decides whether every claimed affine cyclic action is free
and has the stated affine order.  The integral cusp test and
its adapted basis are \ref{maps:src:F.2.1}.  The complete linear admissibility
filter and the canonical representative, including the proved
finite stopping bound, are \ref{maps:src:F.2.2} and \ref{maps:src:F.2.36}--38.  The
characteristic-polynomial filters used by that algorithm are
conjugacy invariants, so none can remove a possible
representative.  The unbounded trace family \eqref{maps:eq:F.CN3} proves
that no fixed finite box could replace the input-dependent
bound.

\item We spell out how the linear classification is used with a
specified gluing.  If $R$ is its returned conjugator and
$P$ is the cusp marking, change the data by
\[
 M_i'=RM_iR^{-1},\qquad v_i'=Rv_i,\qquad
 \kappa'=R\kappa,\qquad P'=RP.                           \tag{F.SY1}\hypertarget{tag-02I6}{}\label{maps:eq:F.SY1}
\]
The affine orders are unchanged.  The coordinate map
$x\mapsto Rx$ on the torus intertwines the finite quotient
charts, since it sends $M_ix+v_i/m_i$ to
$M_i'x'+v_i'/m_i$.  At the cusp, its coordinates in the
marked Kuhn model are unchanged:
$(P')^{-1}x'=P^{-1}x$.  Hence it intertwines every toric
chart, translated face identification and collapse map of
that fixed filling.  It also carries the clutch path to
the path with vector $\kappa'$.  These maps agree on the
collars and give a diffeomorphism of the two fully marked
gluings.  Thus homology can be computed in the original
coordinates, while the canonical data and their group
presentation are also returned.  Independently replacing
$P'$ by another Smith marking is not this comparison.

\item An affine order-one site is ordinary, even if its integer
numerator has been recorded.  Delete it by the explicit comparison of \cref{maps:src:F.2.31}: if it is the $i$th site, replace
$\kappa$ by
$\kappa+(M_{i+1}\cdots M_k)^{-1}v_i$ before deletion.
The cusp matrix and its given marking are unchanged.
There cannot be an order-one site in an admissible two-point
tuple: its other finite matrix would be the inverse of a
nontrivial unipotent matrix.  Nor can a three-point tuple
have two such sites.  Thus this reduction leaves precisely
the genuine two- or three-point inputs of \ref{maps:src:F.2.35} and loses
no case from the input class.

\item For each remaining input, the exact unsaturated normal
closure, cyclic or triangle-group extension and exceptional
spherical Euler relation of \ref{maps:src:F.2.31} give the full fundamental
group.  The finite cap, complement and marked Kuhn chain
maps of \ref{maps:src:F.2.33}--35 give a finite integer Mayer--Vietoris cone
through degree nine.  Integer kernels and Smith reduction
therefore give all $H_0,\ldots,H_8$ with their complete
torsion factors.  Its independent
abelianization check compares the resulting $H_1$ with the
van Kampen answer.

\item Finally \ref{maps:src:F.2.7}--30 prove both directions of the classification.  Genuine three-point orbifold quotients
exclude simple connectivity.  The two-point filters reduce
every three-connected datum to the minimal order-$(5,4)$ spectrum; the
rational normal form, finite stable-lattice and extension
algorithms include every integral representation above it.
The universal affine normalization includes every twist,
and the six exact integral tests decide three-connectivity. Conversely every datum that passes them satisfies all three
topological conditions by those tests.  \ref{maps:src:F.2.30} computes
its common integral homology and applies the almost-complex
and smooth-classification arguments to identify its oriented
type.  There is no residual untested rational-height bound
in this parametrization.  The canonical-representative
algorithm can remove duplicates from it if desired, but
no finite global enumeration is asserted or needed for
these complete per-input decisions.

\end{enumerate}
\end{proof}




\section{Source obstructions and the punctured-base period decision}



\begin{theorem}[exclusion of the abelian-surface source]\hypertarget{tag-02O6}{}\label{maps:src:F.3.4}
\StatementRecord{maps:src:F.3.4}{proved}{02O6}{\SeriesRef{maps:src:F.3.1}; \cite[Definition 4.2.4 and Proposition 4.2.5, p.~44; Theorem 4.2.6, p.~45]{ffDeligne}}


The construction of \cite{E26}, starting with a polarized
algebraic family of abelian surfaces with a toric-rank-two cusp,
taking a fibrewise degree-zero line bundle, and quotienting its
punctured bundle by a lifted translation, cannot have $b_1=b_2=0$
with the stipulated rank-three square-zero Kuhn cusp.
The two sections need not be proportional.

\end{theorem}


\begin{proof}
\begin{enumerate}[label={$\langle1\rangle$\arabic*.}]
\item Let $W$ be the rank-four rational homology local system of the
abelian surfaces.  It is semisimple by
\cite[Definition 4.2.4 and Proposition 4.2.5, p.~44; Theorem 4.2.6, p.~45]{ffDeligne}.
Its toric-rank-two cusp has
$N_W^2=0$ and $\ker N_W=\operatorname{im}N_W$ of dimension two.
Semisimplicity splits $W=W^\Gamma\oplus W'$.
The image of $N_W$ lies in $W'$ whereas $W^\Gamma$ lies in its
kernel.  The equality of image and kernel forces $W^\Gamma=0$.
The coinvariants vanish as well.

\item The rank-six rational homology system $T$ after the two Engel
operations has a natural filtration
\[
 0\subset L\subset E\subset T,\qquad
 L\simeq\mathbb Q,\quad E/L\simeq W,\quad T/E\simeq\mathbb Q.
\]
The lower line is the circle in the punctured line bundle and
the upper quotient is the lifted-translation generator.
No polarization on $T$ is assumed.
The triviality of $T/E$ gives $\operatorname{im}N_T\subset E$.
Maximal square-zero cusp monodromy gives
$\ker N_T=\operatorname{im}N_T$; thus every globally fixed vector
lies in $E$.  Since $W^\Gamma=0$, the fixed vectors of $E$ are
exactly $L$.  Hence $\dim T^\Gamma=1$.
The dual filtration, whose middle term is $W^*$, gives
$\dim(T^*)^\Gamma=1$ by the identical argument.

\item Consequently the two ranks $a,b$ in Theorem \ref{maps:src:F.3.1} are both one.
Its Euler-parity calculation applies without requiring self-duality
of $T$, and contradicts $b_1=b_2=0$.
This excludes the algebraic abelian-surface construction and,
more generally, every such filtration with semisimple middle term.
It does not exclude an arbitrary unpolarized torus period family.

\end{enumerate}
\end{proof}




\begin{theorem}[the unpolarized period domain has LSAP]\hypertarget{tag-02OA}{}\label{maps:src:F.3.8}
\StatementRecord{maps:src:F.3.8}{proved}{02OA}{\eqref{maps:eq:F.O3}, \eqref{maps:eq:F.O4}; \cite[Definition 2.1, p.~4]{IVVOkaOneV5}; \cite[Proposition 7.13 and preceding paragraph, pp.~23--24]{IVVOkaOneV5}}


Let $D_3$ be the component of
\[
 \{F\in\operatorname{Gr}(3,\mathbb C^6):F\cap\overline F=0\}
\]
containing the standard complex structure, so that
$D_3\cong GL^+(6,\mathbb R)/GL(3,\mathbb C)$.
At every point, $D_3$
is dominated by a spanning tube of nine complex lines.  In particular,
it has the local spray approximation property (LSAP) and is Oka--1.
These assertions concern the unpolarized domain.

\end{theorem}


\begin{proof}
\begin{enumerate}[label={$\langle1\rangle$\arabic*.}]
\item Use the graph chart $F_Z=\operatorname{im}\binom Z{I_3}$.
The open set
\[
 \Omega=\{Z\in M_3(\mathbb C):\det(\operatorname{Im}Z)>0\}
 \tag{F.O1}\hypertarget{tag-02HC}{}\label{maps:eq:F.O1}
\]
lies in $D_3$ and contains $iI_3$.  For a real column $w$, write
$K_w x=w\times x$, put $q_w=w^tw$, and define
\[
 A_w=q_wI_3-ww^t+K_w.
 \tag{F.O2}\hypertarget{tag-02HD}{}\label{maps:eq:F.O2}
\]
Take the six vectors
\[
 e_1,\ e_2,\ e_3,\ e_1+e_2,\ e_1+e_3,\ e_2+e_3
\]
and the corresponding six matrices $A_w$, followed by the three
matrices $K_{e_1},K_{e_2},K_{e_3}$.  Call these nine matrices
$A_1,\ldots,A_9$.  Their symmetric parts are a basis of the
six-dimensional space of real symmetric matrices: the first three
give all diagonal matrices and the last three add the three distinct
off-diagonal entries.  The final three skew matrices form a basis
of the skew matrices.  Thus the nine directions are linearly
independent over $\mathbb R$ and also over $\mathbb C$.
For reference, the determinant of their nine row-major coordinate
columns is $-16$.

\item Each $A_w$ is real normal, is zero on $\mathbb Rw$, and has
eigenvalues $q_w\pm i\sqrt{q_w}$ on $w^\perp\otimes\mathbb C$.
For every real $s$ its two transverse squared singular values are
\[
 (1+q_ws)^2+q_ws^2
 =q_w(q_w+1)\left(s+\frac1{q_w+1}\right)^2
       +\frac1{q_w+1}\ge\frac13,
 \tag{F.O3}\hypertarget{tag-02HE}{}\label{maps:eq:F.O3}
\]
since $q_w\in\{1,2\}$; the remaining squared singular value is one.
For a pure skew direction $K_{e_j}$ all singular values of
$I_3+sK_{e_j}$ are at least one.  Their determinants are positive
at $s=0$ and never vanish, so they are positive for all real $s$.

\item Let $\Lambda=\bigcup_{j=1}^9(iI_3+\mathbb C A_j)$ and let
$B_{1/4}$ be the open Frobenius ball of radius $1/4$ in
$M_3(\mathbb C)$.  Then
\[
 T=\Lambda+B_{1/4}\subset\Omega.
 \tag{F.O4}\hypertarget{tag-02HF}{}\label{maps:eq:F.O4}
\]
Indeed, the imaginary part of $iI_3+tA_j+E$ is
$I_3+(\operatorname{Im}t)A_j+\operatorname{Im}E$.
Its unperturbed least singular value is at least $1/\sqrt3$,
whereas $\|\operatorname{Im}E\|_{\rm op}<1/4$.
It therefore stays invertible throughout the straight perturbation
from $E=0$, preserving the positive determinant.
The set $T$ is open and connected and is a union of affine
translates of the spanning tree $\Lambda$.  Thus it is a spanning
tube in the precise sense of \cite[Definition 2.1, p.~4]{IVVOkaOneV5},
not merely a neighborhood of one entire line.  Its inclusion in
$D_3$ is a local biholomorphism at $iI_3$.

\item The real group $GL^+(6,\mathbb R)$ acts holomorphically and
transitively on $D_3$.  Transporting this tube gives the same
domination at every point.  The observation preceding
\cite[Proposition 7.13 and preceding paragraph, pp.~23--24]{IVVOkaOneV5} gives LSAP from
this domination, and its part (b), or Theorem 2.2 of that paper,
gives Oka--1.  The exact program
\texttt{fourfolds\_oka\_tubes.py} checks all nine integer
matrices, their linear independence and the symbolic identity
\eqref{maps:eq:F.O3}.

\end{enumerate}
\end{proof}




\begin{theorem}[holomorphic periods on every open base]\hypertarget{tag-02OB}{}\label{maps:src:F.3.9}
\StatementRecord{maps:src:F.3.9}{proved}{02OB}{\SeriesRef{maps:src:F.3.8}; \cite[Proposition 7.13(d), p.~24]{IVVOkaOneV5}}


Let $B$ be an open connected Riemann surface and
$\rho:\pi_1(B)\to GL^+(6,\mathbb R)$ any representation.
Its associated flat $D_3$-bundle admits a holomorphic section.
More generally, every continuous section which is holomorphic
near a compact Runge set can be deformed through sections to a
holomorphic section, with arbitrary uniform approximation there
and any prescribed finite-order interpolation at finitely many
points of that set.  For an integral representation, this gives
a holomorphic family of complex three-tori with the prescribed
lattice local system and a zero section over $B$.

\end{theorem}


\begin{proof}
\begin{enumerate}[label={$\langle1\rangle$\arabic*.}]
\item The flat bundle
$\mathscr D=\widetilde B\times_\rho D_3\to B$ is a holomorphic
fibre bundle: its transition maps are constant real linear
transformations, acting holomorphically on the Grassmannian.
Its fibre has LSAP by Theorem \ref{maps:src:F.3.8}, so the bundle map is an
Oka--1 map by \cite[Proposition 7.13(d), p.~24]{IVVOkaOneV5}.
The lifting property in Definition 7.7 of that paper, applied
to the identity map of $B$, gives precisely the stated section
deformation and approximation, including the finite jets.


\item There is a continuous section to start with.  An open Riemann
surface has the homotopy type of a one-dimensional CW complex;
a bundle with connected fibre admits a section over its vertices
and edges, and the homotopy lifting property transports this
section to $B$.  Connectedness holds for the chosen component
$D_3$.  Taking the Runge set empty therefore proves existence.

\item For an integral representation let $\mathbb V_\mathbb Z$ be its
lattice local system and let $\mathcal F\subset
\mathbb V_\mathbb Z\otimes\mathcal O_B$ be the resulting
holomorphic rank-three subbundle.  Real transversality says that
projection identifies the real vector bundle with the underlying
real bundle of
$\mathcal Q=(\mathbb V_\mathbb Z\otimes\mathcal O_B)/\mathcal F$.
Hence the image of $\mathbb V_\mathbb Z$ is a full lattice in
every fibre of $\mathcal Q$.  Locally it is uniformly discrete
and acts by holomorphic translations.  The quotient
$\mathcal Q/\mathbb V_\mathbb Z\to B$ is consequently a smooth
proper holomorphic family of compact complex three-tori,
with its evident zero section and the required marking.

\end{enumerate}
\end{proof}




\begin{lemma}[Isolated extension points]\hypertarget{tag-02O2}{}\label{maps:src:F.3.10}
\StatementRecord{maps:src:F.3.10}{proved}{02O2}{\eqref{maps:eq:F.O5}; \cite[Theorem 1.2, preprint p.~3]{ffForstnericPrezelj02}; \cite[Theorem 2.13, p.~9]{IVVOkaOneV5}}


Let $R$ be an open Riemann surface, $V\to R$ a holomorphic vector
bundle, and $\mathscr U\subset\operatorname{Gr}_r(V)$ an open
subset whose projection onto $R$ is surjective.  Let $S\subset R$
be finite.  Suppose that over $R\setminus S$ this projection is
a holomorphic fibre bundle with an LSAP fibre.  Let $K\subset R$
be a compact Runge set containing $S$ in its interior.
Every continuous section of $\mathscr U$ which is holomorphic
near $K$ can be deformed through sections to a holomorphic
section, with arbitrary uniform approximation on $K$ and
interpolation of any given finite jets at finitely many points
of its interior.  No fibre-bundle or Serre-fibration assertion
is required at the points of $S$.

\end{lemma}


\begin{proof}
\begin{enumerate}[label={$\langle1\rangle$\arabic*.}]
\item We give the relative version of the LSAP argument
in \cite[Lemma 5.1, Theorem 6.1 and Proposition 7.13(d)]
{IVVOkaOneV5}.    Enlarge $K$ slightly inside
the domain of holomorphicity, retaining the Runge property.
Choose a strongly subharmonic Morse exhaustion of $R$ starting
with this compact set.  All subsequent boundary arcs, attached
discs and critical one-handles can and will be chosen disjoint
from $S$.  On them $\mathscr U$ is an ordinary holomorphic
bundle with LSAP fibre.

\item The global spray on the old compact set can be taken to be
vertical even though the fibres at $S$ differ.  Indeed, let $f$
be the current holomorphic section on a Stein neighborhood $W$
of that set.  Put $A$ equal to the finite interpolation set,
including $S$, and choose an ideal sheaf $\mathcal I$ vanishing
to one more than the prescribed jet orders at its points.
Cartan's theorem A supplies finitely many holomorphic sections
$T_1,\ldots,T_N$ of
$\operatorname{End}(V)|_W\otimes\mathcal I$ which generate
$\operatorname{End}(V)$ on the compact attaching arcs, after
shrinking $W$ if needed.  The formula
\[
 F(z,a)=\exp\left(\sum_{j=1}^N a_jT_j(z)\right)f(z)
 \tag{F.O5}\hypertarget{tag-02HG}{}\label{maps:eq:F.O5}
\]
defines a holomorphic spray of sections of the Grassmannian
bundle.  Its differential in $a$ surjects onto the vertical
tangent space at every attaching point, since
$\operatorname{End}(V_z)\to
\operatorname{Hom}(f(z),V_z/f(z))$ is onto.
For $a$ sufficiently small it has range in $\mathscr U$ on
a neighborhood of the old compact set, by openness and
compactness.  Every parameter value preserves all the chosen
jets.  

\item For a noncritical extension, subdivide the new annular regions
into small discs lying in bundle trivializations over
$R\setminus S$.  Restrict \eqref{maps:eq:F.O5} to a thin attaching disc.
LSAP of the fibre extends this spray across the attached disc,
approximating the old spray as closely as desired on the
intersection.  The spray gluing used in the proof of
\cite[Lemma 5.1, equations (5.6)--(5.8)]
{IVVOkaOneV5} now applies verbatim to vertical sprays:
the inverse-function step is needed only on the attaching
intersection, where \eqref{maps:eq:F.O5} is dominating.  It produces
reparametrizations of the two sprays and a glued section.
The old part remains in $\mathscr U$ because it is evaluated
in \eqref{maps:eq:F.O5}; the new part remains in $\mathscr U$ because the
LSAP spray takes values in the bundle fibre.  The fixed jets
survive the reparametrizations, being independent of the
spray parameters on the old side.

\item For a critical extension, first extend over the attached real
arc using the given continuous section and the homotopy already
constructed on the old compact set.  Apply Mergelyan
approximation with the same finite jets to a neighborhood of
the union of the old set and this arc, and then apply the
noncritical step.  To see explicitly that this approximation
is by sections, trivialize $V$ holomorphically on the open
Riemann surface: its underlying complex bundle is trivial over
a graph spine, and the frame bundle therefore has a continuous
section.  Its fibre $GL(\operatorname{rank}V,\mathbb C)$ has
the dominating spray $(g,A)\mapsto g\exp A$, so
\cite[Theorem 1.2, preprint p.~3]{ffForstnericPrezelj02} gives a holomorphic
frame.  The section is then a map to the
fixed Grassmannian, to which
\cite[Theorem 2.13, p.~9]{IVVOkaOneV5} applies.  The base
coordinate is kept unchanged.  Close approximation on the
compact arc and old set, followed by shrinking its
neighborhood, preserves membership in the open subset
$\mathscr U$.  This is the critical case of Theorem 6.1 of
that paper, with no approximation performed at a singular
point of the projection: here the projection is a submersion
everywhere, as $\mathscr U$ is open in a Grassmannian bundle.

\item Choose successive errors summably small on the preceding
compact sets and smaller than their distances to the complement
of $\mathscr U$.  The sections converge uniformly on every
compact set to a holomorphic section with range in
$\mathscr U$.  The finite jets persist.  The close homotopies
on successive stages concatenate in the usual exhaustion
construction.  All new attachment regions lie outside $S$;
the old section remains holomorphic on a full neighborhood of
each point of $S$ at every stage.  This proves the lemma.

\end{enumerate}
\end{proof}




\begin{theorem}[the bounded-cusp period decision]\hypertarget{tag-02O3}{}\label{maps:src:F.3.11}
\StatementRecord{maps:src:F.3.11}{proved}{02O3}{\SeriesRef{maps:src:F.2.1}, \eqref{maps:eq:F.O7}, \SeriesRef{maps:src:F.3.8}, \SeriesRef{maps:src:F.3.10}, \SeriesRef{maps:src:F.2.7}, \SeriesRef{maps:src:F.2.15}}


Let
$B^*=\mathbb P^1\setminus\{p_0,p_1,\ldots,p_k\}$,
where $k\ge1$, and let
$\rho:\pi_1(B^*)\to GL^+(6,\mathbb R)$ be a representation.
Assume that the monodromy $M_0$ at $p_0$ satisfies
$(M_0-I)^2=0$ and $\operatorname{rank}(M_0-I)=3$.
Then there is a holomorphic period family with this monodromy
and nondegenerate real transversality everywhere on $B^*$,
equivalently a nondegenerate imaginary part in each local
normalized period chart.
In a marking with $M_0=\left(\begin{smallmatrix}I&I\\0&I\end{smallmatrix}\right)$
its cusp period is
\[
 Z(t)=\frac{\log t}{2\pi i}I_3+C(t),
 \qquad C(t)\text{ holomorphic at }t=0.
 \tag{F.O6}\hypertarget{tag-02HH}{}\label{maps:eq:F.O6}
\]
Any prescribed finite jet of $C$ at zero may be imposed.
For an integral primitive cusp the marking can be integral.
Consequently every admissible datum with three-connected gluing admits such a period family.  No good-reduction
condition at the other, still punctured, points is included
in this assertion.

\end{theorem}


\begin{proof}
\begin{enumerate}[label={$\langle1\rangle$\arabic*.}]
\item Write $N=M_0-I$ and choose a real basis with
$N=\left(\begin{smallmatrix}0&I\\0&0\end{smallmatrix}\right)$.
Its integral version is Theorem \ref{maps:src:F.2.1}.  Choose the component
of the period domain selected by the positive logarithmic
model in \eqref{maps:eq:F.O6}.  This choice matters only when an orientation
has been fixed independently of the marking; every element
of $GL^+(6,\mathbb R)$ preserves either component.

\item Retain the cusp as a point of the \emph{open} surface
\[
 B=\mathbb P^1\setminus\{p_1,\ldots,p_k\}.
\]
The flat complex vector bundle on $B^*$ has a holomorphic
extension $V_D$ across $p_0$: in a local coordinate $t$ use
the single-valued frame obtained by the change
$\exp(-sN)$, where $s=(\log t)/(2\pi i)$.
Changing the logarithm by $2\pi i$ is cancelled by $M_0$.
This specifies the extension directly.

\item Over $B^*$ let $\mathscr U$ be the flat bundle of real-transverse
planes in the chosen component.  At $p_0$ adjoin only the
planes transverse to $\ker N_\mathbb C$, an affine space of
matrices $C\in M_3(\mathbb C)$.  This gives an \emph{open}
subset $\mathscr U\subset\operatorname{Gr}_3(V_D)$.
Indeed, in the chart $F_C=\operatorname{im}\binom C{I_3}$
the flat period is
$\exp(sN)F_C=F_{sI_3+C}$.  For $t\ne0$ membership is exactly
\[
 \det\bigl(y(t)I_3+\operatorname{Im}C\bigr)>0,
 \qquad y(t)=-\frac{\log|t|}{2\pi}.
 \tag{F.O7}\hypertarget{tag-02HI}{}\label{maps:eq:F.O7}
\]
For every fixed $C_0$, and uniformly for $C$ in a sufficiently
small bounded neighborhood of $C_0$, the matrix in \eqref{maps:eq:F.O7} is
invertible with positive determinant when $|t|$ is sufficiently
small.  This follows, for example, from
$y(t)>\|\operatorname{Im}C\|_{\rm op}$ and a straight-line
homotopy to $y(t)I_3$.  Thus a full neighborhood of
$(0,C_0)$ lies in the adjoined set, proving openness.
The projection is consequently a holomorphic submersion,
although its fibre at $p_0$ is $\mathbb C^9$ rather than $D_3$.

\item Start with a polynomial $C_0(t)$ realizing the desired finite
jet.  Formula \eqref{maps:eq:F.O7} makes it a holomorphic section on a small
closed cusp disc $K$, which is Runge in $B$.
This section extends continuously over $B$.  To justify the
relative extension, $D_3$ is connected and simply connected:
it retracts onto $SO(6)/U(3)$, and
$\pi_1(U(3))\to\pi_1(SO(6))$ is onto, since a unitary rotation
in one complex coordinate is a generator of the real
two-plane rotation loop.  A cellular decomposition of the
surface outside the cusp disc now extends its boundary
section over edges by connectedness and over two-cells by
simple connectivity.  The flat bundle is trivial over each
cell.  This argument uses no fibration property at $p_0$.

\item Over $B\setminus\{p_0\}$ the fibre is the LSAP manifold
$D_3$ from Theorem \ref{maps:src:F.3.8}.  Lemma \ref{maps:src:F.3.10}, with $S=\{p_0\}$,
produces a holomorphic section of $\mathscr U$ over all of
$B$, with the chosen jet at $p_0$.  Since its value at $p_0$
lies in the affine chart just adjoined, it has a holomorphic
matrix $C(t)$ there.  Undoing the canonical frame gives
exactly \eqref{maps:eq:F.O6}, with a bounded holomorphic correction.
Everywhere else membership in $\mathscr U$ is precisely
nondegeneracy of the complex-torus periods.
More explicitly, choose three members of a real lattice basis
whose projected periods are complex-linearly independent, and
normalize them to $I_3$.  The remaining three periods form a
matrix $Z$, and real transversality is exactly
$\det\operatorname{Im}Z\ne0$.  This normalization persists
locally.  

\item For admissible data with three-connected gluing, Theorem \ref{maps:src:F.2.7} excludes $k=3$ and
Proposition \ref{maps:src:F.2.15} proves orientation preservation of the
finite matrices for $k=2$.  The cusp determinant is also one.
Thus all three-connected data satisfy the hypotheses.  

\end{enumerate}
\end{proof}





\clearpage
